\documentclass[UTF8,openany,11pt]{ctexbook}
\usepackage[a4paper,margin=2.2cm,headheight=15pt]{geometry}
\usepackage{amsmath,amssymb,mathtools,amsthm,bm}
\usepackage{enumitem}
\usepackage{xcolor}
\usepackage{fancyhdr}
\usepackage[hidelinks,unicode]{hyperref}
\usepackage[nameinlink,noabbrev]{cleveref}
\allowdisplaybreaks[2]
\definecolor{answerblue}{HTML}{1F4E79}
\definecolor{skipgray}{HTML}{666666}
\setlength{\parindent}{2em}
\setlength{\parskip}{0.25em}
\setlist{itemsep=0.25em,topsep=0.35em}
\pagestyle{fancy}
\fancyhf{}
\fancyhead[L]{泛函分析作业完整解答}
\fancyhead[R]{\nouppercase{\rightmark}}
\fancyfoot[C]{\thepage}
\newtheoremstyle{problemstyle}{0.8em}{0.5em}{\normalfont}{}{\bfseries\color{answerblue}}{.}{0.5em}{}
\theoremstyle{problemstyle}
\newtheorem*{problem}{题干}
\newenvironment{solution}{\par\smallskip\noindent\textbf{解答.}\ }{\hfill$\square$\par\medskip}
\newenvironment{skipped}{\par\smallskip\noindent\color{skipgray}\textbf{本题跳过.}\ }{\par\medskip\color{black}}
\newcommand{\R}{\mathbb R}
\newcommand{\C}{\mathbb C}
\newcommand{\N}{\mathbb N}
\newcommand{\F}{\mathbb F}
\newcommand{\cL}{\mathcal L}
\newcommand{\ol}[1]{\overline{#1}}
\DeclareMathOperator{\Ker}{Ker}
\DeclareMathOperator{\Ran}{Ran}
\DeclareMathOperator{\Span}{span}
\DeclareMathOperator{\dist}{dist}
\DeclareMathOperator*{\esssup}{ess\,sup}

\title{泛函分析作业完整解答\\\large 第一次至第二十七次作业}
\author{依据课程作业文件整理}
\date{2026 年 9 月}
\begin{document}
\frontmatter
\maketitle
\chapter*{使用说明}
本答案册覆盖课程作业文件中的全部题号。证明尽量写出所用定义、关键估计和极限过程。题目若只引用未提供讲义中的特定命题、性质或例子，且题面不足以唯一还原原结论，则依照要求保留题号并标注跳过。

文中 $\F$ 表示 $\R$ 或 $\C$；未特别说明时，赋范空间均取题面给定的数域。答案用于学习和核对，提交前应自行重新书写并检查课程采用的符号约定。
\tableofcontents
\mainmatter
\chapter{第一次至第九次作业}
本章中，除非特别说明，标量域为 $\mathbb R$ 或 $\mathbb C$。所有由原题中“讲义第某页”才能确定、而题面没有给出具体陈述的题，均按要求保留题号并跳过。

\section{第一次作业}

\subsection*{1.1 完备子空间与闭子空间}
\begin{problem}
设 $(X,d_X)$ 为度量空间，$Y\subset X$，并用 $d_Y=d_X|_{Y\times Y}$ 表示 $Y$ 上的诱导距离。证明：
\begin{enumerate}[label=(\arabic*)]
\item 若 $(Y,d_Y)$ 完备，则 $Y$ 是 $X$ 的闭子集；
\item 若 $(X,d_X)$ 完备且 $Y$ 是 $X$ 的闭子集，则 $(Y,d_Y)$ 完备。
\end{enumerate}
\end{problem}
设 $d_Y=d_X|_{Y\times Y}$。

（1）设 $(Y,d_Y)$ 完备。取 $x\in\overline Y$，则存在序列 $y_n\in Y$ 使 $d_X(y_n,x)\to0$。收敛序列必为 Cauchy 序列，故 $(y_n)$ 在 $Y$ 中是 Cauchy 序列。由 $Y$ 完备，存在 $y\in Y$ 使 $d_Y(y_n,y)\to0$。另一方面 $d_X(y_n,x)\to0$，而度量空间中极限唯一，所以 $x=y\in Y$。因此 $\overline Y\subset Y$，故 $Y$ 闭。

（2）设 $X$ 完备且 $Y$ 闭。任取 $Y$ 中的 Cauchy 序列 $(y_n)$。因 $d_Y$ 是限制度量，它也是 $X$ 中的 Cauchy 序列，故存在 $x\in X$ 使 $y_n\to x$。由于 $Y$ 闭且 $y_n\in Y$，有 $x\in Y$。于是 $(y_n)$ 在 $Y$ 中收敛，故 $Y$ 完备。

\subsection*{1.2 度量空间的完备化}
\begin{problem}
设 $(X,d)$ 为度量空间。称三元组 $(\widehat X,\widehat d,\iota)$ 为 $(X,d)$ 的完备化，如果 $(\widehat X,\widehat d)$ 完备，且 $\iota:X\to\widehat X$ 是具有稠密像的等距嵌入。证明：
\begin{enumerate}[label=(\arabic*)]
\item 若 $Y$ 完备、$\phi:X\to Y$ 为 $1$-Lipschitz 映射，则存在唯一的 $1$-Lipschitz 映射 $\widehat\phi:\widehat X\to Y$ 使 $\phi=\widehat\phi\circ\iota$；
\item 任意两个完备化之间存在唯一的、与嵌入相容的等距映射；
\item 完备化存在。提示：固定 $x_0\in X$，在 Banach 空间 $C_b(X)$ 中令
$f_x(y)=d(y,x)-d(y,x_0)$，证明 $x\mapsto f_x$ 是等距嵌入，并取其像的闭包；
\item 若 $\widehat X$ 完备，$\iota:X\to\widehat X$ 为 $1$-Lipschitz 映射且满足第 (1) 问的范性质，则 $(\widehat X,\widehat d,\iota)$ 是 $(X,d)$ 的完备化。
\end{enumerate}
\end{problem}
为避免符号混淆，把完备化记为 $(\widehat X,\widehat d,\iota)$。

（1）给定 $1$-Lipschitz 映射 $\phi:X\to Y$，其中 $Y$ 完备。对 $\widehat x\in\widehat X$，由 $\iota(X)$ 稠密，可取 $x_n\in X$ 使 $\iota(x_n)\to\widehat x$。由于
\[
d_Y(\phi(x_n),\phi(x_m))\le d(x_n,x_m)=\widehat d(\iota x_n,\iota x_m),
\]
$(\phi(x_n))$ 是 $Y$ 中的 Cauchy 序列。定义
\[
\widehat\phi(\widehat x):=\lim_{n\to\infty}\phi(x_n).
\]
若另有 $\iota(x_n')\to\widehat x$，则
\[
d_Y(\phi(x_n),\phi(x_n'))\le \widehat d(\iota x_n,\iota x_n')\to0,
\]
故定义与逼近序列无关。若 $\widehat x,\widehat z\in\widehat X$，分别以 $\iota x_n,\iota z_n$ 逼近，则
\[
d_Y(\widehat\phi(\widehat x),\widehat\phi(\widehat z))
=\lim_n d_Y(\phi(x_n),\phi(z_n))
\le\lim_n d(x_n,z_n)=\widehat d(\widehat x,\widehat z).
\]
所以 $\widehat\phi$ 为 $1$-Lipschitz，且显然 $\widehat\phi\circ\iota=\phi$。若 $\psi$ 也是这样的延拓，则 $\psi$ 与 $\widehat\phi$ 在稠密集 $\iota(X)$ 上相等；二者连续，故在 $\widehat X$ 上处处相等。

（2）对 $\iota_2:X\to\widehat X_2$ 应用（1），得到唯一的 $1$-Lipschitz 映射 $\psi:\widehat X_1\to\widehat X_2$，满足 $\psi\iota_1=\iota_2$。交换下标得到 $\chi:\widehat X_2\to\widehat X_1$。于是 $\chi\psi$ 与 $\operatorname{id}_{\widehat X_1}$ 都延拓 $\iota_1$，由唯一性 $\chi\psi=\operatorname{id}$；同理 $\psi\chi=\operatorname{id}$。因为 $\psi,\chi$ 均为 $1$-Lipschitz，
\[
\widehat d_1(u,v)=\widehat d_1(\chi\psi u,\chi\psi v)
\le\widehat d_2(\psi u,\psi v)\le\widehat d_1(u,v),
\]
故每个不等号都是等号，$\psi$ 是唯一的满射等距映射。

（3）固定 $x_0\in X$，定义 $f_x(y)=d(y,x)-d(y,x_0)$。反三角不等式给出 $|f_x(y)|\le d(x,x_0)$，故 $f_x$ 有界；距离函数连续，故 $f_x\in C_b(X)$。又
\[
\|f_x-f_z\|_\infty
=\sup_y|d(y,x)-d(y,z)|\le d(x,z).
\]
取 $y=x$，等号右侧在绝对值中达到：$|d(x,x)-d(x,z)|=d(x,z)$，所以
\[
\|f_x-f_z\|_\infty=d(x,z).
\]
故 $\iota:x\mapsto f_x$ 等距。$C_b(X)$ 在上确界范数下完备，因而闭子空间 $\overline{\iota(X)}$ 完备，而 $\iota(X)$ 在其中稠密。这便给出一个完备化。

（4）题面中的同名空间改记为 $\widehat X$。取 $Y=\overline{\iota(X)}\subset\widehat X$。它是完备空间 $\widehat X$ 的闭子空间，故完备。由范性质，把映射 $\iota:X\to Y$ 延拓为 $1$-Lipschitz 映射 $r:\widehat X\to Y$，满足 $r\iota=\iota$。设 $j:Y\hookrightarrow\widehat X$ 为包含映射。则 $jr$ 与 $\operatorname{id}_{\widehat X}$ 均为 $\iota:X\to\widehat X$ 的 $1$-Lipschitz 延拓，由唯一性 $jr=\operatorname{id}_{\widehat X}$。因此每个 $\widehat x\in\widehat X$ 都属于 $j(Y)=Y$，即 $\iota(X)$ 稠密。故 $(\widehat X,\widehat d,\iota)$ 确为完备化。

\subsection*{1.3 赋范空间的完备化仍是线性空间}
\begin{problem}
证明：赋范向量空间的完备化能够自然地赋予向量空间结构和范数，并在此结构下成为 Banach 空间。
\end{problem}
设 $X$ 的度量完备化由 Cauchy 序列等价类构成，等价关系为 $(x_n)\sim(y_n)$ 当且仅当 $\|x_n-y_n\|\to0$。定义
\[
[x_n]+[y_n]=[x_n+y_n],\qquad \lambda[x_n]=[\lambda x_n],
\qquad \|[x_n]\|:=\lim_n\|x_n\|.
\]
若换取等价代表，三角不等式表明上述结果不变。例如
$|\|x_n\|-\|y_n\||\le\|x_n-y_n\|\to0$，故范数定义良好。线性空间公理由逐项运算直接继承；正定性、齐次性、三角不等式由取极限得到。常值序列给出等距线性嵌入 $X\hookrightarrow\widehat X$，其像稠密。度量完备化本身完备，且该度量正是新范数诱导的度量，所以 $\widehat X$ 是 Banach 空间。

\subsection*{1.4}
\begin{problem}
说明讲义第 8 页注 3 中的赋范线性空间 $X$ 不是完备的，并写出它的完备化空间。
\end{problem}
\textbf{依赖未提供讲义，按用户要求跳过。}题目所称“第 8 页注 3”中的空间 $X$ 未在题面给出，无法唯一确定其完备化。

\subsection*{1.5 复测度的全变差}
\begin{problem}
设 $(X,\mathcal M,\mu)$ 为复测度空间。对 $E\in\mathcal M$ 定义
\[
|\mu|(E)=\sup\left\{\sum_{j\ge1}|\mu(E_j)|:
E=\bigsqcup_{j\ge1}E_j, E_j\in\mathcal M\right\}.
\]
证明：
\begin{enumerate}[label=(\alph*)]
\item $|\mu|$ 是有限正测度；
\item 令 $\mathcal C_m$ 为 $(X,\mathcal M)$ 上所有复测度组成的空间，并定义
$\|\mu\|=|\mu|(X)$，则 $(\mathcal C_m,\|\cdot\|)$ 是 Banach 空间。
\end{enumerate}
\end{problem}
记
\[
|\mu|(E)=\sup\left\{\sum_j|\mu(E_j)|:E=\bigsqcup_jE_j\right\}.
\]
（a）先证可列可加性。若 $E=\bigsqcup_n A_n$，对 $E$ 的任意分割 $E=\bigsqcup_jE_j$，令 $E_{jn}=E_j\cap A_n$，则
\[
\sum_j|\mu(E_j)|=\sum_j\left|\sum_n\mu(E_{jn})\right|
\le\sum_n\sum_j|\mu(E_{jn})|\le\sum_n|\mu|(A_n).
\]
取上确界得 $|\mu|(E)\le\sum_n|\mu|(A_n)$。反之，给定 $\varepsilon>0$，对每个 $n$ 取 $A_n$ 的分割使其变差和大于 $|\mu|(A_n)-\varepsilon2^{-n}$；合并所有分割即得 $E$ 的分割，于是
$|\mu|(E)\ge\sum_n|\mu|(A_n)-\varepsilon$。令 $\varepsilon\downarrow0$，得到等号。因此 $|\mu|$ 是正测度。

还需证明其有限。复测度 $\mu$ 的实部、虚部都是有限的符号测度。由 Hahn 分解，$X$ 可分成有限个集合，使 $\operatorname{Re}\mu$ 与 $\operatorname{Im}\mu$ 在每块上各自符号固定（两次二分共至多四块）。记这些块为 $A_1,\dots,A_4$。对任意分割 $(E_j)$，在固定 $A_k$ 上，
\[
\sum_j|\mu(E_j\cap A_k)|
\le\sum_j\bigl(|\operatorname{Re}\mu(E_j\cap A_k)|+|\operatorname{Im}\mu(E_j\cap A_k)|\bigr)
\]
而由于两部分符号固定，右端等于
$|\operatorname{Re}\mu(A_k)|+|\operatorname{Im}\mu(A_k)|$。对 $k$ 求和得到一个与分割无关的有限上界，故 $|\mu|(X)<\infty$。

（b）范数公理中只有完备性需详证。设 $(\mu_n)$ 在全变差范数下为 Cauchy。对每个 $E\in\mathcal M$，
\[
|\mu_n(E)-\mu_m(E)|\le |\mu_n-\mu_m|(E)\le\|\mu_n-\mu_m\|,
\]
故 $\mu(E):=\lim_n\mu_n(E)$ 存在。若 $E=\bigsqcup_jE_j$，给定 $\varepsilon>0$，取 $n$ 使 $\|\mu-\mu_n\|$ 尚未定义的问题通过下述一致估计处理：先由 Cauchy 性取 $N$，使 $m,n\ge N$ 时 $|\mu_m-\mu_n|(X)<\varepsilon$；令 $m\to\infty$，对任何有限分割有
\[
\sum_{j=1}^r|\mu(E_j)-\mu_n(E_j)|\le\varepsilon.
\]
再取这些有限分割的上确界，得 $|\mu-\mu_n|(X)\le\varepsilon$。这同时说明 $\mu-\mu_n$ 有有限全变差。用 $\mu_n$ 的可列可加性及尾项的全变差控制，即得 $\mu$ 可列可加。故 $\mu$ 是复测度，且 $\|\mu_n-\mu\|\to0$，所以该空间完备。

\section{第二次作业}

\subsection*{2.1 泛函范数不一定达到}
\begin{problem}
在 $C([-1,1])$ 上取上确界范数，并定义
\[
\Lambda(f)=\int_0^1f(t)\,dt-\int_{-1}^0f(t)\,dt.
\]
证明 $\|\Lambda\|=2$，但不存在 $\|f\|_\infty=1$ 的连续函数使 $|\Lambda(f)|=2$。
\end{problem}
对 $\|f\|_\infty\le1$，
\[
|\Lambda f|\le\int_0^1|f|+\int_{-1}^0|f|\le2,
\]
故 $\|\Lambda\|\le2$。令 $f_n$ 在 $[-1,-1/n]$ 上等于 $-1$，在 $[1/n,1]$ 上等于 $1$，并在 $[-1/n,1/n]$ 上线性连接，则 $\|f_n\|_\infty=1$ 且
$\Lambda(f_n)=2-1/n\to2$，故 $\|\Lambda\|=2$。

若某连续 $f$ 满足 $\|f\|_\infty=1$ 且 $|\Lambda f|=2$，更换 $f$ 的符号后可设 $\Lambda f=2$。上述两个三角不等式均须取等，因而 $f=1$ 在 $[0,1]$ 上几乎处处成立、$f=-1$ 在 $[-1,0]$ 上几乎处处成立。连续性把“几乎处处”提升为处处，特别得到 $f(0)=1$ 与 $f(0)=-1$，矛盾。

\subsection*{2.2 $C(0,1]$ 与 $\ell^\infty$}
\begin{problem}
设 $C(0,1]$ 为 $(0,1]$ 上所有有界连续函数组成的空间，并令
$\|f\|=\sup_{0<x\le1}|f(x)|$。证明：
\begin{enumerate}[label=(\arabic*)]
\item $\|\cdot\|$ 是 $C(0,1]$ 上的范数；
\item $\ell^\infty$ 与 $C(0,1]$ 的某个子空间等距同构。
\end{enumerate}
\end{problem}
（1）上确界范数的非负性、齐次性和三角不等式逐点成立后取上确界即得；若 $\|f\|=0$，则每点 $|f(x)|=0$，故 $f=0$。

（2）取两两不交、在 $(0,1]$ 内仅向 $0$ 聚集的闭区间 $I_n$，例如以 $1/n$ 为中心取足够小的区间。取连续函数 $\varphi_n$，满足 $0\le\varphi_n\le1$、支集含于 $I_n$，且某点取值 $1$。定义
\[
T(a)(x)=\sum_{n=1}^\infty a_n\varphi_n(x),\qquad a=(a_n)\in\ell^\infty.
\]
因支集互不相交，每个 $x$ 至多有一项非零，故函数在 $(0,1]$ 上连续，且
$\|T(a)\|_\infty=\sup_n|a_n|$。所以 $T$ 是从 $\ell^\infty$ 到 $C(0,1]$ 的线性等距嵌入，其像即所求子空间。

\subsection*{2.3 $C^1[a,b]$ 上的 Sobolev 型范数}
\begin{problem}
在 $C^1[a,b]$ 上定义
\[
\|f\|_1=\left(\int_a^b(|f|^2+|f'|^2)\,dx\right)^{1/2}.
\]
证明它是范数，并判断 $(C^1[a,b],\|\cdot\|_1)$ 是否完备。
\end{problem}
（1）把 $J(f)=(f,f')$ 看成到 Hilbert 空间 $L^2[a,b]\oplus L^2[a,b]$ 的线性映射，则题设量就是 $\|J(f)\|$。齐次性和三角不等式随即成立。若其为零，则 $f=0$ 几乎处处；由连续性 $f\equiv0$，故它是范数。

（2）不完备。取一个属于 $H^1[a,b]$ 但不属于 $C^1[a,b]$ 的函数，例如区间内部取 $c$，令 $f(x)=|x-c|$。用光滑函数 $f_n(x)=\sqrt{(x-c)^2+n^{-2}}$，则 $f_n\to f$ 于 $L^2$，且
\[
f_n'(x)=\frac{x-c}{\sqrt{(x-c)^2+n^{-2}}}\to \operatorname{sgn}(x-c)
\]
几乎处处并由 $1$ 控制，故导数也在 $L^2$ 中收敛。因此 $(f_n)$ 在题设范数下是 Cauchy 序列。若它在 $C^1$ 中收敛到 $g$，则 $g=f$ 几乎处处，连续性给出 $g=f$ 处处，但 $f$ 在 $c$ 不可微，矛盾。

\subsection*{2.4 两个加权 $L^2$ 范数}
\begin{problem}
在 $C[0,1]$ 上定义
\[
\|f\|_1=\left(\int_0^1|f(x)|^2dx\right)^{1/2},\qquad
\|f\|_2=\left(\int_0^1(1+x)|f(x)|^2dx\right)^{1/2}.
\]
证明这两个范数等价。
\end{problem}
因为 $1\le1+x\le2$，
\[
\|f\|_1^2\le\|f\|_2^2\le2\|f\|_1^2,
\]
故 $\|f\|_1\le\|f\|_2\le\sqrt2\|f\|_1$。两者都是由正权内积诱导的范数，因而等价。

\subsection*{2.5 指数权范数}
\begin{problem}
设 $BC[0,\infty)$ 为 $[0,\infty)$ 上的有界连续函数空间。对 $a>0$ 定义
\[
\|f\|_a=\left(\int_0^\infty e^{-ax}|f(x)|^2dx\right)^{1/2}.
\]
证明它是范数；并证明当 $a,b>0$ 且 $a\ne b$ 时，$\|\cdot\|_a$ 与 $\|\cdot\|_b$ 不等价。
\end{problem}
（1）$e^{-ax}>0$，故积分平方根满足正定性；其余范数公理由加权 $L^2$ 范数得到。

（2）设 $a<b$。总有 $\|f\|_b\le\|f\|_a$。若反向也等价，则存在 $C$ 使 $\|f\|_a\le C\|f\|_b$。取固定非零 $\varphi\in C_c[0,1]$，令 $f_n(x)=\varphi(x-n)$。换元 $u=x-n$ 得
\[
\frac{\|f_n\|_a^2}{\|f_n\|_b^2}
=e^{(b-a)n}\frac{\int_0^1e^{-au}|\varphi(u)|^2du}{\int_0^1e^{-bu}|\varphi(u)|^2du}\to\infty,
\]
矛盾。$a>b$ 时交换二者即可。

\subsection*{2.6 绝对收敛级数判别完备性}
\begin{problem}
设 $(X,\|\cdot\|)$ 为赋范空间。证明：$X$ 是 Banach 空间，当且仅当对任意 $(u_i)\subset X$，只要
$\sum_{i\ge1}\|u_i\|<\infty$，就存在 $u\in X$ 使 $u=\sum_{i\ge1}u_i$。
\end{problem}
若 $X$ 完备且 $\sum\|u_i\|<\infty$，部分和 $s_n=\sum_{i=1}^nu_i$ 满足
$\|s_m-s_n\|\le\sum_{i=n+1}^m\|u_i\|\to0$，故收敛。

反之，设每个绝对收敛级数都收敛。给定 Cauchy 序列 $(x_n)$，选子列 $(x_{n_k})$ 使
$\|x_{n_{k+1}}-x_{n_k}\|<2^{-k}$。于是
\[
x_{n_1}+\sum_{k=1}^\infty(x_{n_{k+1}}-x_{n_k})
\]
收敛，即子列 $x_{n_k}$ 收敛到某 $x$。原序列为 Cauchy，故也收敛到 $x$。因此 $X$ 完备。

\section{第三次作业}

\subsection*{3.1 上确界范数不是内积范数}
\begin{problem}
证明：在 $C[a,b]$ 上不可能定义内积 $(\cdot,\cdot)$，使
$(f,f)^{1/2}=\max_{a\le x\le b}|f(x)|$ 对所有 $f\in C[a,b]$ 成立。
\end{problem}
内积诱导的范数必须满足平行四边形恒等式。取 $f(x)=1$，并取非恒定连续函数 $g$ 满足 $\|g\|_\infty=1$ 且值域含 $-1,1$（例如 $g$ 为区间的仿射标准化）。则
$\|f\|=\|g\|=1$，而 $\|f+g\|=\|f-g\|=2$。于是左边为 $8$，右边 $2\|f\|^2+2\|g\|^2=4$，矛盾。

\subsection*{3.2 $L^2$ 泛函的最大值}
\begin{problem}
在 $L^2[0,T]$ 的单位球面上求泛函
\[
f\longmapsto\int_0^T e^{-(T-t)}f(t)\,dt
\]
的最大值，并求达到最大值的函数 $f$。
\end{problem}
令 $h(t)=e^{-(T-t)}$。由 Cauchy--Schwarz，单位球面上
\[
\left|\int_0^Th(t)f(t)dt\right|\le\|h\|_2
=\left(\int_0^Te^{-2(T-t)}dt\right)^{1/2}
=\sqrt{\frac{1-e^{-2T}}2}.
\]
等号当且仅当 $f$ 与 $h$ 成比例。若求泛函本身的最大值，取
\[
f(t)=\frac{e^{-(T-t)}}{\sqrt{(1-e^{-2T})/2}}.
\]
其 $L^2$ 范数为 $1$，且达到上述正值；负号给出最小值。

\subsection*{3.3 $L^\infty$ 不可分}
\begin{problem}
证明 $L^\infty$ 空间不是可分度量空间。
\end{problem}
在 $[0,1]$ 上，对每个二进制序列 $a=(a_n)$，令
$E_a=\bigcup_{a_n=1}(2^{-n},2^{-n+1}]$，并令 $f_a=\mathbf1_{E_a}$。若 $a\ne b$，某个区间的示性值不同，故
$\|f_a-f_b\|_\infty=1$。这样得到不可数个两两距离为 $1$ 的点。若空间可分，取可数稠密集，则半径 $1/3$ 的球必须覆盖空间，但每个球至多包含一个 $f_a$，不可能覆盖不可数族。因此不可分。一般非原子测度空间或含无限个不交正测集的情形同理。

\subsection*{3.4 Hanner 不等式与一致凸性}
\begin{problem}
证明讲义第 23 页的 Hanner 不等式，并由此说明：当 $1<p<\infty$ 时，$L^p$ 空间一致凸。
\end{problem}
Hanner 不等式为：当 $1\le p\le2$ 时
\[
\|x+y\|_p^p+\|x-y\|_p^p
\ge(\|x\|_p+\|y\|_p)^p+|\|x\|_p-\|y\|_p|^p,
\]
当 $2\le p<\infty$ 时不等号反向。其核心是一维 Hanner 引理：利用齐次性令 $a=1$、$0\le b\le1$，把两边之差记为 $G(b)$；对 $b$ 求导后，$G'(b)$ 的符号由 $t\mapsto t^{p-1}$ 的凸性（$p\ge2$）或凹性（$p\le2$）确定，再由端点 $G(1)=0$ 得所需方向。将该标量结论逐点应用，积分后用 Minkowski 不等式，便得到上述函数空间版本；复数情形将复数视作二维实向量即可。

下面详细说明它如何推出一致凸。设 $\|x\|_p=\|y\|_p=1$ 且 $\|x-y\|_p\ge\varepsilon$。若 $p\ge2$，由 Hanner 不等式
\[
\|x+y\|_p^p+\|x-y\|_p^p\le2^p,
\]
故
\[
\left\|\frac{x+y}{2}\right\|_p
\le\left(1-(\varepsilon/2)^p\right)^{1/p}<1.
\]
若 $1<p<2$，令 $q=p/(p-1)$。由 Hanner 不等式等价得到 Clarkson 估计
\[
\left\|\frac{x+y}{2}\right\|_p^q+
\left\|\frac{x-y}{2}\right\|_p^q
\le\left(\frac{\|x\|_p^p+\|y\|_p^p}{2}\right)^{q/p}=1.
\]
所以 $\|(x+y)/2\|_p\le(1-(\varepsilon/2)^q)^{1/q}<1$。两种情形都给出只依赖于 $p,\varepsilon$ 的严格缺口，故 $L^p$ 一致凸。

\subsection*{3.5 $\sigma$-有限完备测度是可局部化的}
\begin{problem}
证明：若 $(X,\mathcal M,\mu)$ 是完备的 $\sigma$-有限测度空间，则 $\mu$ 是可局部化的。
\end{problem}
采用常用定义：若 $\mathcal E\subset\mathcal M$，则存在 $H\in\mathcal M$，使每个 $E\in\mathcal E$ 都满足 $E\setminus H$ 为零测集，且任何具有此性质的 $G$ 都满足 $H\setminus G$ 为零测集。

由 $\sigma$-有限性，取 $X=\bigcup_nX_n$ 且 $\mu(X_n)<\infty$，并可令 $X_n$ 两两不交。固定 $n$，令
\[
s_n=\sup\{\mu(E\cap X_n):E\in\mathcal E\}.
\]
取 $E_{n,k}\in\mathcal E$ 使 $\mu(E_{n,k}\cap X_n)\to s_n$，置
$H_n=X_n\cap\bigcup_kE_{n,k}$，$H=\bigcup_nH_n$。对任意 $E\in\mathcal E$，若
$\mu((E\cap X_n)\setminus H_n)>0$，则把它加入某个近极大集合会使测度超过 $s_n$，矛盾；故 $E\setminus H$ 为零测集。若 $G$ 也是上界，则每个 $E_{n,k}\setminus G$ 为零集；可列并仍是零集，所以 $H\setminus G$ 是零集。完备性保证所出现零集的子集仍可测，故得到可局部化性。

\subsection*{3.6}
\begin{problem}
证明讲义第 18 页中的结论 (i)、(ii)。
\end{problem}
\textbf{依赖未提供讲义，按用户要求跳过。}题面仅称“第 18 页中的 (i)(ii)”，没有给出结论。

\section{第四次作业}

\subsection*{4.1 Riemann 和作为网收敛}
\begin{problem}
查阅资料写出网（net）及网收敛的定义，并说明区间 $[a,b]$ 上 Riemann 可积可以表述为一族 Riemann 和的网收敛。
\end{problem}
有向集 $(A,\preceq)$ 是满足任意两元素都有共同上界的预序集。网是映射 $A\to X$，记作 $(x_\alpha)$；若对 $x$ 的每个邻域 $U$，存在 $\alpha_0$ 使 $\alpha\succeq\alpha_0$ 时 $x_\alpha\in U$，则称网收敛于 $x$。

令 $A$ 为 $[a,b]$ 的所有带标记分割 $(P,\xi)$。规定 $(P',\xi')\succeq(P,\xi)$ 当 $P'$ 细分 $P$；也可用“网格直径更小”的预序。对应网为 Riemann 和
\[
S(P,\xi;f)=\sum_i f(\xi_i)(x_i-x_{i-1}).
\]
$f$ Riemann 可积且积分为 $I$，正是说对每个 $\varepsilon>0$，存在 $\delta>0$，使所有网格小于 $\delta$ 的带标记分割都满足 $|S-I|<\varepsilon$；这恰是上述 Riemann 和网收敛于 $I$。

\subsection*{4.2 两个距离问题}
\begin{problem}
在装备上确界范数的 $C[0,1]$ 中令
\[
N=\left\{f:\int_0^1f(x)dx=0\right\},\quad
X=\{f:f(0)=0\},\quad M=N\cap X.
\]
证明：
\begin{enumerate}[label=(\arabic*)]
\item 对任意 $u\in C[0,1]$，若 $\bar u=\int_0^1u(x)dx$，则
$d(u,N)=|\bar u|$，且下确界在 $n=u-\bar u$ 处达到；
\item 对 $u(x)=x\in X$，有 $d(x,M)=1/2$，但该下确界不在任何 $m\in M$ 处达到。
\end{enumerate}
\end{problem}
（1）令 $\bar u=\int_0^1u(x)dx$。对任意 $n\in N$，
\[
|\bar u|=\left|\int_0^1(u-n)\right|\le\|u-n\|_\infty,
\]
故 $d(u,N)\ge|\bar u|$。另一方面 $n_0=u-\bar u\in N$，且
$\|u-n_0\|_\infty=|\bar u|$，故等号成立且下确界达到。

（2）对 $m\in M$，令 $h=x-m$。则 $h(0)=0$ 且 $\int_0^1h=1/2$，所以
$1/2=|\int h|\le\|h\|_\infty$，得到下界。构造趋近者：给定 $\varepsilon>0$，取连续 $h_\varepsilon$，满足 $h_\varepsilon(0)=0$，在 $[\varepsilon,1]$ 上接近常数 $c_\varepsilon$，并调节 $c_\varepsilon$ 使 $\int h_\varepsilon=1/2$；例如令
\[
h_\varepsilon(x)=c_\varepsilon\min(x/\varepsilon,1),\qquad
c_\varepsilon=\frac1{2(1-\varepsilon/2)}.
\]
则 $\|h_\varepsilon\|_\infty=c_\varepsilon\downarrow1/2$，而 $m_\varepsilon=x-h_\varepsilon\in M$，故距离为 $1/2$。若某 $m$ 达到下界，则 $h=x-m$ 连续、$h(0)=0$、$\|h\|_\infty=1/2$ 且 $\int h=1/2$。积分达到上界迫使 $h(x)=1/2$ 处处，与 $h(0)=0$ 矛盾，故不达到。

\subsection*{4.3 正交补}
\begin{problem}
若 $A$ 是 Hilbert 空间的任意子集，证明 $A^\perp=(\overline A)^\perp$。若 $M$ 是 Hilbert 空间的线性子空间，证明 $M^{\perp\perp}=\overline M$。
\end{problem}
若 $x\in A^\perp$，则内积函数 $y\mapsto\langle x,y\rangle$ 连续，故它在 $\overline A$ 上仍为零，因而 $x\in(\overline A)^\perp$；反向包含显然，所以 $A^\perp=(\overline A)^\perp$。

对线性子空间 $M$，显然 $\overline M\subset M^{\perp\perp}$。反之，若 $x\notin\overline M$，由 Hilbert 投影定理，$x=P_{\overline M}x+z$，其中 $z\ne0$ 且 $z\in M^\perp$。于是 $\langle x,z\rangle=\|z\|^2\ne0$，故 $x\notin M^{\perp\perp}$。因此 $M^{\perp\perp}=\overline M$。

\subsection*{4.4 Hilbert 直和}
\begin{problem}
设 $H_1,H_2$ 为 Hilbert 空间。在
$H_1\oplus H_2=\{(x_1,x_2):x_i\in H_i\}$ 上定义
\[
\langle(x_1,x_2),(y_1,y_2)\rangle
=\langle x_1,y_1\rangle_{H_1}+\langle x_2,y_2\rangle_{H_2}.
\]
证明 $H_1\oplus H_2$ 是 Hilbert 空间，并求子空间 $\{(x_1,0):x_1\in H_1\}$ 的正交补。
\end{problem}
若 $(x_1^{(n)},x_2^{(n)})$ 为 Cauchy 序列，则
$\|x_i^{(n)}-x_i^{(m)}\|\le\|(x_1^{(n)},x_2^{(n)})-(x_1^{(m)},x_2^{(m)})\|$，故两个坐标各自 Cauchy，并分别收敛到 $x_i\in H_i$。平方范数相加表明原序列收敛到 $(x_1,x_2)$，故直和完备。

若 $(y_1,y_2)$ 与所有 $(x_1,0)$ 正交，则 $\langle y_1,x_1\rangle=0$ 对所有 $x_1$ 成立，取 $x_1=y_1$ 得 $y_1=0$。所以正交补恰为 $\{(0,y_2):y_2\in H_2\}$。

\subsection*{4.5 可数正交直和}
\begin{problem}
设 $(H_n)_{n\ge1}$ 是 Hilbert 空间 $H$ 中两两正交的闭子空间，并定义
\[
\bigoplus_{n=1}^\infty H_n
=\left\{\sum_{n=1}^\infty x_n:x_n\in H_n, 
\sum_{n=1}^\infty\|x_n\|^2<\infty\right\}.
\]
证明 $\bigoplus_{n=1}^\infty H_n$ 是 $H$ 的闭线性子空间。
\end{problem}
令 $M$ 为所有可写成 $\sum_nx_n$、$x_n\in H_n$ 且 $\sum_n\|x_n\|^2<\infty$ 的元素。正交性给出勾股恒等式
\[
\left\|\sum_{n=1}^Nx_n\right\|^2=\sum_{n=1}^N\|x_n\|^2,
\]
所以级数收敛，且 $M$ 对加法、数乘封闭。

若 $z_k=\sum_nx_n^{(k)}\in M$ 且 $z_k\to z$，则闭子空间 $H_n$ 的正交投影 $P_n$ 连续，故
$x_n^{(k)}=P_nz_k\to P_nz=:x_n$。对每个 $N$，
\[
\sum_{n=1}^N\|x_n\|^2=\left\|\sum_{n=1}^NP_nz\right\|^2\le\|z\|^2.
\]
令 $N\to\infty$ 得平方可和。又 $s_N=\sum_{n=1}^Nx_n$ 收敛到某 $s\in M$。对每个 $n$，$P_n(z-s)=0$；而 $z_k\in M$ 且 $z_k\to z$ 可进一步用
$\|z-s_N\|\le\|z-z_k\|+\|z_k-\sum_{n\le N}x_n^{(k)}\|+\sum_{n\le N}\|x_n^{(k)}-x_n\|$
先选 $k$、再选 $N$ 得到 $s_N\to z$。故 $z=s\in M$，所以 $M$ 闭。

\subsection*{4.6 $L^p$ 的凸投影定理}
\begin{problem}
写出并证明 $L^p$ 空间版本的凸投影定理：当 $1<p<\infty$ 时，对 $L^p$ 中每个非空闭凸集 $C$ 及每个 $f\in L^p$，存在唯一的 $g\in C$ 使 $\|f-g\|_p=d(f,C)$。
\end{problem}
设 $1<p<\infty$，$C\subset L^p$ 非空、闭、凸。则对每个 $f\in L^p$，存在唯一 $g\in C$ 使
$\|f-g\|_p=d(f,C)$。

取极小化序列 $g_n\in C$。$L^p$ 一致凸。设 $r=d(f,C)$；若 $r=0$，闭性立即给出 $f\in C$。若 $r>0$ 而 $(g_n)$ 非 Cauchy，则存在 $\varepsilon>0$ 及子列使 $\|g_n-g_m\|\ge\varepsilon$。归一化 $u_n=(f-g_n)/\|f-g_n\|$，当 $n,m$ 足够大时，二者范数为 $1$ 且相距有统一正下界。一致凸性给
\[
\left\|f-\frac{g_n+g_m}{2}\right\|
\le(1-\delta)\frac{\|f-g_n\|+\|f-g_m\|}{2}<r,
\]
与中点属于 $C$ 及 $r$ 的定义矛盾。故 $g_n$ Cauchy，完备性与闭性给极限 $g\in C$。唯一性同样由一致凸性：两个不同极小点的中点会更近。

\subsection*{4.7 无条件收敛但不绝对收敛}
\begin{problem}
设 $(x_n)$ 为 Hilbert 空间中的标准正交基。证明级数
$\sum_{n=1}^\infty x_n/n$ 无条件收敛，但不绝对收敛。
\end{problem}
设 $(x_n)$ 为标准正交基。对任意有限集 $F$，
\[
\left\|\sum_{n\in F}\frac{x_n}{n}\right\|^2=\sum_{n\in F}\frac1{n^2}.
\]
因此按任意排列的尾和都由 $\sum1/n^2$ 的尾项控制，级数无条件收敛。但
$\sum\|x_n/n\|=\sum1/n=\infty$，故不绝对收敛。

\subsection*{4.8 可分性与可数正交基}
\begin{problem}
证明：Hilbert 空间作为度量空间可分，当且仅当它具有可数标准正交基。
\end{problem}
若有可数标准正交基 $(e_n)$，则所有有限有理（复数情形取有理实部虚部）线性组合构成可数稠密集，故 $H$ 可分。

反之，设 $H$ 可分。取可数稠密集 $(x_n)$，依次删去落在此前线性张成闭包中的点，并对余下序列作 Gram--Schmidt，得到可数正交规范族 $(e_n)$。其闭线性张成包含每个 $x_n$，故等于 $H$；所以它是可数标准正交基。

\subsection*{4.9 稠密子空间中的正交基}
\begin{problem}
若 $M$ 是可分 Hilbert 空间 $H$ 的稠密线性子空间，证明 $H$ 有一组完全由 $M$ 中元素组成的标准正交基。若 $M$ 仅为 $H$ 的任意稠密子集，结论是否仍成立？
\end{problem}
取 $M$ 中可数稠密集并按上一题作 Gram--Schmidt。因为 $M$ 是线性子空间，每次有限线性组合仍在 $M$，故所得正交基各元素都属于 $M$。

若 $M$ 只是假定为任意稠密子集，结论不成立。在 $\mathbb R^2$ 中令
\[
D=\{(\cos(2\pi q\sqrt2),\sin(2\pi q\sqrt2)):q\in\mathbb Q\},\quad
M=\{rd:r\in\mathbb Q_{>0},d\in D\}.
\]
$D$ 在单位圆上稠密，故 $M$ 在 $\mathbb R^2$ 稠密。若 $D$ 中两方向正交，则
$2\pi(q-r)\sqrt2=\pi/2+k\pi$，从而 $(q-r)\sqrt2=1/4+k/2$；左边为零或无理数，右边是非零有理数，矛盾。因此 $M$ 不含正交基。

\section{第五次作业}

\subsection*{5.1}
\begin{problem}
证明讲义第 29 页中的命题 9。
\end{problem}
\textbf{依赖未提供讲义，按用户要求跳过。}题面没有给出“第 29 页命题 9”的陈述。

\subsection*{5.2A Gram--Schmidt 与正交多项式}
\begin{problem}
设 $H$ 为 Hilbert 空间且所有多项式均属于 $H$。称 $(P_n)_{n\ge0}$ 为正交多项式系，如果
$\mathcal P=\operatorname{span}\{1,x,x^2,\ldots\}\subset H$、$\deg P_n=n$，且
$(P_m,P_n)=0$（$m\ne n$）。证明：
\begin{enumerate}[label=(\arabic*)]
\item 正交多项式系存在；
\item 若 $(P_n)$ 与 $(\widetilde P_n)$ 是两个正交多项式系，则对每个 $n$ 存在非零标量 $c_n$ 使 $P_n=c_n\widetilde P_n$。
\end{enumerate}
\end{problem}
对 $1,x,x^2,\ldots$ 作 Gram--Schmidt：
\[
P_0=1,\qquad P_n=x^n-\sum_{k<n}\frac{(x^n,P_k)}{(P_k,P_k)}P_k.
\]
减去项次数均小于 $n$，所以 $\deg P_n=n$，构造式又保证两两正交，正交多项式存在。若 $(P_n)$ 与 $(\widetilde P_n)$ 是两族，则二者的前 $n$ 项都张成 $\mathcal P_{n-1}$，而 $P_n,\widetilde P_n$ 都属于一维空间 $\mathcal P_n\cap\mathcal P_{n-1}^{\perp}$，故 $P_n=c_n\widetilde P_n$，$c_n\ne0$。

\subsection*{5.2B Rodrigues 公式}
\begin{problem}
考虑 Sturm--Liouville 方程
\[
p(x)y''+q(x)y'+\lambda y=0,\qquad
p(x)=\alpha x^2+\beta x+\gamma,quad q(x)=\mu x+\nu.
\]
证明：
\begin{enumerate}[label=(\arabic*)]
\item 若 $y_n=\sum_{j=0}^ng_jx^j$ 是方程的 $n$ 次多项式解，则
$[n(n-1)\alpha+n\mu+\lambda]g_n=0$；
\item 若 $w$ 满足 $(wp)'=wq$，则方程可写成
$(wpy')'+\lambda wy=0$；
\item Rodrigues 公式
\[
y_n(x)=\frac1{w(x)}\frac{d^n}{dx^n}\bigl(w(x)p(x)^n\bigr)
\]
在 $\lambda=-\frac{n(n-1)}2p''-nq'$ 时给出方程的解；
\item 验证三种经典正交多项式的 Rodrigues 公式。
\end{enumerate}
\end{problem}
若 $y_n=g_nx^n+$低次项，比较 $py''+qy'+\lambda y=0$ 的 $x^n$ 系数，得
$[n(n-1)\alpha+n\mu+\lambda]g_n=0$。若 $(wp)'=wq$，乘 $w$ 即得
$(wpy')'+\lambda wy=0$。令 $v=wp^n$、$y=w^{-1}v^{(n)}$，用 $p'''=q''=0$ 和 Leibniz 公式整理得
\[
(wpy')'+\lambda wy=
\left(\lambda+nq'+\frac{n(n-1)}2p''\right)v^{(n)}.
\]
故题设 $\lambda$ 下它为解。三种典型选取
\[
(p,q,w)=(1-x^2,-2x,1),\quad(x,1-x,e^{-x}),\quad(1,-2x,e^{-x^2})
\]
分别给出 Legendre、Laguerre、Hermite 的 Rodrigues 公式。

\subsection*{5.3 勒让德多项式}
\begin{problem}
定义 Legendre 多项式
\[
P_n(x)=\frac1{2^nn!}\frac{d^n}{dx^n}(x^2-1)^n.
\]
证明：
\begin{enumerate}[label=(\arabic*)]
\item $(P_n)$ 在 $L^2[-1,1]$ 中两两正交，且可由 $1,x,x^2,\ldots$ 作 Gram--Schmidt 得到；
\item $\int_{-1}^1P_n(x)^2dx=2/(2n+1)$；
\item $(P_n)$ 是 $L^2[-1,1]$ 的正交基；对 $f\in L^2[-1,1]$ 求
$f=\sum_{n\ge0}c_nP_n$ 的 $c_n$，并说明收敛意义；
\item $P_n$ 是算子 $L=-\frac d{dx}(1-x^2)\frac d{dx}$ 的特征函数，特征值为 $n(n+1)$。
\end{enumerate}
\end{problem}
若 $m<n$，Rodrigues 公式分部积分 $n$ 次得
\[
\int_{-1}^1x^mP_n(x)dx=0.
\]
边界项消失，因为 $(x^2-1)^n$ 在端点连同前 $n-1$ 阶导数均为零，且 $D^nx^m=0$。故 $P_n$ 正交于一切低次多项式，亦即是 Gram--Schmidt 的结果（差归一化）。

为计算范数，再把一个 $P_n$ 用 Rodrigues 公式表示并分部积分 $n$ 次。其首项系数及最高阶导数分别为
\[
\frac{(2n)!}{2^n(n!)^2},\qquad P_n^{(n)}=\frac{(2n)!}{2^nn!}.
\]
于是
\[
\int_{-1}^1P_n^2
=\frac{(2n)!}{2^{2n}(n!)^2}\int_{-1}^1(1-x^2)^n dx.
\]
而 Beta 积分给出
$\int_{-1}^1(1-x^2)^n dx=2^{2n+1}(n!)^2/(2n+1)!$，所以
\[
\int_{-1}^1P_n^2=\frac2{2n+1}.
\]
多项式在 $L^2[-1,1]$ 稠密，且 $P_0,\ldots,P_n$ 张成 $\mathcal P_n$，故它们是正交基，并且
\[
f=\sum_{n\ge0}c_nP_n,\qquad c_n=\frac{2n+1}{2}\int_{-1}^1fP_n,
\]
在 $L^2$ 范数意义下收敛。Legendre 方程的自伴随形式为
$-D((1-x^2)P_n')=n(n+1)P_n$，故特征值为 $n(n+1)$。

\subsection*{5.4 切比雪夫多项式}
\begin{problem}
在加权 Hilbert 空间
\[
H=\left\{f:[-1,1]\to\mathbb C:
\int_{-1}^1\frac{|f(x)|^2}{\sqrt{1-x^2}}dx<\infty\right\}
\]
中取自然内积。定义 Chebyshev 多项式
$T_n(x)=\cos(n\theta)$，其中 $x=\cos\theta$、$0\le\theta\le\pi$。证明 $(T_n)_{n\ge0}$ 是 $H$ 的正交基，并且
$\|T_0\|=\sqrt\pi$、$\|T_n\|=\sqrt{\pi/2}$（$n\ge1$）。
\end{problem}
令 $x=\cos\theta$，则
\[
\langle T_m,T_n\rangle=\int_0^\pi\cos(m\theta)\cos(n\theta)d\theta.
\]
所以不同下标正交，$\|T_0\|=\sqrt\pi$，$\|T_n\|=\sqrt{\pi/2}$（$n\ge1$）。若 $f$ 与全部 $T_n$ 正交，则 $f(\cos\theta)$ 的偶延拓与所有 Fourier 指数正交，故为零；所以它们完备。

\subsection*{5.5 Hermite 系统}
\begin{problem}
定义 Hermite 多项式
\[
H_n(x)=(-1)^ne^{x^2}\frac{d^n}{dx^n}e^{-x^2},\qquad
\phi_n(x)=e^{-x^2/2}H_n(x).
\]
证明：
\begin{enumerate}[label=(\arabic*)]
\item $(\phi_n)_{n\ge0}$ 是 $L^2(\mathbb R)$ 的正交基；
\item $\phi_n$ 是 $-d^2/dx^2+x^2$ 的特征函数，特征值为 $2n+1$；提示：使用 $A=d/dx+x$ 与 $A^*=-d/dx+x$；
\item 生成函数满足
\[
G(x,t):=\sum_{n=0}^\infty H_n(x)\frac{t^n}{n!}
=e^{x^2-(x-t)^2}=e^{2xt-t^2};
\]
\item 对酉 Fourier 变换
$\mathcal F(f)(k)=(2\pi)^{-1/2}\int_{\mathbb R}f(x)e^{-ixk}dx$，证明
$\mathcal F(\phi_n)=(-i)^n\phi_n$。
\end{enumerate}
\end{problem}
Rodrigues 公式分部积分说明 $\phi_n=e^{-x^2/2}H_n$ 两两正交。若 $\varphi\perp\phi_n$，则
$\int\varphi x^ne^{-x^2/2}=0$。整函数
$F(z)=\int\varphi(x)e^{-x^2/2}e^{-izx}dx$ 的所有导数在 $0$ 为零，故 $F\equiv0$；Fourier 变换单射性给出 $\varphi=0$，所以系统完备。

令 $A=D+x$、$A^*=-D+x$，则 $-D^2+x^2=A^*A+1$。由
$A\phi_n=2n\phi_{n-1}$、$A^*\phi_{n-1}=\phi_n$ 得
$(-D^2+x^2)\phi_n=(2n+1)\phi_n$。Taylor 展开给出
\[
G(x,t)=e^{x^2-(x-t)^2}=e^{2xt-t^2}=\sum_{n\ge0}H_n(x)t^n/n!.
\]
对 $e^{-x^2/2}G$ 作酉 Fourier 变换并比较 $t^n$ 系数，得到
$\mathcal F\phi_n=(-i)^n\phi_n$。

\section{第六次作业}

\subsection*{6.1 $C^1(I)$ 的嵌入}
\begin{problem}
令 $I=[0,1]$，在 $C^1(I)$ 上定义
$\|f\|_{C^1}=\sup_I|f|+\sup_I|f'|$。证明：
\begin{enumerate}[label=(\arabic*)]
\item $C^1(I)$ 在此范数下是 Banach 空间；
\item 嵌入 $\iota:C^1(I)\to C(I)$ 是有界线性算子；
\item 若 $B$ 是 $C^1(I)$ 的单位球，则 $\iota(B)$ 在 $C(I)$ 中的闭包紧；
\item $\iota(B)$ 是否为 $C(I)$ 中的闭集？
\item $\iota(C^1(I))$ 是否为 $C(I)$ 中的紧集？
\end{enumerate}
\end{problem}
若 $(f_n)$ 在 $C^1$ 范数下 Cauchy，则 $f_n\to f$、$f_n'\to g$ 一致。由
$f_n(t)-f_n(0)=\int_0^tf_n'$ 取极限，得 $f'=g$，故完备。又
$\|\iota f\|_\infty\le\|f\|_{C^1}$，嵌入有界。单位球满足 $|f|\le1$ 及
$|f(t)-f(s)|\le|t-s|$，故像族一致有界且等度连续，其闭包由 Arzelà--Ascoli 定理紧。

单位球的像不是闭集：$|t-1/2|/2$ 可由适当缩放的光滑函数在单位球内一致逼近，却不属于 $C^1$。整个像也不紧，因为它包含任意大的常数函数，因而无界。

\subsection*{6.2 连续核积分算子}
\begin{problem}
令 $K:I\times I\to\mathbb R$ 连续，并定义
\[
(T_Kf)(t)=\int_0^1K(t,s)f(s)ds.
\]
证明 $T_K:C(I)\to C(I)$ 连续；若 $B$ 是 $C(I)$ 的单位球，证明 $T_K(B)$ 的闭包紧。
\end{problem}
\[
\|T_Kf\|_\infty\le\max|K|\,\|f\|_\infty,
\quad |T_Kf(t)-T_Kf(t')|\le\sup_s|K(t,s)-K(t',s)|.
\]
第一式给连续性；第二式结合 $K$ 的一致连续性，说明单位球的像等度连续且一致有界，故闭包紧。

\subsection*{6.3 严格压缩}
\begin{problem}
设 $(X,\rho)$ 为完备度量空间，$M\subset X$ 列紧，且 $f:X\to M$ 满足
\[
\rho(f(x_1),f(x_2))<\rho(x_1,x_2)\qquad(x_1\ne x_2).
\]
证明 $f$ 在 $X$ 中存在唯一不动点。
\end{problem}
令 $x_{n+1}=f(x_n)$，并记 $a_n=\rho(x_{n+1},x_n)$，则 $a_n$ 严格递减或已为零，故趋于 $a\ge0$。因 $x_n\in M$，取子列 $x_{n_k}\to p$，连续性给 $x_{n_k+1}\to f(p)$，故 $a=\rho(p,f(p))$。若 $p\ne f(p)$，再应用严格压缩得
$\rho(f(p),f^2(p))<a$，与相邻距离均趋于 $a$ 矛盾。因此 $p$ 是不动点。两个不动点会直接违反严格压缩，故唯一。

\subsection*{6.4 不定积分族}
\begin{problem}
设 $M$ 是 $C[a,b]$ 中的有界集。证明函数族
\[
\left\{F:F(x)=\int_a^xf(t)dt, f\in M\right\}
\]
是 $C[a,b]$ 中的预列紧集。
\end{problem}
若 $\|f\|_\infty\le C$，则 $|F(x)|\le C(b-a)$ 且
$|F(x)-F(y)|\le C|x-y|$，故由 Arzelà--Ascoli 定理预列紧。

\subsection*{6.5 正弦函数族}
\begin{problem}
证明 $E=\{\sin(nt):n\ge1\}$ 在 $C[0,\pi]$ 中不是列紧集。
\end{problem}
取 $s_n=\pi/(2n)$、$t_n=0$，距离趋零而函数值差恒为 $1$，所以任何频率趋于无穷的子列都不等度连续，因而无一致收敛子列。
更具体地，假设存在一致收敛子列 $\sin(n_kt)$。一致收敛序列与其连续极限合在一起必等度连续：给定 $\varepsilon>0$，先用极限函数的一致连续性控制尾项，再逐个控制有限多个首项。可是令
$s_k=\pi/(2n_k)$、$t_k=0$，则 $|s_k-t_k|\to0$，而
\[
|\sin(n_ks_k)-\sin(n_kt_k)|=|1-0|=1,
\]
与等度连续矛盾。因此原集合中序列 $(\sin nt)$ 没有一致收敛子列，$E$ 不是列紧集。

\subsection*{6.6 完全有界与有界}
\begin{problem}
证明度量空间中的完全有界集一定有界；并以 $\ell^2$ 的标准基集合
$E=\{e_k:k\ge1\}$ 说明有界集未必完全有界。
\end{problem}
用有限个单位球覆盖完全有界集，即得其有界。$\ell^2$ 标准基有界，而
$\|e_j-e_k\|=\sqrt2$，故不能由有限个半径小于 $\sqrt2/2$ 的球覆盖，不完全有界。
详细地说，若 $A\subset X$ 完全有界，取 $\varepsilon=1$，存在 $x_1,\ldots,x_N$ 使
$A\subset\bigcup_{j=1}^NB(x_j,1)$。固定原点 $x_0$ 并令
$R=1+\max_jd(x_j,x_0)$，则任意 $x\in A$ 均满足 $d(x,x_0)<R$，所以 $A$ 有界。
另一方面，$\|e_k\|_2=1$，所以标准基集合有界。若它可由有限个半径 $1/2$ 的球覆盖，由抽屉原理某个球包含不同的 $e_j,e_k$，于是
$\sqrt2=\|e_j-e_k\|_2\le\|e_j-c\|_2+\|e_k-c\|_2<1$，矛盾。因此它不完全有界。

\subsection*{6.7 可分情形与依赖选择}
\begin{problem}
说明：对于可分的非空完备度量空间，Baire 纲定理的证明无需使用依赖选择公理。
\end{problem}
固定可数稠密集。嵌套球构造的每一步都按编号取第一个合适中心及第一个合适有理半径，便成为确定规则，不需依赖选择；完备性保证球心 Cauchy 序列收敛。

\subsection*{6.8 超越数稠密}
\begin{problem}
利用 Baire 纲定理证明超越数在 $\mathbb R$ 中稠密。
\end{problem}
代数数可数，枚举为 $(a_n)$。每个 $\mathbb R\setminus\{a_n\}$ 开且稠密，故其交由 Baire 定理稠密；该交正是超越数集合。
确切地，对任意非空开区间 $I$，把 $I$ 看成 $\mathbb R$ 的开子空间。集合
$I\setminus\{a_n\}$ 在 $I$ 中开且稠密。Baire 定理给出
\[
I\cap\bigcap_{n=1}^\infty(\mathbb R\setminus\{a_n\})\ne\varnothing.
\]
所以每个开区间都含有超越数，这正说明超越数在 $\mathbb R$ 中稠密。

\section{第七次作业}

\subsection*{7.1 Baire 纲定理与依赖选择}
\begin{problem}
设 $X$ 为非空集合，映射 $A:X\to2^X$ 满足 $A(x)\ne\varnothing$。证明存在序列 $(x_n)\subset X$ 使 $x_{n+1}\in A(x_n)$。提示：在 $\mathcal X=X^{\mathbb N}$ 上定义
\[
d(\xi,\eta)=
\begin{cases}
0,&\xi=\eta,\\
2^{-n},&n=\min\{k:x_k\ne y_k\},
\end{cases}
\]
证明 $\mathcal X$ 完备；再令
\[
U_k=\{(x_n):\text{存在 }l>k\text{ 使 }x_l\in A(x_k)\},
\]
证明各 $U_k$ 开且稠密，并从 $\bigcap_kU_k$ 的元素构造所求序列。
\end{problem}
在 $\mathcal X=X^{\mathbb N}$ 上定义题设距离。三角不等式来自“两个序列第一次不同的位置不早于分别与第三个序列第一次不同位置的较小者”。若 $(\xi^r)$ 为 Cauchy 序列，则对每个 $k$，当 $r$ 足够大时其前 $k$ 项稳定；把各稳定坐标拼成 $\xi$，即有 $\xi^r\to\xi$，故 $\mathcal X$ 完备。

$U_k$ 是开集，因为“某个 $l>k$ 满足 $x_l\in A(x_k)$”只依赖有限个坐标。它也稠密：给定任意基本柱集，即指定前 $N$ 个坐标，若 $k>N$ 可继续任意填写到第 $k$ 项，再在某个 $l>k$ 填入 $A(x_k)$ 的一个元素；若 $k\le N$，直接在 $N$ 后某个坐标填入 $A(x_k)$ 的元素。所得序列仍在原柱集中且属于 $U_k$。

由 Baire 定理，取 $\xi=(x_j)\in\bigcap_kU_k$。递归取 $n_1=1$；已取 $n_r$ 后，由 $\xi\in U_{n_r}$，可取 $n_{r+1}>n_r$ 使
$x_{n_{r+1}}\in A(x_{n_r})$。令 $y_r=x_{n_r}$，则 $y_{r+1}\in A(y_r)$，得到依赖选择所要求的序列。

\subsection*{7.2 从离散射线收敛推出无穷远收敛}
\begin{problem}
设连续函数 $f:[0,\infty)\to\mathbb R$ 满足：对每个 $t>0$，
$\lim_{n\to\infty}f(nt)=0$。证明 $\lim_{x\to\infty}f(x)=0$。提示：固定 $\varepsilon>0$，考虑
\[
A_n=\{t>0:|f(mt)|\le\varepsilon, \forall m\ge n\}.
\]
\end{problem}
固定 $\varepsilon>0$，令
\[
A_n=\{t>0:|f(mt)|\le\varepsilon\text{ 对所有 }m\ge n\}.
\]
由假设 $\bigcup_nA_n=(0,\infty)$。每个 $A_n$ 是闭集，因为它是闭集
$\{|f(mt)|\le\varepsilon\}$ 的可数交。$(0,\infty)$ 是 Baire 空间，故某个 $A_n$ 有非空内点；取闭区间 $[a,b]\subset A_n$。把 $n$ 增大而不改变包含关系，可设 $n(b-a)\ge a$。

对任意 $x\ge na$，区间 $[x/b,x/a]$ 的长度足以包含一个整数 $m\ge n$；等价地可取 $m$ 使 $x/m\in[a,b]$。因 $x/m\in A_n$ 且 $m\ge n$，有
$|f(x)|=|f(m(x/m))|\le\varepsilon$。所以 $f(x)\to0$。

\subsection*{7.3 线性泛函连续性的判别}
\begin{problem}
设 $X$ 为无限维赋范空间，$\Lambda:X\to\mathbb R$ 为非零线性泛函。证明以下条件等价：
\begin{enumerate}[label=(\arabic*)]
\item $\Lambda$ 有界；
\item $\ker\Lambda$ 是闭线性子空间；
\item $\ker\Lambda$ 在 $X$ 中不稠密。
\end{enumerate}
\end{problem}
若 $\Lambda$ 有界，则连续，故 $\ker\Lambda=\Lambda^{-1}(\{0\})$ 闭。若核闭且 $\Lambda\ne0$，取 $x_0$ 使 $\Lambda x_0=1$。到闭集的距离
$d(x_0,\ker\Lambda)=c>0$。对任意 $x$，有 $x-\Lambda(x)x_0\in\ker\Lambda$，所以当 $\Lambda(x)\ne0$ 时
\[
c\le d(x_0,\ker\Lambda)\le\left\|x_0-\frac{x}{\Lambda(x)}\right\|
=\frac{\|x-\Lambda(x)x_0\|}{|\Lambda(x)|},
\]
更直接地将 $x/\Lambda(x)$ 与 $x_0$ 的同陪集比较可得
$|\Lambda(x)|\le c^{-1}\|x\|$，故有界。最后，闭的真线性子空间不稠密，所以闭核不稠密；反之，若核不闭，则其闭包仍是线性子空间。由于核为余维一，包含核的线性子空间只有核和全空间；核若非闭，其闭包只能是全空间，故稠密。三者等价。

\subsection*{7.4 乘子判别 $\ell^1$}
\begin{problem}
设数列 $(x_n)$ 满足：对任意满足 $y_n\to0$ 的数列 $(y_n)$，级数
$\sum_{n\ge1}x_ny_n$ 都收敛。证明 $\sum_{n\ge1}|x_n|<\infty$。
\end{problem}
对 $N\ge1$ 定义 $T_N:c_0\to\mathbb F$，
$T_N(y)=\sum_{n=1}^Nx_ny_n$。题设说明对每个 $y\in c_0$，数列 $T_N(y)$ 收敛，故
$\sup_N|T_N(y)|<\infty$。由一致有界原理，$\sup_N\|T_N\|<\infty$。而
\[
\|T_N\|=\sum_{n=1}^N|x_n|,
\]
因为上界由三角不等式得到，下界取有限支集序列 $y_n=\operatorname{sgn}(x_n)$（复数时取相位共轭）。故部分和有统一上界，$\sum|x_n|<\infty$。

\subsection*{7.5}
\begin{problem}
令 $X$ 为 $(0,1)$ 上所有实系数多项式的空间，赋予 $L^1(0,1)$ 范数，并定义
\[
B(f,g)=\int_0^1f(t)g(t)dt,\qquad f,g\in X.
\]
证明：$B$ 满足讲义第 46 页共轭双线性定理中的性质 c，但 $B$ 不是有界双线性泛函。此题旨在说明该定理中“$X$ 为 Banach 空间”的条件不可删去。
\end{problem}
\textbf{依赖未提供讲义，按用户要求跳过。}题目要求验证“讲义第 46 页性质 c”，但题面没有陈述性质 c，无法确定应证结论。

\section{第八次作业}

\subsection*{8.1 Fourier 部分和不必一致收敛}
\begin{problem}
令 $C(\mathbb R/2\pi\mathbb Z,\mathbb C)$ 为装备上确界范数的 $2\pi$ 周期复连续函数空间。
\begin{enumerate}[label=(\arabic*)]
\item 定义 Dirichlet 核
\[
D_n(t)=\sum_{k=-n}^ne^{ikt}
=\frac{\sin((n+\tfrac12)t)}{\sin(t/2)}.
\]
证明 $\|D_n\|_{L^1}\ge\frac8\pi\sum_{k=1}^n1/k$；
\item 对 Fourier 部分和算子
\[
(F_nf)(x)=(D_n*f)(x)
=\sum_{k=-n}^n\int_0^{2\pi}f(t)e^{ik(x-t)}dt
\]
证明 $\|F_n\|=\|D_n\|_{L^1}$（采用题目的卷积归一化）；
\item 利用一致有界性原理证明存在连续周期函数 $f$，使 $F_n(f)$ 不一致收敛。
\end{enumerate}
\end{problem}
（1）利用偶性及 $0<t\le\pi$ 时 $\sin(t/2)\le t/2$，令 $N=n+1/2$，则
\[
\|D_n\|_1
=2\int_0^\pi\frac{|\sin Nt|}{\sin(t/2)}dt
\ge4\int_0^{N\pi}\frac{|\sin u|}{u}du.
\]
对 $k=0,\ldots,n-1$，在 $[k\pi,(k+1)\pi]$ 上有 $u\le(k+1)\pi$，且该区间上 $|\sin u|$ 的积分为 $2$。逐段相加得到
\[
\|D_n\|_{L^1}\ge\frac8\pi\sum_{k=1}^n\frac1k.
\]
关键是调和和发散，故 $\|D_n\|_1\to\infty$。

（2）卷积不等式给
$\|F_nf\|_\infty\le\|D_n\|_1\|f\|_\infty$，故 $\|F_n\|\le\|D_n\|_1$。固定 $x_0$，连续逼近
$t\mapsto\overline{D_n(x_0-t)}/|D_n(x_0-t)|$（在零点附近平滑处理），可使单位上确界范数函数的卷积值任意接近 $\int|D_n|$，故反向不等式成立。这里常数依题目采用的卷积归一化；按题设约定恰为上述等式。

（3）由（1）（2），$\sup_n\|F_n\|=\infty$。一致有界性原理的逆否命题给出某个
$f\in C(\mathbb T)$，使 $\sup_n\|F_nf\|_\infty=\infty$。若 $F_nf$ 一致收敛，则其范数必有界，矛盾。因此该连续函数的 Fourier 级数不一致收敛。

\subsection*{8.2 Riemann--Lebesgue 引理与 Fourier 映射}
\begin{problem}
对 $f\in L^1([0,2\pi],\mathbb C)$ 定义
$\widehat f(k)=\int_0^{2\pi}e^{-ikt}f(t)dt$。
\begin{enumerate}[label=(\arabic*)]
\item 证明 Riemann--Lebesgue 引理：$|\widehat f(k)|\to0$（$|k|\to\infty$）；
\item 证明 Fourier 映射
$F:L^1([0,2\pi],\mathbb C)\to c_0(\mathbb Z,\mathbb C)$ 有稠密像但不是满射。提示：使用上题 Dirichlet 核的 Fourier 系数。
\end{enumerate}
\end{problem}
（1）先对阶梯函数或三角多项式直接计算，得 Fourier 系数在 $|k|\to\infty$ 时趋零。给定
$f\in L^1$ 与 $\varepsilon>0$，取阶梯函数 $g$ 使 $\|f-g\|_1<\varepsilon$，则
\[
|\widehat f(k)-\widehat g(k)|\le\|f-g\|_1<\varepsilon.
\]
令 $|k|$ 足够大使 $|\widehat g(k)|<\varepsilon$，便有 $|\widehat f(k)|<2\varepsilon$。

（2）由（1）及 $|\widehat f(k)|\le\|f\|_1$，$F:L^1\to c_0$ 有界。有限支集序列是三角多项式的 Fourier 系数，故包含于像中；它们在 $c_0$ 稠密，所以像稠密。

若 $F$ 满射，由开映射定理，存在 $C$，使每个 $a\in c_0$ 都有原像 $f$ 满足
$\|f\|_1\le C\|a\|_\infty$。取 $a^{(n)}$ 为 Dirichlet 核的 Fourier 系数，即在 $[-n,n]$ 为 $1$、其余为 $0$，则 $\|a^{(n)}\|_\infty=1$。Fourier 变换在 $L^1$ 上单射，所以其唯一原像为 $D_n$（按同一归一化），这迫使 $\|D_n\|_1\le C$，与 8.1 矛盾。因此不满射。

\subsection*{8.3 三大定理需要完备性}
\begin{problem}
令 $X=\mathbb R^\infty$ 为所有有限支集实序列组成的线性空间，并定义
$\|x\|_1=\sum_i|x_i|$、$\|x\|_\infty=\sup_i|x_i|$。证明：
\begin{enumerate}[label=(\arabic*)]
\item $\Lambda_n(x)=nx_n$ 对每个 $n$ 都有界，且对每个 $x$ 有
$\sup_n|\Lambda_n(x)|<\infty$，但 $\sup_n\|\Lambda_n\|=\infty$；这对两种范数都成立；
\item $\operatorname{id}:(X,\|\cdot\|_1)\to(X,\|\cdot\|_\infty)$ 有界，但其逆无界；
\item $\operatorname{id}:(X,\|\cdot\|_\infty)\to(X,\|\cdot\|_1)$ 图像闭，但算子无界。
\end{enumerate}
\end{problem}
（1）对两种范数，$\Lambda_n(x)=nx_n$ 都有界；固定有限支集 $x$，只有有限多个 $x_n$ 非零，故
$\sup_n|nx_n|<\infty$。在 $\ell^1$ 范数下 $\|\Lambda_n\|=n$（取 $e_n$）；在上确界范数下也为 $n$，故算子范数无界。

（2）因 $\|x\|_\infty\le\|x\|_1$，恒等映射
$(X,\|\cdot\|_1)\to(X,\|\cdot\|_\infty)$ 有界。逆映射无界，因为对
$x^{(N)}=(1,\ldots,1,0,\ldots)$，有 $\|x^{(N)}\|_\infty=1$、$\|x^{(N)}\|_1=N$。

（3）反向恒等映射无界如上。其图像闭：若 $x^{(m)}\to x$ 于上确界范数且
$x^{(m)}\to y$ 于 $\ell^1$ 范数，则后一收敛蕴含上确界收敛，所以 $x=y$，即极限点仍在恒等映射的图像中。

\subsection*{8.4}
\begin{problem}
证明讲义第 50 页注 1 的结论。
\end{problem}
\textbf{依赖未提供讲义，按用户要求跳过。}题面未写出“第 50 页注 1”的结论。

\subsection*{8.5 有下界的满射}
\begin{problem}
设 $X,Y$ 为 Banach 空间，$U\in\mathcal L(X,Y)$，且对每个 $y\in Y$ 方程 $Ux=y$ 有解。若存在 $m>0$ 使
$\|Ux\|\ge m\|x\|$ 对所有 $x\in X$ 成立，证明 $U$ 有连续逆且
$\|U^{-1}\|\le1/m$。
\end{problem}
$U$ 满射且 $\|Ux\|\ge m\|x\|$。若 $Ux=0$，则 $x=0$，故单射；于是逆映射在全体 $Y$ 上定义。对 $y=Ux$，
\[
\|U^{-1}y\|=\|x\|\le m^{-1}\|Ux\|=m^{-1}\|y\|.
\]
所以 $U^{-1}$ 连续且 $\|U^{-1}\|\le1/m$。

\section{第九次作业}

\subsection*{9.1 范数等价定理及 $L^1$ 范数下的 $C[0,1]$}
\begin{problem}
证明范数等价定理（讲义第 52 页推论 3），并利用它证明
$(C[0,1],\|\cdot\|_{L^1})$ 不是 Banach 空间。
\end{problem}
范数等价定理的常用形式为：若同一线性空间 $X$ 在范数 $\|\cdot\|_1,\|\cdot\|_2$ 下都是 Banach 空间，且存在 $C$ 使
$\|x\|_2\le C\|x\|_1$，则两范数等价。事实上恒等映射
$I:(X,\|\cdot\|_1)\to(X,\|\cdot\|_2)$ 是连续双射，由开映射定理，其逆连续，故有反向估计。

在 $C[0,1]$ 上，$\|f\|_{L^1}\le\|f\|_\infty$。若 $L^1$ 范数使它成为 Banach 空间，则上一定理会给出 $\|f\|_\infty\le C\|f\|_1$。取 $f_n(x)=x^n$，则
$\|f_n\|_\infty=1$ 而 $\|f_n\|_1=1/(n+1)\to0$，矛盾。因此不完备。

\subsection*{9.2 用闭图像定理推出逆算子和开映射定理}
\begin{problem}
利用闭图像定理证明逆算子定理，并进一步证明开映射定理。
\end{problem}
设 $T:X\to Y$ 是 Banach 空间间的有界双射。若 $y_n\to y$ 且 $T^{-1}y_n\to x$，则由 $T$ 连续，
$y_n=T(T^{-1}y_n)\to Tx$，极限唯一给 $Tx=y$，即 $T^{-1}y=x$。所以 $T^{-1}$ 图像闭；闭图像定理给出 $T^{-1}$ 有界，这就是逆算子定理。

若 $T:X\to Y$ 只是有界满射，商空间 $X/\ker T$ 是 Banach 空间，诱导映射
$\widetilde T:X/\ker T\to Y$ 是有界双射。由逆算子定理，$\widetilde T^{-1}$ 连续，故 $\widetilde T$ 为开映射。商映射 $X\to X/\ker T$ 也开，因此 $T$ 开，这就是开映射定理。

\subsection*{9.4}
\begin{problem}
完善判别可补充空间的充分条件，即讲义第 57 页中的引理。
\end{problem}
\textbf{依赖未提供讲义，按用户要求跳过。}题面没有写出“第 57 页中的引理”的陈述。

\subsection*{9.5 闭算子的基本性质}
\begin{problem}
设 $X,Y$ 为赋范空间，$D\subset X$ 为线性子空间，$A:D\to Y$ 为线性映射。证明：
\begin{enumerate}[label=(\arabic*)]
\item 若 $A$ 连续且 $D$ 闭，则 $A$ 是闭算子；
\item 若 $A$ 连续且为闭算子，则 $Y$ 完备蕴含 $D$ 闭；
\item 若 $A$ 是单射闭算子，则 $A^{-1}$ 也是闭算子；
\item 若 $X$ 完备，$A$ 是单射闭算子，$R(A)$ 在 $Y$ 中稠密且 $A^{-1}$ 连续，则 $R(A)=Y$。
\end{enumerate}
\end{problem}
（1）设 $x_n\in D$、$x_n\to x$ 于 $X$、$Ax_n\to y$ 于 $Y$。$D$ 闭给 $x\in D$，而连续性给 $Ax_n\to Ax$，故 $y=Ax$，图像闭。

（2）假设 $A$ 连续且图像闭，$Y$ 完备。若 $x_n\in D$ 且 $x_n\to x$，则 $(x_n)$ Cauchy；连续性使 $(Ax_n)$ 在 $Y$ 中 Cauchy，故收敛到某 $y$。闭图像性给 $x\in D$ 且 $Ax=y$，所以 $D$ 闭。

（3）若 $A$ 单射，则
\[
\operatorname{Graph}(A^{-1})=\{(Ax,x):x\in D\}
\]
是 $\operatorname{Graph}(A)$ 交换两个坐标后的像。坐标交换是 $X\times Y$ 与 $Y\times X$ 间的同胚，故前者闭。

（4）$A^{-1}:R(A)\to X$ 连续。若 $y_n=Ax_n\in R(A)$ 且 $y_n\to y$，则
$x_n=A^{-1}y_n$ 为 Cauchy；$X$ 完备，故 $x_n\to x$。闭图像性给 $x\in D$ 且 $Ax=y$，所以 $y\in R(A)$。故 $R(A)$ 闭；它又稠密，只能等于 $Y$。

\subsection*{9.6 弱导数与 $H^1$}
\begin{problem}
令 $T=\partial/\partial x$，初始定义域为 $C^1[a,b]\subset L^2[a,b]$。
\begin{enumerate}[label=(\arabic*)]
\item 证明 $T$ 可闭；提示：使用
$f(x)-f(a)=\int_a^xf'(t)dt$；
\item 若
\[
\operatorname{Graph}(\overline T)=
\{(f,g):\exists f_n\in C^1[a,b], f_n\to f, f_n'\to g\text{ 于 }L^2\},
\]
证明它等于 $\{(f,f'):f\in H^1[a,b]\}$，其中
$H^1[a,b]=\{f\in AC[a,b]:f'\in L^2[a,b]\}$；
\item 在范数 $\|f\|_{H^1}^2=\|f\|_{L^2}^2+\|f'\|_{L^2}^2$ 下，若
$H_0^1[a,b]$ 是 $C_c^\infty(a,b)$ 的闭包，说明 $H_0^1[a,b]$ 与 $H^1[a,b]$ 的差别。
\end{enumerate}
\end{problem}
（1）把 $T$ 看成从 $L^2[a,b]$ 到 $L^2[a,b]$、定义域为 $C^1[a,b]$ 的算子。若
$f_n\to0$、$f_n'\to g$ 于 $L^2$，则对任意 $\varphi\in C_c^\infty(a,b)$，
\[
\int g\varphi=\lim_n\int f_n'\varphi=-\lim_n\int f_n\varphi'=0.
\]
故 $g=0$，所以 $T$ 可闭。

（2）若 $f_n\to f$、$f_n'\to g$ 于 $L^2$，同样分部积分得
$\int f\varphi'=-\int g\varphi$，故 $g$ 是 $f$ 的弱导数。区间上一维 Sobolev 定理说明 $f$ 有绝对连续代表且经典导数为 $g\in L^2$，所以 $f\in H^1$。反之，对 $H^1$ 函数作延拓后用光滑核卷积，得到 $C^1$（甚至光滑）函数同时在函数和导数的 $L^2$ 范数下逼近它。因此闭包图像恰为
$\{(f,f'):f\in H^1[a,b]\}$。

（3）迹映射 $\operatorname{Tr}:H^1[a,b]\to\mathbb F^2$，$f\mapsto(f(a),f(b))$ 连续。每个
$C_c^\infty(a,b)$ 的两个迹都为零，故
$H_0^1\subset\{f\in H^1:f(a)=f(b)=0\}$。反之，零边界值的 $H^1$ 函数先在端点外作零延拓，再卷积并截断，可由 $C_c^\infty(a,b)$ 在 $H^1$ 中逼近。因此
\[
H_0^1[a,b]=\{f\in H^1[a,b]:f(a)=f(b)=0\}.
\]
所以 $H_0^1$ 比 $H^1$ 多了两个齐次边界条件，是后者的闭子空间；例如常数函数 $1$ 属于 $H^1$ 但不属于 $H_0^1$。


\chapter{第十次至第十八次作业}
\section{第十次作业}

\subsection*{10.1　带状区域上的 Poincar\'e 不等式与迹算子}

\begin{problem}
考虑开集
\[
 \Omega=\{(x_1,x_2)\in\mathbb R^2:|x_1|<1, x_2\in\mathbb R\},
 \qquad \Gamma=\{(-1,x_2):x_2\in\mathbb R\}.
\]
\begin{enumerate}
 \item 证明存在常数 $C>0$，使
 \[
  \|f\|_{L^2(\Omega)}\le C\|\nabla f\|_{L^2(\Omega)},
  \qquad
  f\in\{g\in H^1(\Omega)\cap C^1(\overline\Omega):g|_\Gamma=0\}.
 \]
 \item 证明存在常数 $C>0$，使对任意 $f\in H^1(\Omega)\cap C^1(\overline\Omega)$，
 \[
  \|f\|_{L^2(\partial\Omega)}
  :=\|f(-1,\cdot)\|_{L^2(\mathbb R)}+\|f(1,\cdot)\|_{L^2(\mathbb R)}
  \le C\|f\|_{H^1(\Omega)}.
 \]
 \item 假设 $H^1(\Omega)\cap C^1(\overline\Omega)$ 在 $H^1(\Omega)$ 中稠密，证明存在
 唯一线性映射
 \[
  \operatorname{Tr}:H^1(\Omega)\to L^2(\partial\Omega)
 \]
 使 $\operatorname{Tr}f=f|_{\partial\Omega}$ 对
 $f\in H^1(\Omega)\cap C^1(\overline\Omega)$ 成立，并且
 $\|\operatorname{Tr}f\|_{L^2(\partial\Omega)}\le C\|f\|_{H^1(\Omega)}$。
 \item Poincar\'e 不等式能否延伸到
 $f\in\{g\in H^1(\Omega):\operatorname{Tr}(g)|_\Gamma=0\}$？
\end{enumerate}
\end{problem}

记 $\Omega=(-1,1)\times\mathbb R$，$\Gamma=\{-1\}\times\mathbb R$。

\textbf{(1)} 先设 $f\in H^1(\Omega)\cap C^1(\overline\Omega)$ 且 $f|_\Gamma=0$。固定 $x_2$，由微积分基本定理，
\[
 f(x_1,x_2)=\int_{-1}^{x_1}\partial_1f(t,x_2)\,dt .
\]
故由 Cauchy--Schwarz 不等式，
\[
 |f(x_1,x_2)|^2
 \le (x_1+1)\int_{-1}^{x_1}|\partial_1f(t,x_2)|^2dt
 \le 2\int_{-1}^{1}|\partial_1f(t,x_2)|^2dt .
\]
再依次对 $x_1\in(-1,1)$、$x_2\in\mathbb R$ 积分，得到
\[
 \|f\|_{L^2(\Omega)}^2\le4\|\partial_1f\|_{L^2(\Omega)}^2
 \le4\|\nabla f\|_{L^2(\Omega)}^2.
\]
因而可取 $C=2$。

\textbf{(2)} 对固定的 $x_2$ 和任意 $x_1\in[-1,1]$，
\[
 |f(-1,x_2)|\le |f(x_1,x_2)|+
 \int_{-1}^{1}|\partial_1f(t,x_2)|dt.
\]
利用 $(a+b)^2\le2a^2+2b^2$ 以及 Cauchy--Schwarz 不等式，
\[
 |f(-1,x_2)|^2
 \le2|f(x_1,x_2)|^2+4\int_{-1}^{1}|\partial_1f(t,x_2)|^2dt.
\]
对 $x_1$ 取平均，再对 $x_2$ 积分，得
\[
 \|f(-1,\cdot)\|_2^2\le \|f\|_2^2+4\|\partial_1f\|_2^2
 \le4\|f\|_{H^1(\Omega)}^2.
\]
对 $x_1=1$ 同理。因此
\[
 \|f\|_{L^2(\partial\Omega)}
 =\|f(-1,\cdot)\|_2+\|f(1,\cdot)\|_2
 \le4\|f\|_{H^1(\Omega)}.
\]

\textbf{(3)} 令 $D=H^1(\Omega)\cap C^1(\overline\Omega)$。由题设，$D$ 在 $H^1(\Omega)$ 中稠密。
上一步说明限制映射
\[
 R:D\longrightarrow L^2(\partial\Omega),\qquad Rf=f|_{\partial\Omega}
\]
有界。给定 $f\in H^1(\Omega)$，取 $f_n\in D$ 使 $f_n\to f$ 于 $H^1$。则
\[
 \|Rf_n-Rf_m\|_{L^2(\partial\Omega)}\le C\|f_n-f_m\|_{H^1}\to0.
\]
由于 $L^2(\partial\Omega)$ 完备，$Rf_n$ 有极限。定义
\[
 \operatorname{Tr}f:=\lim_{n\to\infty}Rf_n.
\]
若换用另一逼近列，以上估计表明所得极限相同。因此 $\operatorname{Tr}$ 定义良好、线性且
\[
 \|\operatorname{Tr}f\|_{L^2(\partial\Omega)}\le C\|f\|_{H^1(\Omega)}.
\]
唯一性来自 $D$ 的稠密性。

\textbf{(4)} 可以延伸。Sobolev 函数沿几乎每一条平行于 $x_1$ 轴的直线都有绝对连续代表，并且其经典端点值等于迹。若 $\operatorname{Tr}f|_\Gamma=0$，则对几乎处处的 $x_2$，
\[
 f(x_1,x_2)=\int_{-1}^{x_1}\partial_1f(t,x_2)dt.
\]
重复 (1) 的估计即得
\[
 \|f\|_{L^2(\Omega)}\le2\|\nabla f\|_{L^2(\Omega)}.
\]

\subsection*{10.2　一个混合边值问题}

\begin{problem}
考虑方程
\[
 -u''+\alpha u=f\quad (0<x<1),\qquad
 u(0)=0,quad u'(1)=b,quad f\in L^2(0,1),\quad \alpha\ge0.
\]
其弱形式为：求 $u\in H:=\{g\in H^1(0,1):g(0)=0\}$，使对任意 $v\in H$，
\[
 \int_0^1u'v'\,dx+\alpha\int_0^1uv\,dx
 =\int_0^1fv\,dx+bv(1).
\]
利用 Lax--Milgram 定理说明该弱问题存在唯一解。
\end{problem}

在
\[
 H=\{v\in H^1(0,1):v(0)=0\}
\]
上定义
\[
 a(u,v)=\int_0^1u'v'\,dx+\alpha\int_0^1uv\,dx,
 \qquad L(v)=\int_0^1fv\,dx+bv(1).
\]
显然 $a$ 是连续双线性型。又因 $v(0)=0$，
\[
 |v(x)|=\left|\int_0^xv'(t)dt\right|\le\|v'\|_2,
 \qquad \|v\|_2\le\|v'\|_2.
\]
从而
\[
 a(v,v)=\|v'\|_2^2+\alpha\|v\|_2^2
 \ge\|v'\|_2^2\ge\frac12\|v\|_{H^1}^2.
\]
所以 $a$ 强制。另一方面，
\[
 |L(v)|\le\|f\|_2\|v\|_2+|b|\,|v(1)|
 \le(\|f\|_2+|b|)\|v'\|_2
 \le(\|f\|_2+|b|)\|v\|_{H^1}.
\]
故 $L\in H^*$。由 Lax--Milgram 定理，存在唯一 $u\in H$ 使
$a(u,v)=L(v)$ 对所有 $v\in H$ 成立，这正是所给弱问题的唯一解。

\subsection*{10.3　由数值下界推出可逆性}

\begin{problem}
设 $H$ 是 Hilbert 空间，$A\in\mathcal L(H)$，并且存在 $m>0$ 使
\[
 |\langle Ax,x\rangle|\ge m\|x\|^2,\qquad x\in H.
\]
证明 $A$ 存在有界逆算子 $A^{-1}\in\mathcal L(H)$。
\end{problem}

记数值域 $W(A)=\{\langle Ax,x\rangle:\|x\|=1\}$。题设说明
$\operatorname{dist}(0,W(A))\ge m$。数值域的闭包是凸集，因此存在模为 $1$ 的标量
$\omega$，使得（实情形只需取 $\omega=1$ 或 $-1$）
\[
 \operatorname{Re}\langle \omega Ax,x\rangle\ge m\|x\|^2,
 \qquad x\in H.
\]
这里可由凸集分离定理得到：用过原点且垂直于离原点最近方向的直线分离
$0$ 与 $\overline{W(A)}$。

令 $a(x,y)=\langle\omega Ax,y\rangle$。它是有界的强制 sesquilinear 型。对任意
$z\in H$，泛函 $y\mapsto\langle z,y\rangle$ 有界，故 Lax--Milgram 定理给出唯一
$x\in H$ 使
\[
 \langle\omega Ax,y\rangle=\langle z,y\rangle,\qquad y\in H.
\]
于是 $\omega Ax=z$，说明 $A$ 满射；强制性同时给出单射。并且
\[
 m\|x\|^2\le|\langle Ax,x\rangle|\le\|Ax\|\|x\|,
\]
所以 $\|x\|\le m^{-1}\|Ax\|$。令 $x=A^{-1}y$，得
$\|A^{-1}y\|\le m^{-1}\|y\|$，故 $A^{-1}\in\mathcal L(H)$。

\subsection*{10.4　次线性泛函}

\begin{problem}
设 $p$ 是实线性空间 $X$ 上的次线性泛函。证明：
\begin{enumerate}
 \item $p(0)=0$；
 \item $p(-x)\ge-p(x)$；
 \item 给定 $x_0\in X$，存在实线性泛函 $f:X\to\mathbb R$，满足
 $f(x_0)=p(x_0)$ 且 $f(x)\le p(x)$ 对所有 $x\in X$ 成立。
\end{enumerate}
\end{problem}

由次可加性，$p(0)=p(0+0)\le2p(0)$，而正齐次性给出
$p(0)=p(0x)=0p(x)=0$，故 (1) 成立。再由
\[
 0=p(0)=p(x-x)\le p(x)+p(-x)
\]
得到 $p(-x)\ge-p(x)$。

在一维子空间 $M=\mathbb Rx_0$ 上定义
\[
 f_0(tx_0)=t\,p(x_0).
\]
若 $t\ge0$，则 $f_0(tx_0)=p(tx_0)$；若 $t<0$，由上一结论
\[
 f_0(tx_0)=(-t)(-p(x_0))\le(-t)p(-x_0)=p(tx_0).
\]
因此 $f_0\le p$ 于 $M$。Hahn--Banach 定理给出线性延拓 $f:X\to\mathbb R$，满足
$f\le p$ 于 $X$，并且 $f(x_0)=p(x_0)$。

\subsection*{10.5　上极限泛函}

\begin{problem}
设 $X$ 是由实数列 $x=(a_n)$ 组成的实线性空间，坐标逐项定义相等与线性运算。令
\[
 p(x)=\limsup_{n\to\infty}a_n.
\]
证明 $p$ 是 $X$ 上的次线性泛函。
\end{problem}

题面若把 $X$ 解释为“所有实数列”，则命题按字面并不成立：例如 $a_n=n$ 时
$\limsup a_n=+\infty$，$p$ 不是实值泛函。以下按通常且使题目成立的表述，将
$X$ 理解为所有上极限有限的序列（特别地可取 $X=\ell^\infty$）。

对 $x=(a_n),y=(b_n)$，有
\[
 \limsup(a_n+b_n)\le\limsup a_n+\limsup b_n,
\]
因为对每个 $N$，
\[
 \sup_{n\ge N}(a_n+b_n)
 \le\sup_{n\ge N}a_n+\sup_{n\ge N}b_n,
\]
再令 $N\to\infty$ 即可。若 $\lambda\ge0$，则
\[
 p(\lambda x)=\limsup(\lambda a_n)=\lambda\limsup a_n=\lambda p(x).
\]
所以 $p$ 满足次可加性和非负齐次性，是次线性泛函。

\section{第十一次作业}

\subsection*{11.1　有限个线性无关向量的对偶坐标}

\begin{problem}
给定赋范线性空间 $X$ 中 $n$ 个线性无关元素 $x_1,\dots,x_n$，证明存在
$f_1,\dots,f_n\in X^*$，使
\[
 \langle f_i,x_j\rangle=\delta_{ij},\qquad 1\le i,j\le n.
\]
\end{problem}

令 $E=\operatorname{span}\{x_1,\dots,x_n\}$。在线性无关基
$x_1,\dots,x_n$ 下定义坐标泛函
\[
 g_i\!\left(\sum_{j=1}^na_jx_j\right)=a_i.
\]
有限维空间上的线性泛函自动连续，故 $g_i\in E^*$。由 Hahn--Banach 定理，存在
$f_i\in X^*$ 延拓 $g_i$ 且 $\|f_i\|=\|g_i\|$。于是
\[
 \langle f_i,x_j\rangle=g_i(x_j)=\delta_{ij}.
\]

\subsection*{11.2　Banach 极限}

\begin{problem}
令 $\ell^\infty$ 为实有界序列组成的 Banach 空间，定义移位算子
$T:\ell^\infty\to\ell^\infty$ 为
\[
 Tx=(x_{n+1})_{n\ge1},\qquad x=(x_n)_{n\ge1},
\]
并令
\[
 Y=\operatorname{Im}(I-T)=\{x-Tx:x\in\ell^\infty\}.
\]
\begin{enumerate}
 \item 记 $c_0\subset\ell^\infty$ 为收敛到 $0$ 的序列空间，证明 $c_0\subset Y$。
 \item 记 $\mathbf1=(1,1,\dots)$。证明
 \[
  \sup_n|1+x_{n+1}-x_n|\ge1,\qquad x\in\ell^\infty,
 \]
 并推出
 \[
  d(\mathbf1,Y)=\inf_{y\in Y}\|\mathbf1-y\|_\infty=1.
 \]
 \item 由 Hahn--Banach 定理可取 $\Lambda\in(\ell^\infty)^*$，满足
 \[
  \Lambda(\mathbf1)=1,\qquad \|\Lambda\|=1,\qquad
  \Lambda(x-Tx)=0\quad(x\in\ell^\infty).
 \]
 证明任意这样的 $\Lambda$ 满足：
 \begin{enumerate}
  \item $\Lambda(Tx)=\Lambda(x)$；
  \item 若 $x_n\ge0$，则 $\Lambda(x)\ge0$；
  \item $\liminf_{n\to\infty}x_n\le\Lambda(x)\le\limsup_{n\to\infty}x_n$；
  \item 若 $x_n$ 收敛，则 $\Lambda(x)=\lim_{n\to\infty}x_n$。
 \end{enumerate}
 \item 找出 $x,y\in\ell^\infty$，使
 $\Lambda(xy)\ne\Lambda(x)\Lambda(y)$。提示：考虑 $x_n=(-1)^n$ 并证明
 $\Lambda(x)=0$。
 \item 证明不存在 $y=(y_n)\in\ell^1$ 使
 \[
  \Lambda(x)=\sum_{n=1}^\infty x_ny_n,\qquad x\in\ell^\infty.
 \]
 提示：由正性可推出 $y_n\ge0$ 且 $\sum y_n=1$，再与收敛序列上的性质矛盾。
\end{enumerate}
\end{problem}

\textbf{(1)} 题面写成 $c_0\subset Y$ 时结论不成立。取 $a_n=1/n\in c_0$；若
$a=x-Tx$，则
\[
 x_{N+1}=x_1-\sum_{n=1}^N\frac1n,
\]
与 $x\in\ell^\infty$ 矛盾。正确且在后续论证中足够的结论是
\[
 c_{00}\subset Y,\qquad c_0\subset\overline Y.
\]
事实上，若 $a_n=0$ 对 $n>N$，取
$x_n=\sum_{k=n}^Na_k$（$n\le N$）且 $x_n=0$（$n>N$），则
$x-Tx=a$。再利用 $c_{00}$ 在 $c_0$ 中稠密即可。

\textbf{(2)} 对任意 $x\in\ell^\infty$，若
$\sup_n|1+x_{n+1}-x_n|<1$，则存在 $\varepsilon>0$ 使
$1+x_{n+1}-x_n\ge\varepsilon$，即
$x_{n+1}-x_n\ge-1+\varepsilon$；这本身还不足以矛盾。更直接地，对前 $N$ 项取平均：
\[
 \sup_n|1+x_{n+1}-x_n|
 \ge\left|\frac1N\sum_{n=1}^N(1+x_{n+1}-x_n)\right|
 =\left|1+\frac{x_{N+1}-x_1}{N}\right|.
\]
令 $N\to\infty$，右端趋于 $1$，故上确界至少为 $1$。于是对
$y=x-Tx\in Y$，$\|\mathbf1-y\|_\infty\ge1$。取 $y=0$ 又给出反向不等式，故
$d(\mathbf1,Y)=1$。

\textbf{(3)} 在 $\overline Y+\mathbb R\mathbf1$ 上定义
$L(y+c\mathbf1)=c$。由上一问，$L$ 定义良好且范数为 $1$。Hahn--Banach 延拓给出
$\Lambda\in(\ell^\infty)^*$，满足 $\Lambda(\mathbf1)=1$、$\|\Lambda\|=1$ 且
$\Lambda|_Y=0$。

(a) 因 $x-Tx\in Y$，
\[
 \Lambda(x)-\Lambda(Tx)=\Lambda(x-Tx)=0.
\]

(b) 若 $x\ge0$，令 $M=\|x\|_\infty$。$M=0$ 时显然。$M>0$ 时
$\|\mathbf1-x/M\|_\infty\le1$，故
\[
 \Lambda(x)=M\{1-\Lambda(\mathbf1-x/M)\}
 \ge M(1-|\Lambda(\mathbf1-x/M)|)\ge0.
\]

(c) 记 $\alpha_N=\inf_{n\ge N}x_n$、$\beta_N=\sup_{n\ge N}x_n$。
由 (a)，$\Lambda(x)=\Lambda(T^{N-1}x)$；由正性，
\[
 \alpha_N=\Lambda(\alpha_N\mathbf1)\le\Lambda(T^{N-1}x)\le
 \Lambda(\beta_N\mathbf1)=\beta_N.
\]
令 $N\to\infty$，即得
\[
 \liminf x_n\le\Lambda(x)\le\limsup x_n.
\]

(d) 若 $x_n\to a$，上下极限均为 $a$，由 (c) 得 $\Lambda(x)=a$。

\textbf{(4)} 取 $x_n=(-1)^n$。因为 $Tx=-x$，移位不变性给出
$\Lambda(x)=\Lambda(Tx)=-\Lambda(x)$，所以 $\Lambda(x)=0$。令 $y=x$，则
$xy=\mathbf1$，从而
\[
 \Lambda(xy)=1\ne0=\Lambda(x)\Lambda(y).
\]

\textbf{(5)} 假设存在 $y\in\ell^1$ 使
$\Lambda(x)=\sum x_ny_n$。取标准基 $e_n\in c_0\subset\overline Y$，有
$y_n=\Lambda(e_n)=0$。这对每个 $n$ 都成立，于是 $y=0$，但这将导致
$1=\Lambda(\mathbf1)=0$，矛盾。因此不存在这样的 $y$。

\subsection*{11.3　Hahn--Banach 延拓的不唯一性}

\begin{problem}
定义
\[
 C=\{x=(x_n)\in\ell^\infty:\lim_{n\to\infty}x_n\text{ 存在}\},
\]
\[
 S=\{x=(x_n)\in\ell^\infty:\lim_{n\to\infty}x_{2n}
 \text{ 与 }\lim_{n\to\infty}x_{2n+1}\text{ 均存在}\}.
\]
令
\[
 f(x)=\lim_nx_n\quad(x\in C),\qquad
 f_1(x)=\lim_nx_{2n},\quad f_2(x)=\lim_nx_{2n+1}\quad(x\in S).
\]
证明：
\begin{enumerate}
 \item $S,C$ 是 $\ell^\infty$ 的闭子空间；
 \item $f\in C^*$、$f_1,f_2\in S^*$，且
 $\|f\|=\|f_1\|=\|f_2\|=1$，并有 $f_1|_C=f_2|_C$；
 \item 构造例子说明 Hahn--Banach 延拓有无穷多种方式。
\end{enumerate}
\end{problem}

$C$ 和 $S$ 都是线性子空间。若 $x^{(k)}\in C$ 且
$x^{(k)}\to x$ 于 $\ell^\infty$，给定 $\varepsilon>0$，先取 $k$ 使
$\|x^{(k)}-x\|_\infty<\varepsilon/3$，再利用 $x^{(k)}$ 收敛验证 $x$ 是 Cauchy
序列，故 $x\in C$。对偶数、奇数子列分别作同样论证，得到 $S$ 也闭。

显然
\[
 |f(x)|,|f_1(x)|,|f_2(x)|\le\|x\|_\infty,
\]
故三者范数不超过 $1$；代入 $\mathbf1$ 得三者范数恰为 $1$。若 $x\in C$，偶、奇
子列极限都等于 $\lim x_n$，所以 $f_1|_C=f_2|_C=f$。

对每个 $t\in[0,1]$ 定义
\[
 F_t=t f_1+(1-t)f_2\in S^*.
\]
则 $F_t|_C=f$，且由正系数组合 $\|F_t\|\le1$，代入 $\mathbf1$ 又知
$\|F_t\|=1$。取交错序列 $u=(0,1,0,1,\dots)$，不同的 $t$ 给出不同的
$F_t(u)$。故至少在子空间 $S$ 上已有不可数多个保范延拓；再分别用 Hahn--Banach
延拓到 $\ell^\infty$，便得到 $f$ 的无穷多种延拓。

\section{第十二次作业}

\subsection*{12.1　吸收集、代数内点与 Minkowski 泛函}

\begin{problem}
设 $K$ 是实线性空间 $X$ 的子集，$P_K$ 表示 $K$ 的 Minkowski 泛函。证明：
\begin{enumerate}
 \item $K$ 是吸收集，当且仅当 $0$ 是 $K$ 的代数内点；
 \item 若 $K$ 是凸集，则 $K$ 是吸收集，当且仅当
 $P_K(x)<\infty$ 对所有 $x\in X$ 成立；
 \item 若 $K$ 是平衡凸吸收集，令
 \[
  V=\{x\in X:P_K(x)<1\},\qquad W=\{x\in X:P_K(x)\le1\}.
 \]
 则 $V\subset K\subset W$；$V,W$ 也是平衡凸吸收集；$V$ 是 $K$ 的全部代数
 内点所成的集合；并且 $P_K=P_V=P_W$。
\end{enumerate}
\end{problem}

称 $0$ 是 $K$ 的代数内点，是指对每个 $x\in X$，存在 $\varepsilon>0$，使
$tx\in K$ 对所有 $|t|<\varepsilon$ 成立。

\textbf{(1)} 按吸收集的定义，对每个 $x$ 都存在 $\varepsilon_x>0$，使
$tx\in K$ 对所有 $|t|<\varepsilon_x$ 成立；这恰好说明 $0$ 是 $K$ 的代数内点。
反之，若 $0$ 是代数
内点，给定 $x$ 可取 $t>0$ 使 $tx\in K$，于是 $x\in t^{-1}K$，故 $K$ 吸收。

\textbf{(2)} 设 $K$ 凸且含 $0$。若 $K$ 吸收，则对每个 $x$ 存在 $\lambda>0$
使 $x\in\lambda K$，故
\[
 P_K(x)=\inf\{t>0:x\in tK\}<\infty.
\]
反之，若 $P_K(x)<\infty$，则定义上存在某个 $t<\infty$ 使 $x\in tK$，所以
$K$ 吸收。

\textbf{(3)} 设 $K$ 平衡、凸且吸收。若 $P_K(x)<1$，取
$P_K(x)<t<1$ 且 $x\in tK$；由平衡性 $tK\subset K$，所以 $x\in K$。
若 $x\in K$，则 $P_K(x)\le1$。因此
\[
 V:=\{P_K<1\}\subset K\subset W:=\{P_K\le1\}.
\]
$P_K$ 是半范数：正齐次性直接来自定义；若 $x\in rK,y\in sK$，则
\[
 \frac{x+y}{r+s}=\frac r{r+s}\frac xr+\frac s{r+s}\frac ys\in K,
\]
故 $P_K(x+y)\le r+s$，令 $r\downarrow P_K(x),s\downarrow P_K(y)$ 得次可加性。
由此立即验证 $V,W$ 平衡、凸且吸收。

若 $x\in V$，取 $P_K(x)<a<1$。对任意方向 $y$，当
$|t|<(1-a)/P_K(y)$ 时，
\[
 P_K(x+ty)\le P_K(x)+|t|P_K(y)<1,
\]
所以 $x$ 是 $K$ 的代数内点。反之，若 $x$ 是代数内点，沿方向 $x$ 可取
$\varepsilon>0$ 使 $(1+\varepsilon)x\in K$，故
$P_K(x)\le(1+\varepsilon)^{-1}<1$。所以 $V$ 正好是代数内点集合。

由 $V\subset K\subset W$ 得 $P_W\le P_K\le P_V$。而
$x\in tV$ 当且仅当 $P_K(x)<t$，取下确界可得 $P_V=P_K$；同理
$P_W=P_K$。

\subsection*{12.2　Minkowski 泛函的连续性}

\begin{problem}
设 $X$ 为实赋范线性空间，$0\in C\subset X$ 为凸集，定义 Minkowski 泛函
\[
 p(x)=\inf\{\lambda>0:\lambda^{-1}x\in C\},\qquad x\in X.
\]
证明 $p$ 连续当且仅当 $0\in\operatorname{int}C$；此时
\[
 \operatorname{int}C=p^{-1}([0,1)),\qquad
 \overline C=p^{-1}([0,1]).
\]
\end{problem}

题面最后应理解为
\[
 \operatorname{int}C=p^{-1}([0,1)),\qquad
 \overline C=p^{-1}([0,1]);
\]
若没有假设 $C$ 闭，不能把第二个集合写成 $C$（开单位球就是反例）。

先设 $0\in\operatorname{int}C$。存在 $r>0$ 使 $rB_X\subset C$，因此
$p(x)\le r^{-1}\|x\|$。由次线性，
\[
 |p(x)-p(y)|\le\max\{p(x-y),p(y-x)\}\le r^{-1}\|x-y\|,
\]
所以 $p$ 连续。反之，若 $p$ 在 $0$ 连续，则存在 $r>0$ 使 $\|x\|<r$
时 $p(x)<1$。由 $\{p<1\}\subset C$，有 $rB_X\subset C$，故
$0\in\operatorname{int}C$。

连续时，$\{p<1\}$ 是开集且包含于 $C$；反向包含由沿射线稍微放大内点得到，故
$\operatorname{int}C=\{p<1\}$。又 $C\subset\{p\le1\}$ 给出
$\overline C\subset\{p\le1\}$。若 $p(x)\le1$，则 $0<t<1$ 时
$p(tx)<1$，故 $tx\in C$，而 $tx\to x$，于是 $x\in\overline C$。所以
$\overline C=\{p\le1\}$。

\subsection*{12.3　平衡集上的线性泛函估计}

\begin{problem}
设 $X$ 是赋范线性空间，$E\subset X$ 是非空平衡凸集，$f$ 是 $X$ 上的线性泛函。
证明
\[
 |f(x)|\le\sup_{y\in E}\operatorname{Re}f(y),\qquad x\in E.
\]
\end{problem}

实数情形，由平衡性 $-x\in E$，所以
\[
 |f(x)|=\max\{f(x),f(-x)\}\le\sup_{y\in E}f(y).
\]
复数情形，若 $f(x)\ne0$，取 $|\lambda|=1$ 使
$\lambda f(x)=|f(x)|$。因 $\lambda x\in E$，故
\[
 |f(x)|=\operatorname{Re}f(\lambda x)
 \le\sup_{y\in E}\operatorname{Re}f(y).
\]
$f(x)=0$ 时结论显然。

\subsection*{12.4　点与平衡闭凸集的严格分离}

\begin{problem}
设 $X$ 是实赋范线性空间，$E\subset X$ 是非空平衡闭凸集。证明对每个
$x_0\in X\setminus E$，存在 $f\in X^*$ 与 $a>0$，使
\[
 |f(x)|<a<|f(x_0)|,\qquad x\in E.
\]
\end{problem}

因 $E$ 闭且 $x_0\notin E$，有 $d=\operatorname{dist}(x_0,E)>0$。开凸集
$G=E+B(0,d/2)$ 与点 $x_0$ 不交。严格分离定理给出非零 $f\in X^*$，使
\[
 \sup_{z\in G}f(z)<f(x_0).
\]
由于 $E$ 平衡，$\sup_E|f|=\sup_Ef$。把 $f$ 换成 $-f$（如有必要），可令
$f(x_0)>0$。取
\[
 \sup_{x\in E}|f(x)|<a<|f(x_0)|,
\]
便得到结论。

\subsection*{12.5　开平衡凸集与不交凸集的分离}

\begin{problem}
设 $E,F$ 是实赋范线性空间 $X$ 中两个互不相交的非空凸集，并且 $E$ 是开集且平衡。
证明存在 $f\in X^*$，使
\[
 |f(x)|<\inf_{y\in F}|f(y)|,\qquad x\in E.
\]
\end{problem}

由一个集合开时的凸集分离定理，存在非零 $f\in X^*$ 和 $c\in\mathbb R$，使
\[
 f(x)<c\le f(y),\qquad x\in E, y\in F.
\]
因为 $E$ 平衡，$0\in E$ 且 $x\in E\Rightarrow -x\in E$，所以
$|f(x)|<c$，特别地 $c>0$。对 $y\in F$ 有 $|f(y)|\ge f(y)\ge c$。于是
\[
 |f(x)|<c\le\inf_{y\in F}|f(y)|,\qquad x\in E.
\]

\subsection*{12.6　不能严格分离的两个闭凸集}

\begin{problem}
在 $\ell^2$ 中定义
\[
 A=\{x=(x_i):x_i=0\text{ 对所有 }i>1\},
\]
以及
\[
 B=\{x=(x_i)_{i=1}^\infty\in\mathbb R^{\mathbb N}:
 |ix_i-i^{1/3}|\le x_1\text{ 对所有 }i>1\}\cap\ell^2.
\]
证明 $A,B$ 是 $\ell^2$ 中互不相交的非空闭凸集，且 $A-B$ 在 $\ell^2$ 中稠密，
并由此推出 $A,B$ 不能被仿射超平面分离。
\end{problem}

写
\[
 A=\{(a,0,0,\dots):a\in\mathbb R\},
\quad B=\{x\in\ell^2:|ix_i-i^{1/3}|\le x_1, i>1\}.
\]
$A$ 显然非空、闭、凸。$B$ 非空，因为
$b=(0,2^{-2/3},3^{-2/3},\dots)\in\ell^2$ 且等号成立。每个条件等价于两个闭
半空间条件，故 $B$ 是闭凸集。若 $x\in A\cap B$，则 $x_i=0$（$i>1$），从而
$i^{1/3}\le x_1$ 对所有 $i>1$，不可能，故两集不交。

给定 $z\in\ell^2$ 和 $\varepsilon>0$，取 $N$ 足够大，使
\[
 \sum_{i>N}|z_i|^2<\varepsilon^2/4,
 \qquad \sum_{i>N}i^{-4/3}<\varepsilon^2/4.
\]
定义 $b_i=-z_i$（$2\le i\le N$），$b_i=i^{-2/3}$（$i>N$），并取
\[
 b_1\ge\max_{2\le i\le N}|ib_i-i^{1/3}|.
\]
则 $b\in B$。令 $a=(z_1+b_1,0,0,\dots)\in A$。此时
\[
 \|(a-b)-z\|_2^2
 \le2\sum_{i>N}i^{-4/3}+2\sum_{i>N}|z_i|^2<\varepsilon^2.
\]
故 $A-B$ 稠密。若仿射超平面分离 $A,B$，适当改变方向后存在非零
$f\in(\ell^2)^*$ 和 $c$，使 $f(a)\le c\le f(b)$ 对所有 $a\in A,b\in B$
成立。因此 $f(a-b)\le0$ 于稠密集 $A-B$ 上成立。由连续性，$f(x)\le0$ 对所有
$x\in\ell^2$ 成立；再以 $-x$ 代入便得 $f=0$，矛盾。因此两集不能被仿射超平面
分离。

\section{第十三次作业}

\subsection*{13.1　凸集边界点沿射线外推}

\begin{problem}
设 $C$ 是实赋范线性空间 $X$ 中的凸集，$x_0\in\operatorname{int}C$、
$x_1\in\partial C$，并令
\[
 x_2=m(x_1-x_0)+x_0,\qquad m>1.
\]
证明 $x_2\notin C$。
\end{problem}

若 $u\in\operatorname{int}C$、$v\in C$ 且 $0\le t<1$，则
$(1-t)u+tv\in\operatorname{int}C$。事实上，取 $r>0$ 使 $B(u,r)\subset C$，则
\[
 B((1-t)u+tv,(1-t)r)=(1-t)B(u,r)+tv\subset C.
\]
若反设 $x_2\in C$，则
\[
 x_1=\left(1-\frac1m\right)x_0+\frac1m x_2,\qquad 0<\frac1m<1.
\]
由上述事实 $x_1\in\operatorname{int}C$，与 $x_1\in\partial C$ 矛盾。因此
$x_2\notin C$。

\subsection*{13.2　闭半空间表示}

\begin{problem}
设 $X$ 为实赋范线性空间，$A\subset X$ 为非空凸集。证明 $\overline A$ 是所有
包含 $A$ 的闭半空间之交。
\end{problem}

所有包含 $A$ 的闭半空间之交记为 $D$。显然 $\overline A\subset D$。反过来，若
$x_0\notin\overline A$，由严格分离定理，存在 $f\in X^*$ 和 $c\in\mathbb R$ 使
\[
 f(a)\le c<f(x_0),\qquad a\in\overline A.
\]
于是闭半空间 $H=\{x:f(x)\le c\}$ 包含 $A$ 却不包含 $x_0$，故
$x_0\notin D$。这说明 $D=\overline A$。

\subsection*{13.3　双线性映射与张量积}

\begin{problem}
设 $X,Y$ 为实赋范线性空间。证明：
\begin{enumerate}
 \item 对任意赋范线性空间 $Z$，在所有有界双线性映射
 $B:X\times Y\to Z$ 构成的空间 $\mathcal B(X,Y;Z)$ 上可定义范数
 \[
  \|B\|=\sup_{x,y\ne0}\frac{\|B(x,y)\|_Z}{\|x\|_X\|y\|_Y}.
 \]
 \item 映射
 \[
  \mathcal B(X,Y;Z)\longrightarrow\mathcal L(X,\mathcal L(Y,Z)),\qquad
  B\longmapsto(x\mapsto B(x,\cdot))
 \]
 是等距同构。
 \item 对 $(x,y)\in X\times Y$ 定义
 $x\otimes y\in\mathcal B(X,Y;\mathbb R)^*$ 为
 \[
  \langle x\otimes y,B\rangle=B(x,y),\qquad
  B\in\mathcal B(X,Y;\mathbb R).
 \]
 证明 $\|x\otimes y\|=\|x\|_X\|y\|_Y$。提示：用 Hahn--Banach 定理取适当的
 $x^*\in X^*,y^*\in Y^*$，并考虑 $B(u,v)=x^*(u)y^*(v)$。
 \item 令 $X\widehat\otimes Y\subset\mathcal B(X,Y;\mathbb R)^*$ 为包含所有
 $x\otimes y$ 的最小闭线性子空间。证明对任意赋范线性空间 $Z$，映射
 \[
  \mathcal L(X\widehat\otimes Y,Z)\longrightarrow\mathcal B(X,Y;Z),
  \qquad A\longmapsto B_A,quad B_A(x,y)=A(x\otimes y),
 \]
 是等距同构。
\end{enumerate}
\end{problem}

\textbf{(1)} 定义
\[
 \|B\|=\sup_{x,y\ne0}\frac{\|B(x,y)\|}{\|x\|\|y\|}
 =\sup_{\|x\|\le1,\|y\|\le1}\|B(x,y)\|.
\]
非负性与齐次性显然，三角不等式来自逐点三角不等式。若 $\|B\|=0$，则
$B(x,y)=0$ 对所有 $x,y$ 成立，所以这是范数。

\textbf{(2)}（题面中的目标空间应为 $\mathcal L(X,\mathcal L(Y,Z))$。）定义
\[
 \Phi(B)(x)(y)=B(x,y).
\]
则
\[
 \|\Phi(B)\|=\sup_{\|x\|\le1}\sup_{\|y\|\le1}\|B(x,y)\|=\|B\|.
\]
反之，给定 $T\in\mathcal L(X,\mathcal L(Y,Z))$，令 $B_T(x,y)=(Tx)(y)$，则
$B_T$ 有界双线性且 $\Phi(B_T)=T$，故 $\Phi$ 是等距同构。

\textbf{(3)} 泛函 $x\otimes y$ 满足
\[
 |\langle x\otimes y,B\rangle|=|B(x,y)|
 \le\|B\|\|x\|\|y\|,
\]
故 $\|x\otimes y\|\le\|x\|\|y\|$。若 $x,y\ne0$，由 Hahn--Banach
定理取 $x^*\in X^*,y^*\in Y^*$，使
\[
 \|x^*\|=\|y^*\|=1,quad x^*(x)=\|x\|,quad y^*(y)=\|y\|.
\]
令 $B_0(u,v)=x^*(u)y^*(v)$，则 $\|B_0\|=1$ 且
\[
 \|x\otimes y\|\ge|B_0(x,y)|=\|x\|\|y\|.
\]
故等号成立。

\textbf{(4)} 把题面所定义的闭线性包记为 $X\widehat\otimes_\pi Y$。给定
$A\in\mathcal L(X\widehat\otimes_\pi Y,Z)$，令 $B_A(x,y)=A(x\otimes y)$。由 (3)，
\[
 \|B_A(x,y)\|\le\|A\|\|x\|\|y\|,
\]
故 $\|B_A\|\le\|A\|$。简单张量的线性组合稠密，而 $A$ 在它们上的值完全由
$B_A$ 决定；按投射张量范数的定义，反向不等式也成立。

给定 $B\in\mathcal B(X,Y;Z)$，先定义
\[
 A_B\left(\sum_{j=1}^nx_j\otimes y_j\right)=\sum_{j=1}^nB(x_j,y_j).
\]
张量关系保证定义与表示无关，且
\[
 \left\|A_B\left(\sum_jx_j\otimes y_j\right)\right\|
 \le\|B\|\sum_j\|x_j\|\|y_j\|.
\]
对所有表示取下确界，得到 $\|A_Bu\|\le\|B\|\|u\|_\pi$，故唯一连续延拓
到完备化。显然 $B_{A_B}=B$ 且 $A_{B_A}=A$，所以映射为等距同构。

\subsection*{13.4　自反性的三空间性质}

\begin{problem}
设 $X$ 是 Banach 空间，$Y\subset X$ 是闭子空间。若 $Y$ 和商空间 $X/Y$ 都是
自反空间，证明 $X$ 也是自反空间。
\end{problem}

记商映射 $q:X\to X/Y$，规范嵌入为 $J_X$。给定 $F\in X^{**}$。限制
$F|_{Y^\perp}$ 可视为 $(X/Y)^{**}$ 中元素，因为
$q^*:(X/Y)^*\to Y^\perp$ 是等距同构。由 $X/Y$ 自反，存在
$\dot x\in X/Y$，使
\[
 F(q^*\varphi)=\varphi(\dot x),\qquad \varphi\in(X/Y)^*.
\]
取 $x_0\in X$ 使 $q(x_0)=\dot x$，令 $G=F-J_Xx_0$，则
$G|_{Y^\perp}=0$。

限制映射 $r:X^*\to Y^*$ 由 Hahn--Banach 定理满射，且核为 $Y^\perp$。
因此 $G$ 经商空间 $X^*/Y^\perp\cong Y^*$ 给出 $Y^{**}$ 中元素。因 $Y$
自反，存在 $y_0\in Y$ 使
\[
 G(x^*)=x^*(y_0),\qquad x^*\in X^*.
\]
于是 $F=J_X(x_0+y_0)$。任意 $F$ 都在 $J_X$ 的像中，故 $X$ 自反。

\subsection*{13.5　$\ell^p$ 的自反性}

\begin{problem}
设 $1<p<\infty$。证明 $\ell^p$ 是自反 Banach 空间。
\end{problem}

设 $1<p<\infty$，$q=p/(p-1)$。由 H\"older 不等式与 $\ell^p$ 对偶定理，
\[
 (\ell^p)^*\cong\ell^q,\qquad
 L_y(x)=\sum_{n=1}^\infty x_ny_n,\qquad \|L_y\|=\|y\|_q.
\]
同理 $(\ell^q)^*\cong\ell^p$。在这些识别下，规范嵌入
$J:\ell^p\to(\ell^p)^{**}$ 正是
\[
 x\longmapsto\bigl(y\mapsto\sum x_ny_n\bigr),
\]
它穷尽 $(\ell^q)^*$，故 $J$ 满射，$\ell^p$ 自反。

\section{第十四次作业}

\subsection*{14.1　James 空间}

\begin{problem}
关于 James 空间，证明 $(J,\|\cdot\|_J)$ 是完备空间，并且 $\ell^2$ 在 James 空间
$J$ 中稠密。
\end{problem}

\textbf{依赖未提供讲义，按用户要求跳过。}题面没有给出 James 空间 $J$ 及其范数
$\|\cdot\|_J$ 的定义，无法从题面唯一还原具体命题。

\subsection*{14.2　Schauder 基的部分和投影}

\begin{problem}
设 $X$ 是可分实 Banach 空间，$(e_i)_{i\ge1}$ 是 $X$ 的 Schauder 基，即对任意
$x\in X$，存在唯一数列 $(x_i)$ 使
\[
 \lim_{n\to\infty}\left\|x-\sum_{i=1}^nx_ie_i\right\|=0.
\]
定义
\[
 \Pi_n(x)=\sum_{i=1}^nx_ie_i.
\]
证明：
\begin{enumerate}
 \item $\Pi_n$ 是 $X$ 上的线性投影算子；若 $n\ge m\ge1$，则
 $\Pi_n\Pi_m=\Pi_m\Pi_n=\Pi_m$。
 \item 定义
 \[
  |||x|||=\sup_{n\ge1}\|\Pi_nx\|.
 \]
 则 $\|x\|\le|||x|||$，并且 $(X,|||\cdot|||)$ 是 Banach 空间。
 \item 存在常数 $c>0$，使
 \[
  \sup_{n\ge1}\|\Pi_nx\|\le c\|x\|,\qquad x\in X.
 \]
\end{enumerate}
\end{problem}

\textbf{(1)} 由 Schauder 展开的唯一性，坐标逐项线性，故 $\Pi_n$ 线性；且
\[
 \Pi_n^2=\Pi_n,\qquad
 \Pi_n\Pi_m=\Pi_m\Pi_n=\Pi_m\quad(n\ge m).
\]
令 $E_n=\operatorname{span}\{e_1,\dots,e_n\}$，
$F_n=\overline{\operatorname{span}\{e_{n+1},e_{n+2},\dots\}}$。Schauder 展开说明
$X=E_n\oplus F_n$。若 $x_k\to x$ 且 $\Pi_nx_k\to y$，则 $y\in E_n$，且
$x_k-\Pi_nx_k\to x-y\in F_n$，所以直和分解的唯一性给出 $y=\Pi_nx$。
故 $\Pi_n$ 的图闭；闭图定理说明 $\Pi_n\in\mathcal L(X)$。

\textbf{(2)} 对每个 $x$，$\Pi_nx\to x$，所以
\[
 |||x|||:=\sup_n\|\Pi_nx\|<\infty,\qquad \|x\|\le|||x|||.
\]
这显然是范数。若 $(x_k)$ 关于新范数为 Cauchy 列，则它关于原范数也为 Cauchy，
故 $x_k\to x\in X$。固定 $n$，由 $\Pi_n$ 连续，
$\Pi_nx_k\to\Pi_nx$。给定 $\varepsilon>0$，取 $K$ 使 $k,l\ge K$ 时
$|||x_k-x_l|||<\varepsilon$。固定 $k\ge K$ 并令 $l\to\infty$，对每个 $n$ 有
\[
 \|\Pi_n(x_k-x)\|=\lim_l\|\Pi_n(x_k-x_l)\|\le\varepsilon.
\]
取上确界得 $|||x_k-x|||\le\varepsilon$，所以新范数完备。

\textbf{(3)} 恒等映射
$I:(X,|||\cdot|||)\to(X,\|\cdot\|)$ 连续、双射。两边都是 Banach 空间，
故开映射定理说明逆映射连续。因此存在 $c>0$ 使
\[
 \sup_n\|\Pi_nx\|=|||x|||\le c\|x\|.
\]

\section{第十五次作业}

\subsection*{15.1　拓扑向量空间中的 Minkowski 泛函}

\begin{problem}
设 $X$ 为实拓扑向量空间，$K\subset X$ 是平衡凸吸收集，$P_K$ 是其 Minkowski
泛函。证明：
\begin{enumerate}
 \item
 \[
  \operatorname{Int}K\subset\{x:P_K(x)<1\}\subset K
  \subset\{x:P_K(x)\le1\}\subset\overline K.
 \]
 \item $P_K$ 连续当且仅当 $K$ 包含一个含零点的开集。此时
 \[
  \operatorname{Int}K=\{x:P_K(x)<1\},
  \qquad \overline K=\{x:P_K(x)\le1\}.
 \]
\end{enumerate}
\end{problem}

一般有 $\{P_K<1\}\subset K\subset\{P_K\le1\}$。若
$x\in\operatorname{Int}K$，连续映射 $t\mapsto tx$ 在 $t=1$ 附近仍落在 $K$，
故存在 $t>1$ 使 $tx\in K$，于是 $P_K(x)<1$。若 $P_K(x)\le1$，则
$0<t<1$ 时 $P_K(tx)<1$，故 $tx\in K$；令 $t\uparrow1$ 得 $x\in\overline K$。
所以
\[
 \operatorname{Int}K\subset\{P_K<1\}\subset K
 \subset\{P_K\le1\}\subset\overline K.
\]

若 $K$ 包含零点开邻域 $U$，则 $P_K\le1$ 于 $U$。给定 $\varepsilon>0$，
$x\in\varepsilon U$ 时 $P_K(x)\le\varepsilon$，故 $P_K$ 在 $0$ 连续，进而处处
连续。反之，若 $P_K$ 连续，则 $\{P_K<1\}$ 是含零开集并包含于 $K$。
此时包含链给出
\[
 \operatorname{Int}K=\{P_K<1\},\qquad
 \overline K=\{P_K\le1\}.
\]

\subsection*{15.2　由半范族生成的局部凸拓扑}

\begin{problem}
设 $X$ 是实线性空间，$\mathcal P=\{p_\beta\}_{\beta\in\Lambda}$ 是 $X$ 上的一族
半范。定义
\[
 U(x;\varepsilon;p_1,\dots,p_n)
 =\{y\in X:p_k(y-x)<\varepsilon, 1\le k\le n\},
\]
并令 $\mathcal U_{\mathcal P}(x)$ 为所有这类集合组成的族，
$\mathcal U_{\mathcal P}=\bigcup_{x\in X}\mathcal U_{\mathcal P}(x)$，
$\mathcal T_X=\mathcal T(\mathcal U_{\mathcal P})$。证明：
\begin{enumerate}
 \item $(X,\mathcal T_X)$ 是拓扑空间，且 $\mathcal U_{\mathcal P}$ 是
 $\mathcal T_X$ 的一组基；
 \item $(X,\mathcal T_X)$ 是局部凸空间；
 \item 每个 $p_\beta$ 都是 $(X,\mathcal T_X)$ 上的连续函数；
 \item 若 $p_\beta(x)=0$ 对所有 $\beta\in\Lambda$ 成立时必有 $x=0$，则
 $(X,\mathcal T_X)$ 是 Hausdorff 空间。
\end{enumerate}
\end{problem}

\textbf{(1)} 任意两个含 $x$ 的基本集合之交仍包含同型基本集合：合并两组半范并
取较小的 $\varepsilon$ 即可。若
$y\in U(x;\varepsilon;p_1,\dots,p_n)$，令
\[
 \delta=\min_k\{\varepsilon-p_k(y-x)\}>0.
\]
则 $U(y;\delta;p_1,\dots,p_n)$ 包含于原集合，因为
$p_k(z-x)\le p_k(z-y)+p_k(y-x)<\varepsilon$。故这些集合构成拓扑基。

\textbf{(2)} 每个零点基本邻域都是平衡凸集，所以该拓扑局部凸。并且
\[
 U(0;\varepsilon/2;p_1,\dots,p_n)+U(0;\varepsilon/2;p_1,\dots,p_n)
 \subset U(0;\varepsilon;p_1,\dots,p_n),
\]
结合半范的齐次性可验证加法、数乘连续，因此这是局部凸拓扑向量空间。

\textbf{(3)} 由反三角不等式，
\[
 |p_\beta(y)-p_\beta(x)|\le p_\beta(y-x),
\]
所以每个 $p_\beta$ 连续。

\textbf{(4)} 若 $x\ne y$，取 $p_\beta$ 使 $p_\beta(x-y)>0$，并令
$\varepsilon<p_\beta(x-y)/2$。则
$U(x;\varepsilon;p_\beta)$ 与 $U(y;\varepsilon;p_\beta)$ 不交；否则三角不等式会给出
$p_\beta(x-y)<2\varepsilon$。因此该拓扑是 Hausdorff 的。

\section{第十六次作业}

\subsection*{16.1　Hilbert 空间中弱拓扑的非序列现象}

\begin{problem}
设 $H$ 为无穷维可分实 Hilbert 空间，$(e_n)_{n\in\mathbb N}$ 是一组标准正交基。
证明：
\begin{enumerate}
 \item $e_n\rightharpoonup0$；
 \item 集合 $A=\{\sqrt n\,e_n:n\in\mathbb N\}$ 是序列弱闭的，但其弱闭包包含
 $0$。提示：若 $U$ 是 $0$ 的弱开邻域，则存在 $\varepsilon>0$ 和
 $y_1,\dots,y_m\in H$，使
 \[
  V=\{x\in H:\max_{1\le i\le m}|\langle x,y_i\rangle|<\varepsilon\}\subset U.
 \]
 证明 $z_n:=\max_i|\langle e_n,y_i\rangle|$ 平方可加，并推出 $V\cap A\ne\varnothing$。
 \item $A$ 中的任何子列都不弱收敛到 $0$。
\end{enumerate}
\end{problem}

\textbf{(1)} 对任意 $y\in H$，Bessel 不等式给出
\[
 \sum_{n=1}^\infty|\langle e_n,y\rangle|^2\le\|y\|^2.
\]
故 $\langle e_n,y\rangle\to0$。这正是 $e_n\rightharpoonup0$。

\textbf{(2)} 先证 $A=\{\sqrt n,e_n:n\in\mathbb N\}$ 序列弱闭。若
$\sqrt{n_k}e_{n_k}$ 弱收敛，则弱收敛序列必范数有界，因此 $(\sqrt{n_k})$ 有界，
即整数列 $(n_k)$ 只取有限多个值。一个取值于有限集的收敛序列最终必须停留在一个点，
其极限仍属于 $A$。

再证 $0\in\overline A^{,w}$。给定 $0$ 的弱邻域 $U$，可取
$y_1,\dots,y_m\in H$ 及 $\varepsilon>0$，使
\[
 V=\{x:\max_{1\le i\le m}|\langle x,y_i\rangle|<\varepsilon\}\subset U.
\]
令 $z_n=\max_i|\langle e_n,y_i\rangle|$。则
\[
 \sum_nz_n^2\le\sum_{i=1}^m\sum_n|\langle e_n,y_i\rangle|^2
 \le\sum_{i=1}^m\|y_i\|^2<\infty.
\]
必存在无穷多个 $n$ 使 $\sqrt n,z_n<\varepsilon$；否则最终
$z_n^2\ge\varepsilon^2/n$，与平方可加矛盾。对这样的 $n$，
$\sqrt n,e_n\in V\cap A\subset U\cap A$。故每个弱邻域都与 $A$ 相交。

\textbf{(3)} 取 $A$ 中任意由互异指标组成的子列
$\sqrt{n_k}e_{n_k}$。可再取子列使 $n_{k_j}\ge2^j$。定义
\[
 y=\sum_{j=1}^\infty n_{k_j}^{-1/2}e_{n_{k_j}}.
\]
因为 $\sum_jn_{k_j}^{-1}\le\sum_j2^{-j}<\infty$，故 $y\in H$。但
\[
 \langle\sqrt{n_{k_j}}e_{n_{k_j}},y\rangle=1,
\]
所以原子列不可能弱收敛到 $0$。若允许重复指标的子列，则恒定子列显然也不弱收敛到
$0$。

\subsection*{16.2　$\ell^1$ 的 Schur 性质}

\begin{problem}
\begin{enumerate}
 \item 证明 $\ell^1$ 中的标准基 $e_n=(\delta_{in})_{i\ge1}$ 不弱收敛到 $0$；
 \item 把 $\ell^1$ 看成 $c_0$ 的对偶空间，证明 $e_n$ 弱星收敛到 $0$；
 \item 证明 Schur 定理：$\ell^1$ 中的序列弱收敛到 $0$，当且仅当它按范数收敛到
 $0$。提示：使用滑动驼峰技巧。
\end{enumerate}
\end{problem}

\textbf{(1)} 把 $\ell^\infty$ 识别为 $(\ell^1)^*$。取
$\mathbf1=(1,1,\dots)\in\ell^\infty$，则
\[
 \langle e_n,\mathbf1\rangle=1
\]
不趋于 $0$，故 $e_n$ 不弱收敛到 $0$。

\textbf{(2)} 把 $\ell^1$ 视为 $c_0^*$。对任意 $x=(x_i)\in c_0$，
\[
 \langle e_n,x\rangle=x_n\to0.
\]
因此 $e_n\overset{w^*}\longrightarrow0$。

\textbf{(3)} 范数收敛当然推出弱收敛。下面证明反向。设
$x^{(j)}\rightharpoonup0$ 于 $\ell^1$，但不按范数收敛。取子列（仍记为
$x^{(j)}$）及 $\eta>0$，使 $\|x^{(j)}\|_1\ge\eta$。
弱收敛蕴含逐坐标收敛，因此可以递归选取
\[
 0=N_0<N_1<N_2<\cdots,\qquad j_1<j_2<\cdots,
\]
满足
\[
 \sum_{i\le N_{k-1}}|x_i^{(j_k)}|<\eta/8,\qquad
 \sum_{i>N_k}|x_i^{(j_k)}|<\eta/8.
\]
第一式靠先选充分大的 $j_k$，第二式靠随后选充分大的 $N_k$。于是第 $k$ 个块
$I_k=(N_{k-1},N_k]$ 上有
\[
 \sum_{i\in I_k}|x_i^{(j_k)}|\ge3\eta/4.
\]
定义 $a=(a_i)\in\ell^\infty$：在 $I_k$ 上令
$a_i=\operatorname{sgn}(x_i^{(j_k)})$。则 $\|a\|_\infty\le1$，并且
\[
 \langle x^{(j_k)},a\rangle
 \ge\sum_{i\in I_k}|x_i^{(j_k)}|
 -\sum_{i\notin I_k}|x_i^{(j_k)}|
 \ge\frac\eta2.
\]
这与 $x^{(j)}\rightharpoonup0$ 矛盾。故弱收敛必推出范数收敛。一般弱收敛到
$x$ 的情形应用于 $x^{(j)}-x$ 即可。

\subsection*{16.3　讲义第 101 页定理 4}

\begin{problem}
证明讲义第 101 页定理 4 的结论 (3)、(4)。
\end{problem}

\textbf{依赖未提供讲义，按用户要求跳过。}题面未给出定理 4 的陈述及结论 (3)、(4)。

\section{第十七次作业}

\subsection*{17.1　对偶单位球上弱星拓扑的可度量化}

\begin{problem}
设 $X$ 是可分赋范线性空间，$(x_n)_{n\in\mathbb N}$ 是 $X$ 的闭单位球中的稠密
序列。在 $X^*$ 的闭单位球 $B^*$ 上定义
\[
 d(x^*,y^*)=\sum_{n=1}^\infty2^{-n}
 |\langle x^*-y^*,x_n\rangle|.
\]
证明 $d$ 是 $B^*$ 上的距离，并证明该距离诱导的拓扑就是 $B^*$ 上的弱星拓扑。
\end{problem}

对 $x^*,y^*\in B^*$，级数
\[
 d(x^*,y^*)=\sum_{n=1}^\infty2^{-n}|\langle x^*-y^*,x_n\rangle|
\]
绝对收敛，因为每项绝对值不超过 $2$。非负性、对称性、三角不等式逐项成立。
若 $d(x^*,y^*)=0$，则二者在稠密集 $\{x_n\}$ 上取值相同。对任意单位球中的
$x$，取 $x_{n_k}\to x$，由 $\|x^*-y^*\|\le2$，
\[
 |\langle x^*-y^*,x\rangle|
 \le2\|x-x_{n_k}\|+|\langle x^*-y^*,x_{n_k}\rangle|\to0.
\]
再用齐次性推广到所有 $x\in X$，得 $x^*=y^*$。故 $d$ 是距离。

若 $x_\alpha^*\overset{w^*}\to x^*$ 于 $B^*$，给定 $\varepsilon>0$，先取 $N$
使 $2\sum_{n>N}2^{-n}<\varepsilon/2$，再由前 $N$ 个坐标收敛使
\[
 \sum_{n=1}^N2^{-n}|\langle x_\alpha^*-x^*,x_n\rangle|<\varepsilon/2.
\]
故 $d(x_\alpha^*,x^*)<\varepsilon$。

反过来，若 $d(x_\alpha^*,x^*)\to0$，则对固定 $n$，
\[
 |\langle x_\alpha^*-x^*,x_n\rangle|\le2^n d(x_\alpha^*,x^*)\to0.
\]
给定 $x\in B_X$，选 $x_n$ 使 $\|x-x_n\|<\varepsilon/4$，则
\[
 |\langle x_\alpha^*-x^*,x\rangle|
 \le2\|x-x_n\|+|\langle x_\alpha^*-x^*,x_n\rangle|<\varepsilon
\]
最终成立。故在每个 $x$ 上点态收敛，即弱星收敛。两种拓扑一致。

\subsection*{17.2　$C_0(M)$、点质量与弱星闭包}

\begin{problem}
设 $M$ 为局部紧 Hausdorff 空间。称连续实值函数 $f:M\to\mathbb R$ 在无穷远处
消失，如果对任意 $\varepsilon>0$，存在紧集 $K\subset M$，使
$\sup_{x\in M\setminus K}|f(x)|<\varepsilon$。记所有这类函数组成的空间为
$C_0(M)$。
\begin{enumerate}
 \item 证明 $C_0(M)$ 在上确界范数下是 Banach 空间。
 \item 根据 Riesz 表示定理，$C_0(M)^*$ 与符号 Radon 测度空间对应；这里符号
 Radon 测度满足：对每个 Borel 集 $B\subset M$ 和 $\varepsilon>0$，存在紧集
 $K\subset B$，使对任意 Borel 集 $A\subset B$，
 $|\mu(A)-\mu(A\cap K)|<\varepsilon$。

 定义
 \[
  \delta:M\to C_0(M)^*,\qquad \delta(x)=\delta_x,quad
  \delta_x(f)=f(x).
 \]
 在 $C_0(M)^*$ 上装备弱星拓扑。证明
 $\delta:M\to\delta(M)$ 是同胚。点质量集合包含于 Radon 概率测度集
 \[
  P(M)=\{\mu:\mu\ge0, \|\mu\|=\mu(M)=1\}.
 \]
 求 $\delta(M)=\{\delta_x:x\in M\}$ 在 $P(M)$ 中的弱星闭包。
\end{enumerate}
\end{problem}

\textbf{(1)} 设 $f_n$ 是 $C_0(M)$ 中的一致 Cauchy 列。它在有界连续函数空间中
一致收敛到某个连续函数 $f$。给定 $\varepsilon>0$，取 $n$ 使
$\|f-f_n\|_\infty<\varepsilon/2$；再取紧集 $K$ 使
$|f_n(x)|<\varepsilon/2$ 于 $M\setminus K$。于是 $|f(x)|<\varepsilon$ 于
$M\setminus K$，故 $f\in C_0(M)$。所以 $C_0(M)$ 完备。

\textbf{(2)} 对每个 $x$，$\delta_x$ 是正泛函且 $\|\delta_x\|=1$，所以对应于
Radon 概率测度（仍记为 $\delta_x$）。若 $x_\alpha\to x$ 于 $M$，则对每个
$f\in C_0(M)$，
\[
 \delta_{x_\alpha}(f)=f(x_\alpha)\to f(x)=\delta_x(f),
\]
故 $\delta$ 连续。

局部紧 Hausdorff 空间有如下分离性质：给定 $x\in U$（$U$ 开），存在
$f\in C_c(M)$ 使 $0\le f\le1$、$f(x)=1$ 且 $\operatorname{supp}f\subset U$。
若 $\delta_{x_\alpha}\overset{w^*}\to\delta_x$，取这样的 $f$，则最终
$f(x_\alpha)>1/2$，从而 $x_\alpha\in U$。故 $x_\alpha\to x$，说明逆映射连续，
$\delta:M\to\delta(M)$ 是同胚。

下面求闭包。若 $\delta_{x_\alpha}\overset{w^*}\to\mu\in P(M)$，则对每个实值
$f\in C_0(M)$，
\[
 \int f^2d\mu=\lim_\alpha f(x_\alpha)^2
 =\left(\lim_\alpha f(x_\alpha)\right)^2
 =\left(\int f,d\mu\right)^2.
\]
因此每个 $f$ 关于概率测度 $\mu$ 的方差均为零，即 $f$ 在 $\mu$ 几乎处处为常数。
若 $\mu$ 的支集含两个不同点，可用 $C_c(M)$ 函数把它们分开，并在两个正测度邻域上
取不同值，方差便为正，矛盾。故 $\mu$ 的支集是单点，$\mu=\delta_x$。所以
\[
 \overline{\delta(M)}^{,w^*}\cap P(M)=\delta(M).
\]
若在整个 $C_0(M)^*$ 中取闭包，则 $M$ 非紧时还有 $0$：取逃离每个紧集的网
$x_K$，便有 $f(x_K)\to0$ 对所有 $f\in C_0(M)$。同样的平方恒等式说明不会出现
别的极限。因此整个对偶中的闭包在 $M$ 紧时为 $\delta(M)$，在 $M$ 非紧时为
$\delta(M)\cup\{0\}$。

\subsection*{17.3　可分对偶下弱闭包中的序列逼近}

\begin{problem}
设 $X$ 是 Banach 空间且 $X^*$ 可分，$S\subset X$ 是有界集，
$x\in\overline S^{,w}$。证明存在序列 $(x_n)\subset S$，使
$x_n\rightharpoonup x$。
\end{problem}

设 $S\subset X$ 有界，取 $M>0$ 使 $\|s\|\le M$ 对所有 $s\in S$ 成立。
弱闭包仍包含于 $MB_X$，故 $\|x\|\le M$。因为 $X^*$ 可分，可在其单位球中取
稠密序列 $(f_k)$。对每个 $n$，集合
\[
 U_n=\{y\in X:|f_k(y-x)|<1/n, 1\le k\le n\}
\]
是 $x$ 的弱邻域。由 $x\in\overline S^{,w}$，可取 $x_n\in S\cap U_n$。

固定 $f\in X^*$ 和 $\varepsilon>0$。若 $f=0$ 无须证明；否则选 $k$ 使
$\|f/\|f\|-f_k\|<\varepsilon/(4M\|f\|)$（把常数作显然调整亦可）。当
$n\ge k$ 时，
\[
 |f(x_n-x)|
 \le \|f\|\,|f_k(x_n-x)|
 +\|f-\|f\|f_k\|\,\|x_n-x\|.
\]
第一项随 $n\to\infty$ 趋于 $0$，第二项由 $\|x_n-x\|\le2M$ 控制并可任意小。
所以 $f(x_n)\to f(x)$。这对所有 $f\in X^*$ 成立，即 $x_n\rightharpoonup x$。

\section{第十八次作业}

\subsection*{18.1　紧开拓扑与 Arzel\`a--Ascoli 定理}

\begin{problem}
设 $X$ 为拓扑空间，$Y$ 为度量空间，$C(X,Y)$ 为从 $X$ 到 $Y$ 的连续映射空间。
紧开拓扑（记为 $\mathcal T_{\mathrm{co}}$）是以
\[
 L(K,V)=\{f\in C(X,Y):f(K)\subset V\}
\]
为子基的最小拓扑，其中 $K\subset X$ 紧、$V\subset Y$ 开。

首先证明：$U\subset C(X,Y)$ 在紧开拓扑中为开集，当且仅当对每个 $f\in U$，
存在有限多个紧集 $K_1,\dots,K_m\subset X$ 和开集 $V_1,\dots,V_m\subset Y$，使
\[
 f\in\bigcap_{i=1}^mL(K_i,V_i)\subset U.
\]
进而证明：
\begin{enumerate}
 \item 若 $X$ 紧，则紧开拓扑与度量
 \[
  d(f,g)=\sup_{x\in X}d_Y(f(x),g(x))
 \]
 诱导的拓扑一致。

 提示一：若 $f(K_i)\subset V_i$，证明存在 $\varepsilon>0$，使
 $B_\varepsilon(f(x))\subset V_i$ 对所有 $x\in K_i$、$1\le i\le m$ 成立，
 从而 $d(f,g)<\varepsilon$ 蕴含 $g(K_i)\subset V_i$。

 提示二：给定 $f\in C(X,Y)$ 与 $\varepsilon>0$，取有限多个
 $x_1,\dots,x_m\in X$，使
 \[
  X=\bigcup_{i=1}^mK_i,\qquad
  K_i=\{x\in X:d_Y(f(x_i),f(x))\le\varepsilon/4\}.
 \]
 令 $V_i=B_{\varepsilon/2}(f(x_i))$、
 $U=\{g:g(K_i)\subset V_i, 1\le i\le m\}$，证明
 $f\in U$ 且 $d(f,g)<\varepsilon$ 对 $g\in U$ 成立。

 \item 对每个紧集 $K\subset X$ 定义
 \[
  p_K(f)=\sup_{x\in K}|f(x)|,\qquad f\in C(X,\mathbb R).
 \]
 证明这些半范数生成紧开拓扑，即紧开拓扑是使每个 $p_K$ 连续的最小拓扑。
 \item 证明 $(C(X,\mathbb R),\mathcal T_{\mathrm{co}})$ 是局部凸拓扑
 向量空间。
 \item 证明 $F\subset C(X,Y)$ 在紧开拓扑中预紧，当且仅当对每个紧集
 $K\subset X$，限制族
 \[
  F_K=\{f|_K:f\in F\}\subset C(K,Y)
 \]
 预紧。提示：令 $\mathcal K$ 为 $X$ 的紧子集全体，证明映射
 \[
  C(X,Y)\to\prod_{K\in\mathcal K}C(K,Y),\qquad
  f\mapsto(f|_K)_{K\in\mathcal K}
 \]
 是到其像的同胚，并应用 Tychonoff 定理。
 \item 证明如下 Arzel\`a--Ascoli 定理的变形：$F\subset C(X,Y)$ 在紧开拓扑中
 预紧，当且仅当 $F$ 逐点预紧，并且对任意紧集 $K\subset X$，$F_K$ 等度连续。
\end{enumerate}
\end{problem}

先说明紧开拓扑的基。按定义，集合 $L(K,V)$ 构成一个子基，所以它们的有限交
\[
 \bigcap_{i=1}^mL(K_i,V_i)
\]
构成一个基。因此 $U$ 开当且仅当对每个 $f\in U$，存在有限多个 $K_i,V_i$，使
\[
 f\in\bigcap_{i=1}^mL(K_i,V_i)\subset U.
\]
这正是题面所述条件。

\textbf{(1)} 设 $X$ 紧。先证每个紧开基本邻域包含一个一致度量球。若
$f(K_i)\subset V_i$，则 $f(K_i)$ 是紧集且含于开集 $V_i$。函数
$z\mapsto d_Y(z,Y\setminus V_i)$ 在 $f(K_i)$ 上取得正的最小值 $\eta_i$；若
$Y\setminus V_i=\varnothing$，该项可忽略。取
$0<\varepsilon<\min_i\eta_i$。若 $d(f,g)<\varepsilon$，则对 $x\in K_i$，
$g(x)\in V_i$，故
\[
 B_d(f,\varepsilon)\subset\bigcap_iL(K_i,V_i).
\]
所以紧开拓扑不强于一致拓扑。

反过来，给定 $f$ 和 $\varepsilon>0$。紧集 $f(X)$ 可由有限多个
$B_{\varepsilon/4}(f(x_i))$ 覆盖。令
\[
 K_i=\{x\in X:d_Y(f(x),f(x_i))\le\varepsilon/4\}.
\]
它们是紧集并覆盖 $X$。置
\[
 V_i=B_{\varepsilon/2}(f(x_i)),\qquad
 U=\bigcap_{i=1}^mL(K_i,V_i).
\]
则 $f\in U$。若 $g\in U$ 且 $x\in K_i$，则
\[
 d_Y(f(x),g(x))
 \le d_Y(f(x),f(x_i))+d_Y(f(x_i),g(x))
 <\varepsilon/4+\varepsilon/2<\varepsilon.
\]
故 $U\subset B_d(f,\varepsilon)$。两种拓扑一致。

\textbf{(2)} 现在取 $Y=\mathbb R$。半范族
\[
 p_K(f)=\sup_{x\in K}|f(x)|
\]
生成的基本邻域是有限交
\[
 \{g:p_{K_i}(g-f)<\varepsilon_i, 1\le i\le m\}.
\]
在每个紧集 $K_i$ 上，(1) 的证明说明一致拓扑恰为紧开拓扑；把有限多个条件取交，
即知这些半范邻域与紧开基本邻域互相包含。因此半范族生成的拓扑就是紧开拓扑。

也可直接验证连续性的最小性：若某拓扑使所有 $p_K$ 连续，则上述半范球均开，因而
它至少与紧开拓扑一样强；反之，(1) 应用于每个紧子空间 $K$ 表明每个 $p_K$ 在紧开
拓扑下连续。

\textbf{(3)} 每个 $p_K$ 是半范。由半范族生成的拓扑自动有平衡凸的零点邻域基。
加法连续性来自
\[
 p_K((f_1+g_1)-(f_2+g_2))
 \le p_K(f_1-f_2)+p_K(g_1-g_2),
\]
数乘连续性来自
\[
 p_K(af-bg)\le |a|p_K(f-g)+|a-b|p_K(g).
\]
故 $C(X,\mathbb R)$ 是局部凸拓扑向量空间。

\textbf{(4)} 这里“预紧”按相应一致结构中的全有界性理解。令 $\mathcal K$ 为 $X$
中所有紧子集，考虑限制映射
\[
 R:C(X,Y)\longrightarrow\prod_{K\in\mathcal K}C(K,Y),\qquad
 R(f)=(f|_K)_{K\in\mathcal K}.
\]
它是单射。产品中的有限多个坐标邻域拉回后恰为紧开邻域，故 $R$ 是到其像的同胚，
并保持相应的一致结构。

若 $F$ 预紧，则每个连续坐标投影把 $F$ 映为预紧集，所以每个 $F_K$ 预紧。
反之，若所有 $F_K$ 预紧，则其完备化中的闭包 $\overline{F_K}$ 紧。由 Tychonoff
定理，$\prod_K\overline{F_K}$ 紧，故产品中的 $R(F)$ 预紧；拉回即知 $F$ 在紧开
一致结构中预紧。

\textbf{(5)} 先假设 $F$ 在紧开拓扑中预紧。对单点紧集 $\{x\}$，(4) 给出
$\{f(x):f\in F\}$ 预紧，故 $F$ 逐点预紧。对任意紧集 $K$，$F_K$ 在一致度量下
预紧。给定 $\varepsilon>0$，可取有限个 $g_1,\dots,g_N\in C(K,Y)$，使每个
$f|_K$ 与某个 $g_j$ 的一致距离小于 $\varepsilon/3$。有限族连续函数
$g_1,\dots,g_N$ 等度连续：对每个 $x\in K$，存在邻域 $U_x$ 使对所有 $j$ 和
$y\in U_x\cap K$，
$d_Y(g_j(x),g_j(y))<\varepsilon/3$。于是对相应的 $f$，
\[
 d_Y(f(x),f(y))
 \le d_Y(f(x),g_j(x))+d_Y(g_j(x),g_j(y))+d_Y(g_j(y),f(y))
 <\varepsilon.
\]
故 $F_K$ 等度连续。

反之，假设 $F$ 逐点预紧，且每个 $F_K$ 等度连续。固定紧集 $K$ 和
$\varepsilon>0$。由等度连续性，对每个 $x\in K$ 可取邻域 $U_x$，使
\[
 d_Y(f(x),f(y))<\varepsilon/3
 \quad(f\in F, y\in U_x\cap K).
\]
紧性给出有限覆盖 $U_{x_1},\dots,U_{x_m}$。每个点值集
$\{f(x_i):f\in F\}$ 预紧，故有有限 $\varepsilon/3$ 网。把一个函数 $f$ 按它在
$x_1,\dots,x_m$ 处所落入的有限网格分类，只有有限种类型。同一类型的两个函数
$f,g$，对 $y\in U_{x_i}$ 有
\[
 d_Y(f(y),g(y))
 \le d_Y(f(y),f(x_i))+d_Y(f(x_i),g(x_i))+d_Y(g(x_i),g(y))<\varepsilon.
\]
每个非空类型选一个代表，便得到 $F_K$ 的有限 $\varepsilon$ 网。所以 $F_K$ 预紧。
由 (4)，$F$ 在紧开拓扑中预紧。

\subsection*{18.2　由 Arzel\`a--Ascoli 推出 Banach--Alaoglu}

\begin{problem}
设 $X$ 是赋范线性空间。利用 18.1 中的 Arzel\`a--Ascoli 定理变形推出
Banach--Alaoglu 定理。

提示：把 $X^*$ 的闭单位球视为 $C(X,\mathbb R)$ 的子集，它是等度连续的；证明
$X^*$ 上的紧开拓扑强于弱星拓扑，即每个弱星开集在紧开拓扑中仍为开集。
\end{problem}

把对偶闭单位球
\[
 B^*=\{x^*\in X^*:\|x^*\|\le1\}
\]
视为 $C(X,\mathbb R)$ 的子集。它逐点预紧，因为对固定 $x\in X$，
\[
 \{x^*(x):x^*\in B^*\}\subset[-\|x\|,\|x\|].
\]
它还是等度连续的：对所有 $x^*\in B^*$，
\[
 |x^*(x)-x^*(y)|\le\|x-y\|.
\]
因而在每个紧集上的限制族都等度连续。由 18.1(5)，$B^*$ 在紧开拓扑中预紧。

$B^*$ 在紧开拓扑中是闭的。事实上，若一个网 $x_\alpha^*\in B^*$ 在每个紧集上一致
收敛到 $f$，特别地它逐点收敛；逐点极限保持线性，且
\[
 |f(x)|=\lim_\alpha|x_\alpha^*(x)|\le\|x\|,
\]
故 $f\in B^*$。于是 $B^*$ 在紧开拓扑下紧。

弱星拓扑是点态收敛拓扑；每个弱星子基本开集只涉及有限多个点，而有限集是紧集，
所以弱星拓扑弱于紧开拓扑。恒等映射
\[
 (B^*,\mathcal T_{\mathrm{co}})
 \longrightarrow(B^*,w^*)
\]
连续。紧空间的连续像紧，因此 $B^*$ 在弱星拓扑下紧。这就是 Banach--Alaoglu
定理（实数情形；复数情形对实部、虚部分别处理完全相同）。

\chapter{第十九次至第二十七次作业}
\section{第十九次作业}

\subsection*{19.1}
\begin{problem}
若 $X$ 是自反的 Banach 空间，则规范嵌入映射 $\iota$ 满足
$\iota:(X,\mathcal U_w)\to(X^{**},\mathcal U_w)$ 是同胚映射。
\end{problem}

设规范嵌入
\[
 \iota:X\longrightarrow X^{**},\qquad \iota(x)(x^*)=x^*(x)
\]
关于两边的弱拓扑是同胚。

\begin{proof}
由于 $X$ 自反，$\iota$ 是从 $X$ 到 $X^{**}$ 的等距双射。先证明其弱连续性。若网
$x_\alpha\rightharpoonup x$ 于 $X$，则对每个 $x^*\in X^*$，
\[
 \iota(x_\alpha)(x^*)=x^*(x_\alpha)\longrightarrow x^*(x)=\iota(x)(x^*).
\]
这正是 $\iota(x_\alpha)\rightharpoonup\iota(x)$ 于 $X^{**}$。

反过来，若 $\iota(x_\alpha)\rightharpoonup\iota(x)$ 于 $X^{**}$，则弱收敛的定义要求对每个
$F\in X^{***}$ 有 $F(\iota(x_\alpha))\to F(\iota(x))$。特别取规范嵌入
$\iota_{X^*}(x^*)\in X^{***}$，便有
\[
 x^*(x_\alpha)=\iota_{X^*}(x^*)(\iota(x_\alpha))
 \longrightarrow \iota_{X^*}(x^*)(\iota(x))=x^*(x).
\]
故 $x_\alpha\rightharpoonup x$。所以 $\iota$ 与 $\iota^{-1}$ 都连续，结论成立。
\end{proof}

\subsection*{19.2（Milman--Pettis 定理）}
\begin{problem}
证明 Milman--Pettis 定理：一致凸 Banach 空间是自反的。使用网的概念，按以下步骤证明：

有向集是非空集合 $A$，装备一个自反且传递的关系 $\preceq$，并且对任意
$\alpha,\beta\in A$，都存在 $\gamma\in A$ 使
$\alpha\preceq\gamma$、$\beta\preceq\gamma$。关系不要求反对称。例如，拓扑空间中一点
$x_0$ 的开邻域全体在关系 $U\preceq V\Longleftrightarrow V\subset U$ 下构成有向集。

空间 $X$ 中的网是映射 $A\to X$，记作 $(x_\alpha)_{\alpha\in A}$。称该网在拓扑空间
$X$ 中收敛于 $x$，若对 $x$ 的每个开邻域 $U$，都存在 $\alpha_0\in A$，使
$\alpha_0\preceq\alpha$ 时 $x_\alpha\in U$。映射 $f:X\to Y$ 连续，当且仅当它把每个收敛于
$x$ 的网送到收敛于 $f(x)$ 的网。

网 $(x_\alpha)_{\alpha\in A}$ 的子网是形如
$(x_{h(\beta)})_{\beta\in B}$ 的网，其中 $h:B\to A$ 是单调最终映射：单调是指
$\beta_1\preceq\beta_2$ 蕴含 $h(\beta_1)\preceq h(\beta_2)$；最终是指对每个
$\alpha\in A$，都存在 $\beta\in B$ 使 $\alpha\preceq h(\beta)$。拓扑空间紧，当且仅当其中
每个网都有收敛子网。

赋范空间中的网 $(x_\alpha)$ 称为 Cauchy 网，如果
$(\|x_\alpha-x_\beta\|)_{(\alpha,\beta)\in A\times A}$ 收敛于 $0$，其中 $A\times A$ 取乘积序。
若 $X$ 是 Banach 空间，则 $X$ 中每个 Cauchy 网都收敛。

（1）设 $X$ 为一致凸赋范线性空间，$(x_\alpha)_{\alpha\in A}$ 为 $X$ 单位球中的网，且
网 $(\|x_\alpha+x_\beta\|)_{(\alpha,\beta)\in A\times A}$ 收敛于 $2$，证明
$(x_\alpha)_{\alpha\in A}$ 是 Cauchy 网；

（2）设 $X$ 为赋范线性空间，$x^{**}\in X^{**}$ 且 $\|x^{**}\|=1$，证明存在 $X$ 单位球中的网
$(x_\alpha)_{\alpha\in A}$，使 $(\iota(x_\alpha))_{\alpha\in A}$ 在 $X^{**}$ 中按弱星拓扑收敛于
$x^{**}$；

（3）设 $(x_\alpha)_{\alpha\in A}$ 为 $X$ 单位球中的网，且
$\iota(x_\alpha)\overset{w^*}{\longrightarrow}x^{**}$、$\|x^{**}\|=1$。证明
$(\iota(x_\alpha+x_\beta))_{(\alpha,\beta)\in A\times A}$ 按弱星拓扑收敛于 $2x^{**}$；若
$X$ 一致凸，利用（1）证明 $(x_\alpha)$ 是 Cauchy 网；

（4）设 $X$ 是一致凸 Banach 空间，$x^{**}\in X^{**}$ 且 $\|x^{**}\|=1$。按（2）取网
$(x_\alpha)$，由（3）知它收敛于某个 $x\in X$，从而推出 $\iota(x)=x^{**}$。
\end{problem}

\begin{proof}
记 $B_X=\{x:\|x\|\leq1\}$。

\textbf{(1)} 给定 $\varepsilon>0$。一致凸性说明存在 $\delta>0$，使得 $u,v\in B_X$ 且
$\|u-v\|\geq\varepsilon$ 时
\[
 \left\|\frac{u+v}{2}\right\|\leq1-\delta,
 \quad\text{即}\quad \|u+v\|\leq2(1-\delta).
\]
已知 $\|x_\alpha+x_\beta\|\to2$，故存在 $\alpha_0$，当
$\alpha,\beta\succeq\alpha_0$ 时
$\|x_\alpha+x_\beta\|>2(1-\delta)$。由上述命题的逆否命题，
$\|x_\alpha-x_\beta\|<\varepsilon$。所以 $(x_\alpha)$ 是 Cauchy 网。

\textbf{(2)} 由 Goldstine 定理，$\iota(B_X)$ 在 $B_{X^{**}}$ 中关于弱星拓扑稠密。
为把稠密性明确写成网，令
\[
 \mathcal A=\{(U,x):U\text{ 是 }x^{**}\text{ 的弱星邻域},\ x\in B_X,
 \ \iota(x)\in U\}.
\]
以邻域反向包含为主序：$(U,x)\preceq(V,y)$ 当 $V\subset U$。任意两个指标都有共同上界，
故 $\mathcal A$ 是有向集。令 $x_{(U,x)}=x$。对任意弱星邻域 $U_0$，从指标中取一个
$U\subset U_0$ 后，所有更大的指标所对应的像都落在 $U_0$ 中。因此
$\iota(x_\alpha)\overset{w^*}{\longrightarrow}x^{**}$。

\textbf{(3)} 对任意 $x^*\in X^*$，
\[
 \iota(x_\alpha+x_\beta)(x^*)=x^*(x_\alpha)+x^*(x_\beta)
 \longrightarrow2x^{**}(x^*),
\]
所以双指标网 $\iota(x_\alpha+x_\beta)\overset{w^*}{\longrightarrow}2x^{**}$。
范数关于弱星拓扑下半连续，因而
\[
 2=\|2x^{**}\|\leq\liminf_{(\alpha,\beta)}\|x_\alpha+x_\beta\|.
\]
另一方面 $x_\alpha,x_\beta\in B_X$，故这些范数不超过 $2$。于是
$\|x_\alpha+x_\beta\|\to2$。由 (1)，$(x_\alpha)$ 是 Cauchy 网。

\textbf{(4)} 因 $X$ 完备，存在 $x\in X$ 使 $x_\alpha\to x$（范数收敛）。于是
$\iota(x_\alpha)\to\iota(x)$ 也范数收敛，从而弱星收敛。弱星拓扑是 Hausdorff 的，而同一网
已经弱星收敛到 $x^{**}$，故 $\iota(x)=x^{**}$。

上面处理了 $\|x^{**}\|=1$ 的情形。零元显然在 $\iota(X)$ 中；对任意非零 $x^{**}$，把
$x^{**}/\|x^{**}\|$ 应用上述结论后再乘 $\|x^{**}\|$，可知 $x^{**}\in\iota(X)$。
因此 $\iota$ 满射，$X$ 自反。
\end{proof}

\section{第二十次作业}

\subsection*{20.1}
\begin{problem}
设 $X,Y$ 为 Banach 空间，在 $\mathcal L(X,Y)$ 上考虑算子范数。证明：
（1）所有满射构成的集合是开集；
（2）所有单射构成的集合不一定是开集；
（3）所有具有闭像的单射构成的集合是开集。
\end{problem}

\begin{proof}
\textbf{(1) 满射算子组成开集。}设 $T:X\to Y$ 满射。由开映射定理，存在 $c>0$，使每个
$y\in Y$ 都能找到 $x\in X$ 满足
\[
 Tx=y,\qquad \|x\|\leq c\|y\|.
\]
取 $S$ 满足 $\|S-T\|<1/c$。给定 $y$，令 $r_0=y$；递归选择 $x_n$ 使
$Tx_n=r_n$ 且 $\|x_n\|\leq c\|r_n\|$，再令 $r_{n+1}=(T-S)x_n$。若
$q=c\|T-S\|<1$，则
\[
 \|r_{n+1}\|\leq q\|r_n\|,\qquad \|x_n\|\leq c q^n\|y\|.
\]
故 $x=\sum_{n\geq0}x_n$ 在 $X$ 中收敛，而且
\[
 S\sum_{n=0}^N x_n=\sum_{n=0}^N(Tx_n-r_{n+1})=y-r_{N+1}\to y.
\]
于是 $Sx=y$，所以 $S$ 满射。

\textbf{(2) 单射算子组成的集合未必开。}在 $\ell^2$ 上令
$T(x_n)=(x_n/n)$。$T$ 单射。对每个 $N$ 定义
\[
 T_N(x_1,x_2,\ldots)=(x_1,x_2/2,\ldots,x_{N-1}/(N-1),0,x_{N+1}/(N+1),\ldots).
\]
则 $T_Ne_N=0$，故 $T_N$ 非单射；而 $\|T_N-T\|=1/N\to0$。

\textbf{(3) 闭像单射组成开集。}若 $T$ 单射且 $\operatorname{Ran}T$ 闭，则
$T:X\to\operatorname{Ran}T$ 是 Banach 空间间的有界双射。由有界逆定理，存在 $m>0$ 使
$\|Tx\|\geq m\|x\|$。若 $\|S-T\|<m$，则
\[
 \|Sx\|\geq\|Tx\|-\|(S-T)x\|\geq(m-\|S-T\|)\|x\|.
\]
因此 $S$ 单射且有正下界。若 $Sx_n\to y$，上式使 $(x_n)$ 为 Cauchy 列；其极限 $x$ 满足
$Sx=y$，故 $\operatorname{Ran}S$ 闭。
\end{proof}

\subsection*{20.2}
\begin{problem}
设 $X,Y$ 为 Banach 空间，$T,S\in\mathcal L(X,Y)$。证明：
（1）$(T+S)^*=T^*+S^*$，且 $(aT)^*=aT^*$；
（2）若 $T^{-1}\in\mathcal L(Y,X)$，则 $(T^*)^{-1}$ 存在，并且
$(T^*)^{-1}=(T^{-1})^*$。
\end{problem}

\begin{proof}
对 $y^*\in Y^*$ 和 $x\in X$，
\[
 ((T+S)^*y^*)(x)=y^*((T+S)x)=T^*y^*(x)+S^*y^*(x),
\]
故 $(T+S)^*=T^*+S^*$。同理
$((aT)^*y^*)(x)=a(T^*y^*)(x)$，所以 $(aT)^*=aT^*$（复数域采用线性泛函约定；若内积
约定导致伴随共轭线性，则相应写作 $\bar aT^*$）。

若 $T^{-1}\in\mathcal L(Y,X)$，利用 $(AB)^*=B^*A^*$ 得
\[
 T^*(T^{-1})^*=(T^{-1}T)^*=I_{X^*},\qquad
 (T^{-1})^*T^*=(TT^{-1})^*=I_{Y^*}.
\]
所以 $(T^*)^{-1}=(T^{-1})^*$。
\end{proof}

\subsection*{20.3}
\begin{problem}
设 $L_A:\mathbb R^n\to\mathbb R^m$ 满足 $L_Ax=Ax$，其中
$A\in M_{m,n}(\mathbb R)$、$x\in\mathbb R^n$。求 $(L_A)^*$。
\end{problem}

在欧氏内积下，对 $x\in\mathbb R^n,y\in\mathbb R^m$，
\[
 \langle L_Ax,y\rangle_{\mathbb R^m}=\langle Ax,y\rangle
 =x^{\mathsf T}A^{\mathsf T}y
 =\langle x,A^{\mathsf T}y\rangle_{\mathbb R^n}.
\]
故 $(L_A)^*=L_{A^{\mathsf T}}$。

\subsection*{20.4}
\begin{problem}
设 $X$ 是 Banach 空间，$Y$ 是 $X$ 的闭子空间，$\pi:X\to X/Y$ 为商映射
$\pi(x)=[x]$，$\ell:Y\to X$ 为包含映射 $\ell(y)=y$。证明：

（1）$\pi^*:(X/Y)^*\to Y^\perp$ 是等距同构；

（2）$\ell^*$ 是满射且 $\ker\ell^*=Y^\perp$，从而诱导映射
$\ell^*:X^*/Y^\perp\to Y^*$ 是等距同构。
\end{problem}

\begin{proof}
\textbf{(1)} 对 $\varphi\in(X/Y)^*$，$\pi^*\varphi=\varphi\circ\pi$ 在 $Y$ 上为零，故像落在
$Y^\perp$。又因 $\pi(B_X)$ 在 $B_{X/Y}$ 中稠密，
\[
 \|\pi^*\varphi\|=\sup_{\|x\|\leq1}|\varphi([x])|=\|\varphi\|.
\]
所以 $\pi^*$ 是等距单射。给定 $f\in Y^\perp$，定义
$\widetilde f([x])=f(x)$。若 $[x]=[x']$，则 $x-x'\in Y$，故定义良好；并且
$|f(x)|\leq\|f\|\|[x]\|$（对 $x+y$ 取下确界），所以 $\widetilde f\in(X/Y)^*$ 且
$\pi^*\widetilde f=f$。因此 $\pi^*$ 是到 $Y^\perp$ 的等距同构。

\textbf{(2)} $\ell^*f=f|_Y$。Hahn--Banach 定理说明每个 $g\in Y^*$ 都有保持范数的延拓
$f\in X^*$，故 $\ell^*$ 满射。显然
$\ker\ell^*=\{f:f|_Y=0\}=Y^\perp$。由商空间上的第一同构定理，诱导映射
\[
 X^*/Y^\perp\longrightarrow Y^*,\qquad[f]\longmapsto f|_Y
\]
是线性双射。其等距性来自 Hahn--Banach：一方面
$\|f|_Y\|\leq\inf_{h\in Y^\perp}\|f+h\|$；另一方面 $f|_Y$ 有保持范数的延拓 $F$，且
$F-f\in Y^\perp$，于是反向不等式成立。
\end{proof}

\section{第二十一次作业}

\subsection*{21.1}
\begin{problem}
设 $T:\ell^2\to\ell^2$ 满足
$T((x_i)_{i\geq1})=(i^{-1}x_i)_{i\geq1}$。证明 $\operatorname{Ran}T$ 不是
$\ell^2$ 的闭子空间。
\end{problem}

令 $T(x_i)=(x_i/i)$。取 $y=(1/i)_{i\geq1}$，则 $y\in\ell^2$。若 $y=Tx$，必有
$x_i=iy_i=1$，这不属于 $\ell^2$，所以 $y\notin\operatorname{Ran}T$。另一方面令
$x^{(N)}=(1,\ldots,1,0,\ldots)$（前 $N$ 项为 $1$），则 $x^{(N)}\in\ell^2$，并且
$Tx^{(N)}=(1,1/2,\ldots,1/N,0,\ldots)\to y$ 于 $\ell^2$。所以像不是闭集。

\subsection*{21.2}
\begin{problem}
设 $X,Y$ 是 Banach 空间，$T\in\mathcal L(X,Y)$ 是单射，$N$ 是 $Y$ 的闭子空间，且
$Y=\operatorname{Ran}T\oplus N$。证明 $\operatorname{Ran}T$ 是闭子空间。
\end{problem}

设 $Q:Y\to Y/N$ 是商映射。因 $Y=\operatorname{Ran}T\oplus N$，$QT:X\to Y/N$ 是双射；
它有界，且 $X,Y/N$ 都是 Banach 空间，故由有界逆定理存在 $c>0$ 使
$\|x\|\leq c\|QTx\|\leq c\|Tx\|$。若 $Tx_n\to y$，则
$\|x_n-x_m\|\leq c\|Tx_n-Tx_m\|$，故 $x_n\to x$，进而 $y=Tx$。所以
$\operatorname{Ran}T$ 闭。

\subsection*{21.3}
\begin{problem}
设 $X,Y$ 是 Banach 空间，$S,T\in\mathcal L(X,Y)$。若存在 $c>0$ 使
$\|Sx\|\leq c\|Tx\|$ 对所有 $x\in X$ 成立，且 $\ker T=\ker S$，证明：若
$\operatorname{Ran}S$ 闭，则 $\operatorname{Ran}T$ 闭。
\end{problem}

因 $\operatorname{Ran}S$ 闭，诱导算子
\[
 \widetilde S:X/\ker S\to\operatorname{Ran}S,\qquad[x]\mapsto Sx
\]
是 Banach 空间间的有界双射。因此存在 $C>0$ 使
\[
 \operatorname{dist}(x,\ker S)\leq C\|Sx\|.
\]
由 $\ker S=\ker T$ 及 $\|Sx\|\leq c\|Tx\|$，得
$\operatorname{dist}(x,\ker T)\leq Cc\|Tx\|$。这说明
$\widetilde T:X/\ker T\to\operatorname{Ran}T$ 的逆有界，因而像完备，作为 $Y$ 的线性子空间必闭。

\subsection*{21.4}
\begin{problem}
设 $X,Y$ 是 Hilbert 空间，$S,T\in\mathcal L(X,Y)$。若存在 $c>0$ 使
$\|Sx\|\leq c\|Tx\|$ 对所有 $x\in X$ 成立，且
$T(\ker S)\subset T((\ker S)^\perp)$，证明：若 $\operatorname{Ran}S$ 闭，则
$\operatorname{Ran}T$ 闭。
\end{problem}

设 $M=(\ker S)^\perp$。因 $\operatorname{Ran}S$ 闭，$S|_M:M\to\operatorname{Ran}S$
是有界双射，故存在 $C>0$ 使
\[
 \|z\|\leq C\|Sz\|\leq Cc\|Tz\|\qquad(z\in M). \tag{21.1}
\]
先说明 $\operatorname{Ran}T=T(M)$。给定 $x=u+v$，其中 $u\in\ker S,v\in M$。按假设，
存在 $w\in M$ 使 $Tu=Tw$，所以 $Tx=T(v+w)\in T(M)$；反向包含显然。
若 $Tx_n\to y$，把每个 $Tx_n$ 写成 $Tz_n$（$z_n\in M$）。由 (21.1)，
\[
 \|z_n-z_m\|\leq Cc\|Tz_n-Tz_m\|,
\]
故 $z_n\to z\in M$，于是 $y=Tz$。所以 $\operatorname{Ran}T$ 闭。

\section{第二十二次作业}

\subsection*{22.1}
\begin{problem}
设 $X,Y$ 为 Banach 空间，$K:X\to Y$ 为紧算子。证明：
（1）若 $\operatorname{Im}K$ 是闭集，则 $\dim\operatorname{Im}K<\infty$；
（2）$\operatorname{Im}K\subset Y$ 是可分子空间；
（3）若 $Y$ 可分，则存在 Banach 空间 $X$ 和具有稠密像的紧算子 $K:X\to Y$。
\end{problem}

\begin{proof}
\textbf{(1)} 设 $M=\operatorname{Ran}K$ 闭，则 $M$ 是 Banach 空间。把 $K$ 看作到 $M$ 的满射，
由开映射定理存在 $c>0$ 使 $B_M\subset K(cB_X)$。因 $K$ 紧，$K(cB_X)$ 相对紧，从而
$B_M$ 相对紧。Riesz 引理说明无限维赋范空间的单位球不可能相对紧，故 $\dim M<\infty$。

\textbf{(2)} 对每个 $n$，相对紧集 $K(nB_X)$ 是可分的；取其可数稠密集 $D_n$。因为
\[
 \operatorname{Ran}K=\bigcup_{n=1}^\infty K(nB_X),
\]
$D=\bigcup_nD_n$ 可数且在 $\operatorname{Ran}K$ 中稠密，所以像可分。

\textbf{(3)} 取 $Y$ 中稠密序列 $(y_n)$，不妨令 $y_n\neq0$。定义
\[
 K:\ell^1\to Y,\qquad K(a)=\sum_{n=1}^\infty a_n\frac{y_n}{2^n\|y_n\|}.
\]
该级数绝对收敛，$K$ 有界。令 $K_N$ 为前 $N$ 项截断，则 $K_N$ 有限秩，且
$\|K-K_N\|\leq2^{-N}$，故 $K$ 是紧算子。又
$y_n=2^n\|y_n\|K(e_n)$，所以所有 $y_n$ 都在像中，像因而稠密。
\end{proof}

\subsection*{22.2}
\begin{problem}
设 $X$ 为 Banach 空间，$C\subset X$ 为闭子集。证明以下结论等价：

（1）$C$ 是紧集；

（2）存在序列 $(x_n)\subset C$，使 $\|x_n\|\to0$，且
$C\subset\overline{\operatorname{conv}}\{x_n:n\in\mathbb N\}$。

提示 1：若 $\|x_n\|\to0$，则
\[
 \overline{\operatorname{conv}}\{x_n:n\in\mathbb N\}
 =\left\{\sum_{n=1}^{\infty}\lambda_nx_n:
 \lambda_n\geq0,\ \sum_{n=1}^{\infty}\lambda_n\leq1\right\}.
\]

提示 2：为证明（1）推出（2），选取紧集 $C_k\subset X$ 及有限集 $A_k\subset C_k$，使
$C_1=C$，并且对每个 $k\in\mathbb N$ 有
\[
 2C_k\subset\bigcup_{x\in A_k}B_{4^{-k}}(x),\qquad
 C_{k+1}:=\bigcup_{x\in A_k}\bigl((2C_k\cap B_{4^{-k}}(x))-x\bigr).
\]
证明对每个 $c\in C$，存在 $x_k\in A_k$ 使
$c=\sum_{k=1}^{\infty}2^{-k}x_k$；并注意当 $x\in A_{k+1}$ 时
$\|x\|\leq4^{-k}$。
\end{problem}

\textbf{题面校正。}题面若确实要求 $x_n\in C$，命题对一般紧集是假的，例如
$C=\{x\}$ 且 $x\neq0$。标准结论应为：存在 $x_n\in X$，$\|x_n\|\to0$，且
$C\subset\overline{\operatorname{conv}}\{x_n:n\in\mathbb N\}$。以下证明这一标准版本。

\begin{proof}
先证“存在零列”推出紧性。设 $D=\{x_n:n\geq1\}\cup\{0\}$。因 $x_n\to0$，$D$ 紧。
更直接地，对任意 $\varepsilon>0$ 取 $N$ 使 $\|x_n\|<\varepsilon/2$（$n>N$）。任意
$z\in\overline{\operatorname{conv}}D$ 可写成极限形式
$z\approx\sum_n\lambda_nx_n$，其中 $\lambda_n\geq0,\sum_n\lambda_n\leq1$。尾项范数至多
$\varepsilon/2$；前 $N$ 项所在集合
\[
 \Bigl\{\sum_{n=1}^N\lambda_nx_n:\lambda_n\geq0,\ \sum_{n=1}^N\lambda_n\leq1\Bigr\}
\]
是有限维紧集，故有有限 $\varepsilon/2$-网。于是 $\overline{\operatorname{conv}}D$ 全有界；它闭且
$X$ 完备，故紧。闭子集 $C$ 因而紧。

反之设 $C$ 紧。按提示递归构造紧集 $C_k$ 和有限集 $A_k\subset C_k$：令 $C_1=C$；由
$2C_k$ 紧，取有限 $4^{-k}$-网 $A_k$，并置
\[
 C_{k+1}=\bigcup_{a\in A_k}\bigl((2C_k\cap B_{4^{-k}}(a))-a\bigr).
\]
有限并保持紧性，且 $C_{k+1}\subset B_{4^{-k}}(0)$。给定 $c_1=c\in C_1$，依次选
$a_k\in A_k$ 使 $\|2c_k-a_k\|<4^{-k}$，并令 $c_{k+1}=2c_k-a_k\in C_{k+1}$。递推得
\[
 c=\sum_{k=1}^N2^{-k}a_k+2^{-N}c_{N+1}.
\]
而 $\|c_{N+1}\|\leq4^{-N}$，故余项趋于 $0$，于是
$c=\sum_{k\geq1}2^{-k}a_k$。把各有限集 $A_k$ 的元素依次列成一个序列 $(x_n)$；因
$\sup_{a\in A_k}\|a\|\to0$（$A_k\subset C_k\subset B_{4^{-(k-1)}}$，首项有限不影响），
有 $x_n\to0$。上述展开的系数总和 $\sum2^{-k}=1$，故
$c\in\overline{\operatorname{conv}}\{x_n\}$。任意 $c\in C$ 均如此，结论成立。
\end{proof}

\subsection*{22.3（Pitt 定理）}
\begin{problem}
设 $1\leq p<q<\infty$，按照以下步骤证明所有有界线性算子
$A:\ell^q\to\ell^p$ 都是紧算子：

（1）固定 $A:\ell^q\to\ell^p$、$\|A\|=1$，以及在 $\ell^q$ 中弱收敛于 $0$ 的序列
$(x_n)$，说明只需证明 $\|Ax_n\|_p\to0$；

（2）若 $(y_n)$ 在 $\ell^p$ 中弱收敛于 $0$，证明对每个 $y\in\ell^p$，
\[
 \limsup_{n\to\infty}\|y+y_n\|_p^p
 =\|y\|_p^p+\limsup_{n\to\infty}\|y_n\|_p^p;
\]

（3）固定 $\varepsilon>0$，选取 $x\in\ell^q$ 使
$\|x\|_q=1$ 且 $1-\varepsilon<\|Ax\|_p<1$。证明对任意 $\lambda>0$，
\[
 \bigl(\|Ax\|_p^p+\lambda^p\limsup_n\|Ax_n\|_p^p\bigr)^{1/p}
 \leq
 \bigl(\|x\|_q^q+\lambda^q\limsup_n\|x_n\|_q^q\bigr)^{1/q}.
\]
提示：令 $y_n=\lambda Ax_n$，使用（2）以及
$\|Ax+\lambda Ax_n\|_p\leq\|x+\lambda x_n\|_q$；

（4）证明存在 $C>0$，使对任意 $\lambda,\varepsilon>0$，
\[
 \lambda^p\limsup_n\|Ax_n\|_p^p
 \leq(1+\lambda^qC^q)^{p/q}-(1-\varepsilon)^p.
\]
提示：取 $C\geq\sup_n\|x_n\|_q$；

（5）在（4）中令 $\lambda=C^{-1}\varepsilon^{1/q}$，得到
\[
 \limsup_n\|Ax_n\|_p^p
 \leq C^p\varepsilon^{1-p/q}
 \left(\frac{(1+\varepsilon)^{p/q}-1}{\varepsilon}
 +\frac{1-(1-\varepsilon)^p}{\varepsilon}\right),
\]
再令 $\varepsilon\to0$，证明 $\|Ax_n\|_p\to0$。
\end{problem}

\begin{proof}
设 $A:\ell^q\to\ell^p$ 有界。缩放后可设 $\|A\|=1$。因 $1<q<\infty$ 时 $\ell^q$
自反，证明 $A$ 把弱收敛到 $0$ 的序列送到范数收敛到 $0$ 的序列，即可由序列判别得到紧性。

先证题目中的等式。若 $y_n\rightharpoonup0$ 于 $\ell^p$，则每一坐标 $y_n(j)\to0$。
固定 $N$，前 $N$ 个坐标逐项取极限；尾部用
$|a+b|^p\leq(1+\eta)|a|^p+C_\eta|b|^p$ 控制，再令 $N\to\infty$、$\eta\downarrow0$，得到
\[
 \limsup_n\|y+y_n\|_p^p=\|y\|_p^p+\limsup_n\|y_n\|_p^p. \tag{22.1}
\]

现设 $x_n\rightharpoonup0$ 于 $\ell^q$，并令 $C\geq\sup_n\|x_n\|_q$。给定 $\varepsilon>0$，
由 $\|A\|=1$ 可取 $\|x\|_q=1$ 且 $1-\varepsilon<\|Ax\|_p\leq1$。对任意 $\lambda>0$，
由 (22.1)（分别用于 $p$ 和 $q$）以及 $\|A(x+\lambda x_n)\|_p\leq\|x+\lambda x_n\|_q$，有
\[
 \bigl(\|Ax\|_p^p+\lambda^pL\bigr)^{1/p}
 \leq\bigl(1+\lambda^qM\bigr)^{1/q},
\]
其中 $L=\limsup\|Ax_n\|_p^p$，$M=\limsup\|x_n\|_q^q\leq C^q$。因此
\[
 \lambda^pL\leq(1+\lambda^qC^q)^{p/q}-(1-\varepsilon)^p.
\]
令 $\lambda=C^{-1}\varepsilon^{1/q}$（$C=0$ 时结论显然），得
\[
 L\leq C^p\varepsilon^{1-p/q}
 \left(\frac{(1+\varepsilon)^{p/q}-1}{\varepsilon}
       +\frac{1-(1-\varepsilon)^p}{\varepsilon}\right).
\]
括号在 $\varepsilon\downarrow0$ 时有界，而 $1-p/q>0$，故 $L=0$，即
$\|Ax_n\|_p\to0$。于是 $A$ 紧。
\end{proof}

\section{第二十三次作业}

\subsection*{23.1}
\begin{problem}
设 $X,Y$ 为 Banach 空间，并且存在从 $X$ 到 $Y$ 的 Fredholm 算子。证明：
（1）$X$ 自反当且仅当 $Y$ 自反；
（2）$X$ 可分当且仅当 $Y$ 可分。
\end{problem}

设 $T:X\to Y$ 为 Fredholm 算子。取闭补空间
$X=\ker T\oplus X_0$，以及有限维补空间
$Y=\operatorname{Ran}T\oplus Y_0$。则 $T|_{X_0}:X_0\to\operatorname{Ran}T$ 是有界双射，故
$X_0\cong\operatorname{Ran}T$。

自反性对闭子空间、有限维扩张和有限直和保持。因此
\[
 X\text{ 自反}\iff X_0\text{ 自反}
 \iff\operatorname{Ran}T\text{ 自反}\iff Y\text{ 自反}.
\]
同理，可分性对闭子空间、连续像及有限维扩张保持，所以
$X$ 可分当且仅当 $X_0$ 可分，当且仅当 $\operatorname{Ran}T$ 可分，当且仅当 $Y$ 可分。

\subsection*{23.2}
\begin{problem}
设 $X$ 为实 Banach 空间。证明：任意两个 $X$ 的余维数为 $1$ 的闭子空间同构。

提示：考虑两个子空间 $Y,Z$ 以及 $(Y\cap Z)\times\mathbb R$。
\end{problem}

设 $Y,Z\subset X$ 都是余维 $1$ 的闭子空间。若 $Y=Z$ 无须证明。否则 $Y\not\subset Z$ 且
$Z\not\subset Y$，故 $Y\cap Z$ 在 $Y$、$Z$ 中都余维 $1$。取
$y_0\in Y\setminus Z,z_0\in Z\setminus Y$，则
\[
 Y=(Y\cap Z)\oplus\mathbb Ry_0,\qquad
 Z=(Y\cap Z)\oplus\mathbb Rz_0.
\]
这些有限维补空间上的投影有界（也可用定义 $Y=\ker f,Z=\ker g$ 后显式写投影）。映射
$u+ty_0\mapsto u+tz_0$ 是有界线性双射，其逆同样有界，因此 $Y\cong Z$。

\subsection*{23.3}
\begin{problem}
设 $X$ 为无限维实 Banach 空间。证明以下命题等价：

（1）$X$ 与 $X\times\mathbb R$ 同构；

（2）存在 $X$ 的一个余维数为 $1$ 的子空间与 $X$ 同构；

（3）$X$ 的所有余维数为 $1$ 的闭子空间都与 $X$ 同构；

（4）存在指数为 $1$ 的 Fredholm 算子 $A:X\to X$。
\end{problem}

\begin{proof}
$(1)\Rightarrow(2)$：在 $X\times\mathbb R$ 中，$X\times\{0\}$ 余维 $1$ 且同构于 $X$；经
$X\cong X\times\mathbb R$ 拉回即得。

$(2)\Rightarrow(3)$：由 23.2，任意两个余维 $1$ 闭子空间同构。

$(3)\Rightarrow(1)$：取任意非零 $f\in X^*$ 和 $x_0$ 使 $f(x_0)=1$，则
$X=\ker f\oplus\mathbb Rx_0\cong\ker f\times\mathbb R$。由 (3)，$\ker f\cong X$，所以
$X\cong X\times\mathbb R$。

$(1)\Rightarrow(4)$：取同构 $J:X\to X\times\mathbb R$，再取投影
$P:X\times\mathbb R\to X$。则 $A=PJ$ 满射，且 $\ker A=J^{-1}(\{0\}\times\mathbb R)$
一维，故 $A$ 是指数 $1$ 的 Fredholm 算子。

$(4)\Rightarrow(1)$：设 $k=\dim\ker A$，$c=\operatorname{codim}\operatorname{Ran}A$，则
$k-c=1$。取闭补空间 $X=X_0\oplus\ker A$ 及
$X=\operatorname{Ran}A\oplus\mathbb R^c$。因 $X_0\cong\operatorname{Ran}A$，
\[
 X\cong X_0\oplus\mathbb R^k
   \cong(X_0\oplus\mathbb R^c)\oplus\mathbb R
   \cong X\oplus\mathbb R.
\]
故四个命题等价。
\end{proof}

\subsection*{23.4}
\begin{problem}
记 $I=[0,1]$，固定 $p\geq1$，定义
\[
 W^{1,p}(I)=\left\{f:I\to\mathbb R: f\text{ 绝对连续且 }
 \int_0^1|f'(t)|^pdt<\infty\right\},
\]
并赋予范数
\[
 \|f\|_{W^{1,p}}=\left(\int_0^1(|f(t)|^p+|f'(t)|^p)dt\right)^{1/p}.
\]
证明：

（1）$W^{1,p}(I)$ 在该范数下是 Banach 空间；

（2）从 $W^{1,p}(I)$ 到带上确界范数的 $C(I)$ 的包含映射是有界线性算子；

（3）上述包含映射在 $p>1$ 时是紧算子，而在 $p=1$ 时不是。

提示：当 $p>1$ 时证明单位球等度连续并应用 Arzelà--Ascoli 定理；当 $p=1$ 时考虑
$f_n(t)=t^n$。
\end{problem}

\begin{proof}
\textbf{(1)} 定义 $J:W^{1,p}(I)\to L^p(I)\times L^p(I)$ 为 $Jf=(f,f')$，其范数恰为题设范数。
若 $(f_n,f_n')\to(g,h)$ 于 $L^p\times L^p$，则取子列使 $f_n\to g,f_n'\to h$ 几乎处处；且
\[
 f_n(t)=f_n(0)+\int_0^t f_n'(s)\,ds.
\]
由下面 (2) 的估计，$(f_n(0))$ 为 Cauchy 列，设极限为 $a$。Hölder 不等式给出
$\int_0^t f_n'\to\int_0^t h$，故 $g$ 几乎处处等于绝对连续函数
$a+\int_0^t h$，其导数为 $h$。所以 $J$ 的像闭，$W^{1,p}$ 完备。

\textbf{(2)} 存在 $s\in[0,1]$ 使 $|f(s)|\leq\|f\|_p$；否则积分将大于 $\|f\|_p^p$。于是
\[
 |f(t)|\leq |f(s)|+\int_0^1|f'(u)|\,du
 \leq\|f\|_p+\|f'\|_p\leq2^{1-1/p}\|f\|_{W^{1,p}}.
\]
故嵌入到 $C(I)$ 有界。

\textbf{(3)} 若 $p>1$，令 $p'=p/(p-1)$。单位球中任意 $f$ 满足一致有界性，且
\[
 |f(t)-f(s)|\leq\int_s^t|f'(u)|du
 \leq\|f'\|_p|t-s|^{1/p'}\leq|t-s|^{1/p'}.
\]
故单位球等度连续。Arzelà--Ascoli 定理说明其在 $C(I)$ 中相对紧，嵌入为紧算子。

当 $p=1$，取 $f_n(t)=t^n$。有
$\|f_n\|_{W^{1,1}}=\int_0^1t^n dt+\int_0^1nt^{n-1}dt=1+1/(n+1)$，故序列有界。
若存在一致收敛子列，其极限必连续；但逐点极限在 $[0,1)$ 为 $0$，在 $1$ 为 $1$，不连续，矛盾。
故嵌入不紧。
\end{proof}

\section{第二十四次作业}

\subsection*{24.1}
\begin{problem}
记 $X=(X,\|\cdot\|;\mathbb C)$，并把同一赋范空间视为实赋范空间
$X_{\mathbb R}=(X,\|\cdot\|;\mathbb R)$。若
$\Lambda\in\mathcal L_{\mathbb C}(X;\mathbb C)$，证明：

（1）$\operatorname{Re}\Lambda\in\mathcal L(X_{\mathbb R};\mathbb R)$；

（2）映射
\[
 \Phi:\mathcal L_{\mathbb C}(X;\mathbb C)
 \longrightarrow\mathcal L(X_{\mathbb R};\mathbb R),
 \qquad \Phi(\Lambda)=\operatorname{Re}\Lambda
\]
是等距同构。
\end{problem}

设 $X_{\mathbb R}$ 表示忘掉复数乘法后的实空间。
\begin{proof}
若 $\Lambda\in X^*$ 为复线性，则 $\operatorname{Re}\Lambda$ 显然实线性，且
$|\operatorname{Re}\Lambda(x)|\leq|\Lambda(x)|\leq\|\Lambda\|\|x\|$。
反过来，给定 $u\in(X_{\mathbb R})^*$，定义
\[
 \Lambda_u(x)=u(x)-iu(ix).
\]
直接计算得 $\Lambda_u(ix)=i\Lambda_u(x)$，故它复线性，且 $\operatorname{Re}\Lambda_u=u$；唯一性来自
$\Lambda(x)=\operatorname{Re}\Lambda(x)-i\operatorname{Re}\Lambda(ix)$。最后，对每个 $x$ 取
$\theta$ 使 $e^{i\theta}\Lambda(x)=|\Lambda(x)|$，则
\[
 |\Lambda(x)|=|\operatorname{Re}\Lambda(e^{i\theta}x)|\leq\|\operatorname{Re}\Lambda\|\|x\|.
\]
结合反向显然不等式，$\|\Lambda\|=\|\operatorname{Re}\Lambda\|$。故 $\Phi$ 是等距同构。
\end{proof}

\subsection*{24.2}
\begin{problem}
设 $X=(X,\|\cdot\|;\mathbb R)$，对 $z\in X^c$ 定义
\[
 \|z\|_{X^c}:=\sup_{\theta\in\mathbb R}
 \left(\|\operatorname{Re}(e^{i\theta}z)\|_X^2+
       \|\operatorname{Im}(e^{i\theta}z)\|_X^2\right)^{1/2}.
\]
证明：

（1）$(X^c,\|\cdot\|_{X^c})$ 是复赋范空间；若 $X$ 是 Banach 空间，则 $X^c$ 也是
Banach 空间；

（2）设 $A:X\to Y$ 为实线性算子，并定义
$A^c:X^c\to Y^c$，
\[
 A^c(x+iy)=Ax+iAy\qquad(x,y\in X).
\]
若 $A\in\mathcal L(X,Y)$，证明 $A^c\in\mathcal L_{\mathbb C}(X^c,Y^c)$ 且
$\|A^c\|=\|A\|$。
\end{problem}

写 $z=x+iy$，并令
\[
 \|z\|_{X^c}=\sup_{\theta\in\mathbb R}
 \bigl(\|\operatorname{Re}(e^{i\theta}z)\|_X^2+
       \|\operatorname{Im}(e^{i\theta}z)\|_X^2\bigr)^{1/2}.
\]
每个固定 $\theta$ 对应的表达式是 $X\times X$ 的积范数；取上确界保持三角不等式。
若范数为零，取 $\theta=0$ 得 $x=y=0$。对 $a=re^{i\phi}$，变量替换
$\theta\mapsto\theta+\phi$ 给出 $\|az\|=r\|z\|$，故这是复范数。
并且取 $\theta=0$ 及使用三角不等式可得
\[
 \max(\|x\|,\|y\|)\leq\|x+iy\|_{X^c}
 \leq\sqrt2(\|x\|+\|y\|).
\]
所以它与积空间范数等价；若 $X$ 完备，则 $X^c$ 完备。

若 $A\in\mathcal L(X,Y)$，则 $A^c(x+iy)=Ax+iAy$ 复线性，且对每个 $\theta$，
\[
 \|\operatorname{Re}(e^{i\theta}A^cz)\|^2+
 \|\operatorname{Im}(e^{i\theta}A^cz)\|^2
 \leq\|A\|^2\bigl(\|\operatorname{Re}(e^{i\theta}z)\|^2+
 \|\operatorname{Im}(e^{i\theta}z)\|^2\bigr).
\]
故 $\|A^c\|\leq\|A\|$。在实向量 $z=x$ 上有
$\|A^cx\|=\|Ax\|$、$\|x\|_{X^c}=\|x\|$，取上确界得反向不等式。因此
$\|A^c\|=\|A\|$。

\subsection*{24.3}
\begin{problem}
设 $\Lambda_1,\Lambda_2\in\mathcal L(X;\mathbb R)$，对 $x,y\in X$ 定义
\[
 \Lambda^c(x+iy)=\Lambda_1(x)-\Lambda_2(y)
 +i\bigl(\Lambda_2(x)+\Lambda_1(y)\bigr).
\]
证明：

（1）映射
\[
 \Phi:(\mathcal L(X;\mathbb R))^c\longrightarrow
 \mathcal L_{\mathbb C}(X^c;\mathbb C),
 \qquad \Phi(\Lambda_1+i\Lambda_2)=\Lambda^c
\]
是同构；

（2）若 $X$ 是 Hilbert 空间，则 $\Phi$ 还是等距的。
\end{problem}

给定 $\Lambda_1,\Lambda_2\in X^*$，题设公式直接验证 $\Lambda^c$ 对 $x+iy$ 复线性且有界。
若 $F\in(X^c)^*$，令
$\Lambda_1(x)=\operatorname{Re}F(x)$、$\Lambda_2(x)=\operatorname{Im}F(x)$；由复线性
$F(x+iy)=F(x)+iF(y)$，恰好恢复题设公式。因此 $\Phi$ 为双射；由前述等价范数估计，它和逆都
有界，故为同构。

若 $X$ 为实 Hilbert 空间，复化内积取
\[
 \langle x+iy,u+iv\rangle_c
 =\langle x,u\rangle+\langle y,v\rangle
 +i(\langle y,u\rangle-\langle x,v\rangle).
\]
由 Riesz 表示，$\Lambda_1(x)=\langle x,u\rangle$、
$\Lambda_2(x)=\langle x,v\rangle$，且 $\Lambda^c$ 由向量 $u+iv$ 表示。实对偶复化范数和复
Hilbert 对偶范数都等于 $(\|u\|^2+\|v\|^2)^{1/2}$，故 $\Phi$ 等距。

\section{第二十五次作业}

\subsection*{25.1}
\begin{problem}
设右移算子 $R:\ell^2\to\ell^2$ 为
\[
 R((x_i)_{i\geq1})=(0,x_1,x_2,\ldots).
\]
证明其连续谱和剩余谱分别为
\[
 \sigma_c(R)=\{\lambda\in\mathbb C:|\lambda|=1\},\qquad
 \sigma_r(R)=\{\lambda\in\mathbb C:|\lambda|<1\},
\]
从而点谱 $\sigma_p(R)$ 为空。

提示：连续谱部分可参考讲义中的例子；剩余谱的关键是证明当 $|\lambda|<1$ 时
\[
 \operatorname{Ran}(\lambda I-R)
 =\{(\bar\lambda^i)_{i\geq0}\}^{\perp}.
\]
\end{problem}

令 $R(x_1,x_2,\ldots)=(0,x_1,x_2,\ldots)$。
\begin{proof}
若 $|\lambda|>1=\|R\|$，则
$\lambda I-R=\lambda(I-R/\lambda)$ 可由 Neumann 级数求逆，故在预解集。

若 $|\lambda|<1$，方程 $(\lambda I-R)x=0$ 逐坐标给出 $\lambda x_1=0$、
$\lambda x_{n+1}-x_n=0$；无论 $\lambda=0$ 与否都推出 $x=0$，故算子单射。其伴随为左移
$R^*(y_1,y_2,\ldots)=(y_2,y_3,\ldots)$。向量
$v_\lambda=(1,\bar\lambda,\bar\lambda^2,\ldots)\in\ell^2$ 满足
$(\bar\lambda I-R^*)v_\lambda=0$。因此
\[
 \overline{\operatorname{Ran}(\lambda I-R)}
 =\ker(\bar\lambda I-R^*)^\perp=v_\lambda^\perp\neq\ell^2.
\]
所以 $|\lambda|<1$ 属于剩余谱。

若 $|\lambda|=1$，同样的递推说明 $\lambda I-R$ 单射；此时上述形式的特征向量不在 $\ell^2$，
且直接递推可知 $\ker(\bar\lambda I-R^*)=\{0\}$，所以像稠密。另一方面取
\[
 x^{(N)}=N^{-1/2}(1,\lambda^{-1},\ldots,\lambda^{-(N-1)},0,\ldots),
\]
则 $\|x^{(N)}\|=1$ 且 $\|(\lambda I-R)x^{(N)}\|=\sqrt{2/N}\to0$。若算子有有界逆便不可能
发生此事，故它不满射。于是 $|\lambda|=1$ 属于连续谱，点谱为空。
\end{proof}

\subsection*{25.2}
\begin{problem}
设 $X$ 为复 Banach 空间，$A\in\mathcal L_{\mathbb C}(X)$ 是双射。设实数
$\varepsilon,r$ 满足
\[
 0<\varepsilon<\|A^{-1}\|^{-1}\leq\|A\|<r.
\]
证明
$\sigma(A)\subset\{\lambda\in\mathbb C:\varepsilon<|\lambda|<r\}$。
\end{problem}

若 $|\lambda|>\|A\|$，则
$\lambda I-A=\lambda(I-A/\lambda)$ 可逆。若
$|\lambda|<\|A^{-1}\|^{-1}$，则
\[
 \lambda I-A=-A(I-\lambda A^{-1})
\]
也由 Neumann 级数可逆。因此谱只可能满足
$\|A^{-1}\|^{-1}\leq|\lambda|\leq\|A\|$。结合题设严格不等式即
$\sigma(A)\subset\{\varepsilon<|\lambda|<r\}$。

\subsection*{25.3}
\begin{problem}
设 $X$ 为复 Banach 空间，$\Omega\subset\mathbb C$ 为开集，
$f:\Omega\to X$ 为全纯函数。证明：

（1）导数 $f':\Omega\to X$ 也是全纯函数；

（2）$f$ 是光滑函数；

（3）取 $z_0\in\Omega$、$r>0$ 使 $\overline{B_r(z_0)}\subset\Omega$，定义
$\gamma(t)=z_0+re^{2\pi it}$。证明对 $w\in B_r(z_0)$，
\[
 f^{(n)}(w)=\frac{n!}{2\pi i}
 \int_\gamma\frac{f(z)}{(z-w)^{n+1}}\,dz.
\]
\end{problem}

取闭圆盘 $\overline{B_r(z_0)}\subset\Omega$。Banach 值 Cauchy 公式为
\[
 f(w)=\frac1{2\pi i}\int_\gamma\frac{f(z)}{z-w}\,dz.
\]
它可通过对两边施加任意 $x^*\in X^*$、使用标量 Cauchy 公式，再由 Hahn--Banach 分离点得到。
对 $w$ 在更小的闭圆盘内，核函数关于 $w$ 的各阶导数一致有界，所以可在 Bochner 积分号下
逐次求导，得到
\[
 f^{(n)}(w)=\frac{n!}{2\pi i}\int_\gamma\frac{f(z)}{(z-w)^{n+1}}\,dz. \tag{25.1}
\]
右端作为 $w$ 的函数连续，且再次求导仍合法。因此 $f'$ 全纯，归纳得所有阶导数存在且连续，
$f$ 是光滑函数；(25.1) 也正是第 (3) 问公式。

\subsection*{25.4}
\begin{problem}
设 $X$ 为复 Banach 空间，$(a_n)_{n\in\mathbb N}$ 为 $X$ 中的序列，且
\[
 \rho:=\frac{1}{\limsup_{n\to\infty}\|a_n\|^{1/n}}>0.
\]
证明幂级数
$f(z)=\sum_{n=0}^{\infty}a_nz^n$ 对每个 $|z|<\rho$ 收敛，并且
$f:B_\rho(0)\to X$ 是全纯函数。取 $0<r<\rho$，定义
$\gamma(t)=re^{2\pi it}$，证明
\[
 a_n=\frac{f^{(n)}(0)}{n!}
 =\frac1{2\pi i}\int_\gamma\frac{f(z)}{z^{n+1}}\,dz.
\]
\end{problem}

固定 $|z|<r<\rho$。由根值判别，存在 $N$ 与 $c<1/r$，使 $n\geq N$ 时
$\|a_n\|^{1/n}\leq c$，故
$\|a_nz^n\|\leq(c|z|)^n$，级数绝对收敛。并且在每个闭圆盘 $|z|\leq r<\rho$ 上一致收敛。
逐项求导级数也由同样的根值判别一致收敛，所以
\[
 f'(z)=\sum_{n\geq1}na_nz^{n-1},
\]
$f$ 全纯。对圆周 $\gamma(t)=re^{2\pi it}$ 逐项积分（由一致收敛），
\[
 \frac1{2\pi i}\int_\gamma\frac{f(z)}{z^{n+1}}dz
 =\sum_{k\geq0}a_k\frac1{2\pi i}\int_\gamma z^{k-n-1}dz=a_n.
\]
又逐项求导在 $0$ 给出 $f^{(n)}(0)=n!a_n$，故题设两式成立。

\subsection*{25.5}
\begin{problem}
设 $X$ 为复 Banach 空间，$A:X\to X$ 为有界复线性算子，
$p(z)=\sum_{k=0}^na_kz^k$ 为复系数多项式。直接证明
\[
 \sigma(p(A))=p(\sigma(A)),
 \qquad p(A)=\sum_{k=0}^na_kA^k.
\]

提示：为证 $p(\sigma(A))\subset\sigma(p(A))$，固定 $\lambda\in\sigma(A)$，写
$p(z)-p(\lambda)=(z-\lambda)q(z)$。为证反向包含，不妨设 $a_n\neq0$；固定
$\mu\in\sigma(p(A))$，并把 $p(z)-\mu$ 分解为
$a_n\prod_{i=1}^n(z-\lambda_i)$，证明至少有一个 $A-\lambda_iI$ 不是双射。
\end{problem}

若 $\lambda\in\sigma(A)$，则
$p(z)-p(\lambda)=(z-\lambda)q(z)$，故
\[
 p(A)-p(\lambda)I=(A-\lambda I)q(A).
\]
若左端可逆，则其两个可交换因子都可逆（例如
$(A-\lambda I)^{-1}=q(A)[p(A)-p(\lambda)I]^{-1}$），与
$\lambda\in\sigma(A)$ 矛盾。因此 $p(\lambda)\in\sigma(p(A))$。

反之若 $\mu\in\sigma(p(A))$，在复数域分解
$p(z)-\mu=a\prod_{j=1}^n(z-\lambda_j)$。若每个
$A-\lambda_jI$ 都可逆，则它们彼此交换，乘积
$p(A)-\mu I$ 也可逆，矛盾。因此某个 $\lambda_j\in\sigma(A)$，并且
$\mu=p(\lambda_j)$。故 $\sigma(p(A))=p(\sigma(A))$。

\section{第二十六次作业}

\subsection*{26.1}
\begin{problem}
在 25.2 的条件下，证明
\[
 A^{-1}=
 \frac1{2\pi i}\int_{|z|=r}\frac{R(z,A)}{z}\,dz
 -\frac1{2\pi i}\int_{|z|=\varepsilon}\frac{R(z,A)}{z}\,dz.
\]
\end{problem}

取正向边界 $\Gamma=\{|z|=r\}-\{|z|=\varepsilon\}$。由 25.2，谱位于环域
$\varepsilon<|z|<r$。对在谱邻域全纯的函数 $g(z)=1/z$，全纯函数演算给出
\[
 g(A)=\frac1{2\pi i}\int_\Gamma g(z)R(z,A)\,dz.
\]
又 $g(A)=A^{-1}$，而内边界在 $\Gamma$ 中取负向，故
\[
 A^{-1}=\frac1{2\pi i}\int_{|z|=r}\frac{R(z,A)}z,dz
 -\frac1{2\pi i}\int_{|z|=\varepsilon}\frac{R(z,A)}z,dz.
\]

\subsection*{26.2}
\begin{problem}
设 $H$ 是复 Hilbert 空间，$A\in\mathcal L_{\mathbb C}(H)$，$E\subset H$ 为闭复线性子空间。
称 $E$ 在 $A$ 下不变，是指 $x\in E$ 蕴含 $Ax\in E$。证明：
$E$ 在 $A$ 下不变，当且仅当 $E^\perp$ 在 $A^*$ 下不变。
\end{problem}

若 $E$ 在 $A$ 下不变，取 $y\in E^\perp,x\in E$，则
\[
 \langle x,A^*y\rangle=\langle Ax,y\rangle=0,
\]
故 $A^*y\in E^\perp$。反之若 $E^\perp$ 在 $A^*$ 下不变，则对
$x\in E,y\in E^\perp$，仍有
$\langle Ax,y\rangle=\langle x,A^*y\rangle=0$。因此
$Ax\in(E^\perp)^\perp=E$，其中最后一步使用了 $E$ 闭。

\subsection*{26.3}
\begin{problem}
若 $A\in\mathcal L_{\mathbb C}(X)$ 满足
\[
 \ker((\lambda I-A)^m)=\ker((\lambda I-A)^{m+1}),
\]
证明对每个 $k\geq1$，
\[
 \ker((\lambda I-A)^m)=\ker((\lambda I-A)^{m+k}).
\]
\end{problem}

令 $B=\lambda I-A$。已知 $\ker B^m=\ker B^{m+1}$。先证相邻项以后保持相等：若
$x\in\ker B^{m+2}$，则 $Bx\in\ker B^{m+1}=\ker B^m$，所以
$B^{m+1}x=0$，即 $x\in\ker B^{m+1}$。反向包含显然，故
$\ker B^{m+2}=\ker B^{m+1}$。归纳即得
$\ker B^{m+k}=\ker B^m$ 对所有 $k\geq1$ 成立。

\subsection*{26.4}
\begin{problem}
设 $A\in\mathcal L_{\mathbb C}(X)$ 是紧算子，
$\lambda\in\sigma(A)\setminus\{0\}$。取 $\varepsilon>0$ 使
$\{z:|z-\lambda|\leq\varepsilon\}\setminus\{\lambda\}\subset\rho(A)$，定义
\[
 P_\lambda=\frac1{2\pi i}\int_{|z-\lambda|=\varepsilon}R(z,A)\,dz,
 \qquad Y=\operatorname{Ran}P_\lambda.
\]
证明对每个 $m\geq1$，
\[
 \ker((\lambda I-A)^m)=\ker((\lambda I-A|_Y)^m).
\]
\end{problem}

Riesz 投影 $P_\lambda$ 与 $A$ 交换，故 $Y=\operatorname{Ran}P_\lambda$ 在 $A$ 下不变。
显然
\[
 \ker((\lambda I-A|_Y)^m)\subset\ker((\lambda I-A)^m).
\]
反过来，若 $(\lambda I-A)^mx=0$，在 $X=Y\oplus\ker P_\lambda$ 的谱分解中写
$x=y+z$。限制算子 $A|_{\ker P_\lambda}$ 的谱不含 $\lambda$，故
$\lambda I-A$ 在 $\ker P_\lambda$ 上可逆。由
$(\lambda I-A)^my+(\lambda I-A)^mz=0$，两项分别属于互补的不变子空间，因此
$(\lambda I-A)^mz=0$；可逆性推出 $z=0$。所以 $x=y\in Y$，得到反向包含。

\subsection*{26.5}
\begin{problem}
补充证明讲义第 187 页例 3 的证明过程。
\end{problem}

\textbf{依赖未提供讲义，按用户要求跳过。}题面只写“补充证明讲义第 187 页例 3”，未给出例题
陈述，无法唯一还原待证结论。

\subsection*{26.6}
\begin{problem}
设 $A\in\mathcal L_{\mathbb C}(X)$，$r>\|A\|$，并定义
$e^A=\sum_{k=0}^{\infty}A^k/k!$。证明：

（1）
\[
 e^A=\frac1{2\pi i}\int_{|z|=r}e^zR(z,A)\,dz;
\]

（2）$\sigma(e^A)=\{e^\lambda:\lambda\in\sigma(A)\}$；

（3）对任意 $s,t\in\mathbb R$，$e^{(s+t)A}=e^{sA}e^{tA}$。
\end{problem}

\begin{proof}
\textbf{(1)} 在 $|z|=r>\|A\|$ 上有一致收敛的 Neumann 展开
$R(z,A)=\sum_{n\geq0}A^nz^{-n-1}$，而
$e^z=\sum_{k\geq0}z^k/k!$。逐项积分并用
$\frac1{2\pi i}\int_{|z|=r}z^{k-n-1}dz=\delta_{kn}$，得
\[
 \frac1{2\pi i}\int_{|z|=r}e^zR(z,A)dz
 =\sum_{n=0}^\infty\frac{A^n}{n!}=e^A.
\]

\textbf{(2)} 全纯谱映射定理给出
$\sigma(e^A)=e^{\sigma(A)}$。为说明其适用性：若
$\mu\notin e^{\sigma(A)}$，则 $e^z-\mu$ 在谱的邻域无零点，故
$1/(e^z-\mu)$ 在该邻域全纯；函数演算给出它在 $A$ 处的值，并成为
$e^A-\mu I$ 的逆。反之若 $\mu=e^\lambda$、$\lambda\in\sigma(A)$ 而
$e^A-\mu I$ 可逆，把 $e^z-e^\lambda=(z-\lambda)h(z)$ 代入函数演算，将导致
$A-\lambda I$ 可逆，矛盾。

\textbf{(3)} 两个幂级数绝对收敛，可按总次数重排：
\begin{align*}
 e^{sA}e^{tA}
 &=\sum_{m,n\geq0}\frac{s^mt^nA^{m+n}}{m!n!}\\
 &=\sum_{k\geq0}\frac{A^k}{k!}\sum_{m=0}^k\binom{k}{m}s^mt^{k-m}\\
 &=\sum_{k\geq0}\frac{(s+t)^kA^k}{k!}=e^{(s+t)A}.
\end{align*}
\end{proof}

\section{第二十七次作业}

\subsection*{27.1}
\begin{problem}
证明讲义第 188 页所列的 Hilbert 空间上伴随算子的性质 1--9。
\end{problem}

\textbf{依赖未提供讲义，按用户要求跳过。}题面只引用“讲义第 188 页性质 1--9”，没有列出九条
性质，无法确定需要证明的具体命题及所采用的内积约定。

\subsection*{27.2}
\begin{problem}
设 $H$ 是非零复 Hilbert 空间，$A:H\to H$ 为正规算子。

（1）证明
\[
 \operatorname{Re}\lambda\geq0\quad(\lambda\in\sigma(A))
 \quad\Longleftrightarrow\quad
 \operatorname{Re}\langle x,Ax\rangle\geq0\quad(x\in H).
\]
提示：从右向左时，若 $\operatorname{Re}\lambda<0$，对
$\operatorname{Re}\langle x,Ax-\lambda x\rangle$ 使用 Cauchy--Schwarz 不等式；从左向右时，
证明 $\|e^{-tA}\|\leq1$（$t\geq0$），并对 $t\mapsto\|e^{-tA}x\|^2$ 求导。

（2）证明
\[
 \sup_{\|x\|=1}\operatorname{Re}\langle x,Ax\rangle
 =\sup_{\lambda\in\sigma(A)}\operatorname{Re}\lambda,
 \qquad
 \inf_{\|x\|=1}\operatorname{Re}\langle x,Ax\rangle
 =\inf_{\lambda\in\sigma(A)}\operatorname{Re}\lambda.
\]

（3）证明
\[
 \sigma(A)\cap i\mathbb R=\varnothing
 \quad\Longleftrightarrow\quad A+A^*\text{ 是双射}.
\]

（4）证明
\[
 \sigma(A+A^*)=\{\lambda+\bar\lambda:\lambda\in\sigma(A)\}.
\]
提示：对 $A-\mu I$（$\mu\in\mathbb R$）使用（3）的结论。

（5）证明（1）--（4）中“$A$ 是正规算子”的假设不能去掉。
提示：找出 $A\in\mathbb R^{2\times2}$ 以及 $x\in\mathbb R^2$，使
$\sigma(A)=\{0\}$ 且 $\langle x,Ax\rangle>0$。
\end{problem}

以下记
\[
 B=\operatorname{Re}A:=\frac{A+A^*}{2}.
\]
若 $A$ 正规，则连续函数演算说明 $B$ 自伴，并且
\[
 \sigma(B)=\{\operatorname{Re}\lambda:\lambda\in\sigma(A)\}. \tag{27.1}
\]

\begin{proof}
\textbf{(1)} 因 $\operatorname{Re}\langle x,Ax\rangle=\langle x,Bx\rangle$，由自伴算子的谱定理，
$\langle x,Bx\rangle\geq0$ 对所有 $x$ 成立，当且仅当
$\sigma(B)\subset[0,\infty)$。结合 (27.1)，这又当且仅当
$\operatorname{Re}\lambda\geq0$ 对所有 $\lambda\in\sigma(A)$ 成立。

也可核对题目提示的一个方向：若数值实部非负而 $\operatorname{Re}\lambda<0$，则
\[
 \|(A-\lambda I)x\|\|x\|
 \geq\operatorname{Re}\langle x,(A-\lambda I)x\rangle
 \geq-\operatorname{Re}\lambda\,\|x\|^2,
\]
故 $A-\lambda I$ 有正下界；对正规算子，$A-\lambda I$ 亦正规，正下界同时作用于其伴随，
从而像稠密且闭，即满射，与 $\lambda\in\sigma(A)$ 矛盾。

\textbf{(2)} 自伴算子 $B$ 的 Rayleigh 商满足
\[
 \sup_{\|x\|=1}\langle x,Bx\rangle=\max\sigma(B),\qquad
 \inf_{\|x\|=1}\langle x,Bx\rangle=\min\sigma(B).
\]
这是谱定理的直接结果：上界由谱测度积分得到；取谱靠近最大（最小）值的非零谱投影中的单位
向量可任意逼近该端点。再用 (27.1) 和
$\langle x,Bx\rangle=\operatorname{Re}\langle x,Ax\rangle$，即得题设两个等式。

\textbf{(3)} $A+A^*=2B$。由 (27.1)，
\begin{align*}
 0\notin\sigma(A+A^*)
 &\iff0\notin\sigma(B)\\
 &\iff0\notin\{\operatorname{Re}\lambda:\lambda\in\sigma(A)\}\\
 &\iff\sigma(A)\cap i\mathbb R=\varnothing.
\end{align*}
有界算子双射等价于 $0$ 不在其谱中，故结论成立。

\textbf{(4)} 由 $A+A^*=2B$ 和 (27.1)，
\[
 \sigma(A+A^*)=2\sigma(B)
 =\{2\operatorname{Re}\lambda:\lambda\in\sigma(A)\}
 =\{\lambda+\bar\lambda:\lambda\in\sigma(A)\}.
\]

\textbf{(5)} 正规性不能删去。取
\[
 A=\begin{pmatrix}0&1\\0&0\end{pmatrix}.
\]
则 $A^2=0$，所以 $\sigma(A)=\{0\}$，但 $A$ 不是正规算子。取
$x=(1,-1)^{\mathsf T}/\sqrt2$，有
$\operatorname{Re}\langle x,Ax\rangle=-1/2<0$，故 (1) 失败；同时数值实部的上确界为
$1/2$、下确界为 $-1/2$，而谱实部上下确界均为 $0$，故 (2) 失败。
又 $\sigma(A)\cap i\mathbb R=\{0\}\neq\varnothing$，但
\[
 A+A^*=\begin{pmatrix}0&1\\1&0\end{pmatrix}
\]
可逆，故 (3) 失败；其谱为 $\{-1,1\}$，而
$\{\lambda+\bar\lambda:\lambda\in\sigma(A)\}=\{0\}$，故 (4) 失败。
\end{proof}

\end{document}

